% Generated mechanically from frozen input inputs/p11/ja/defos.tex.
% Do not edit this generated file; edit the bound locale source instead.

% OK, start here.
%

\setcounter{chapter}{90}
\chapter{変形理論}



\phantomsection
\label{defos-section-phantom}


\section{序論}
\label{defos-section-introduction}

\noindent
本章の目的は、加群や射などの変形理論への（比較的）穏やかな入門を
与えることである。本章では、素朴な余接複体を用いて証明できる結果を
扱う。余接複体に関する章では、これらの結果をさらに少し拡張する。
この主題について、より進んだ読者は Illusie の論考を参照されたい。
\cite{cotangent} を見よ。





\section{環の変形と素朴な余接複体}
\label{defos-section-deformations}

\noindent
本節では、素朴な余接複体を用いて変形理論を少し展開する。まず、核が
二乗零イデアル $I$ である全射環準同型 $A' \to A$ から始める。
さらに、環準同型 $A \to B$、$B$-加群 $N$、および $A$-加群準同型
$c : I \to N$ が与えられていると仮定する。本節で問うのは、次の図式の
疑問符に入るものを見いだせるかどうかである。
\begin{equation}
\label{defos-equation-to-solve}
\vcenter{
\xymatrix{
0 \ar[r] & N \ar[r] & {?} \ar[r] & B \ar[r] & 0 \\
0 \ar[r] & I \ar[u]^c \ar[r] & A' \ar[u] \ar[r] & A \ar[u] \ar[r] & 0
}
}
\end{equation}
また、解が存在するとき、それがどの程度一意であるかも問う。より正確には、
核が二乗零イデアルで $N$ と同一視され、かつ $A' \to B'$ が与えられた
写像 $c$ を誘導するような、$A'$-代数の全射 $B' \to B$ を求める。
このとき $B'$ を (\ref{defos-equation-to-solve}) の {\it 解} と呼ぶ。

\begin{lemma}
\label{defos-lemma-huge-diagram}
次の可換図式が与えられ、
$$
\xymatrix{
& 0 \ar[r] & N_2 \ar[r] & B'_2 \ar[r] & B_2 \ar[r] & 0 \\
& 0 \ar[r]|\hole & I_2 \ar[u]_{c_2} \ar[r] &
A'_2 \ar[u] \ar[r]|\hole & A_2 \ar[u] \ar[r] & 0 \\
0 \ar[r] & N_1 \ar[ruu] \ar[r] & B'_1 \ar[r] & B_1 \ar[ruu] \ar[r] & 0 \\
0 \ar[r] & I_1 \ar[ruu]|\hole \ar[u]^{c_1} \ar[r] &
A'_1 \ar[ruu]|\hole \ar[u] \ar[r] & A_1 \ar[ruu]|\hole \ar[u] \ar[r] & 0
}
$$
その手前と奥が (\ref{defos-equation-to-solve}) の解であるとする。このとき
\begin{enumerate}
\item
$\Ext^1_{B_1}(\NL_{B_1/A_1}, N_2)$
には正準元が存在し、その消滅は、図式に収まる環準同型
$B'_1 \to B'_2$ が存在するための必要十分条件である。
\item 図式に収まる写像 $B'_1 \to B'_2$ が存在するならば、そのような
写像全体の集合は
$\Hom_{B_1}(\Omega_{B_1/A_1}, N_2)$
の下で主等質空間をなす。
\end{enumerate}
\end{lemma}

\begin{proof}
$E=B_1$ を集合とみなす。核を $J$ とする全射
$A_1[E] \to B_1$ を考える。これは
$$
\NL_{B_1/A_1} = (J/J^2 \to \Omega_{A_1[E]/A_1} \otimes_{A_1[E]} B_1)
$$
という式によって素朴な余接複体を定義するために用いられるもの
である（『代数』第 \ref{algebra-section-netherlander} 節）。
$\Omega_{A_1[E]/A_1} \otimes B_1$ は自由 $B_1$-加群なので、
$$
\Ext^1_{B_1}(\NL_{B_1/A_1}, N_2) =
\frac{\Hom_{B_1}(J/J^2, N_2)}
{\Hom_{B_1}(\Omega_{A_1[E]/A_1} \otimes B_1, N_2)}
$$
右辺の加群の中に障害を構成する。
$J'=\Ker(A'_1[E]\to B_1)$ とおく。核が $I_1A'_1[E]$ である全射
$J'\to J$ が存在することに注意する。各 $e\in E$ に対し、対応する
変数を $x_e\in A_1[E]$ と書く。$B_1$ における $x_e$ の像の持ち上げ
$y_e\in B'_1$ と、$B_2$ における $x_e$ の像の持ち上げ
$z_e\in B'_2$ を選ぶ。これらの選択は $A'_1$-代数準同型
$$
A'_1[E] \to B'_1 \quad\text{and}\quad A'_1[E] \to B'_2
$$
を定める。前者は写像 $J'\to N_1$、$f'\mapsto f'(y_e)$ を与え、
後者は写像 $J'\to N_2$、$f'\mapsto f'(z_e)$ を与える。計算により、
これらの写像は $(J')^2$ を零に送る。図式の左側の正方形
（$c_1$ と $c_2$ を含むもの）が可換なので、$N_2$ への写像として、
これらは $I_1A'_1[E]$ 上で一致する。$B'_1$ は
$J'\to A'_1[E]$ と $J'\to N_1$ の押し出しであることに注意する。
したがって、写像 $J'\to N_1\to N_2$ と $J'\to N_2$ が一致すれば、
図式に収まる写像 $B'_1\to B'_2$ が得られる。そこで、写像
$$
J/J^2 \to N_2,\quad f \mapsto f'(z_e) - \nu(f'(y_e))
$$
の類を障害とする。ここで $\nu:N_1\to N_2$ は与えられた写像であり、
$f'\in J'$ は $f$ の持ち上げである。以上の考察により、これは
良定義である。なお、ある $\delta_{i,e}\in N_i$ によって、
$z_e$ の選択を $z_e+\delta_{2,e}$ に、$y_e$ の選択を
$y_e+\delta_{1,e}$ に変更する自由がある。この変更により上の写像は
$$
f \mapsto f'(z_e + \delta_{2, e}) - \nu(f'(y_e + \delta_{1, e})) =
f'(z_e) - \nu(f'(z_e)) +
\sum (\delta_{2, e} - \nu(\delta_{1, e}))\frac{\partial f}{\partial x_e}
$$
となる。これはまさに、写像 $J/J^2\to N_2$ を合成
$J/J^2\to\Omega_{A_1[E]/A_1}\otimes B_1\to N_2$ だけ変更することを
意味する。ただし後者の写像は $\text{d}x_e$ を
$\delta_{2,e}-\nu(\delta_{1,e})$ に送る。したがって障害は
良定義であり、持ち上げが存在することと障害が零であることは
同値である。

\medskip\noindent
(2) は次の観察から従う。図式に収まる二つの写像
$\varphi,\psi:B'_1\to B'_2$ が与えられると、
$\varphi-\psi$ は $A_1$-導分である写像 $D:B_1\to N_2$ を経由する。
\begin{align*}
D(fg) & = \varphi(f'g') - \psi(f'g') \\
& =
\varphi(f')\varphi(g') - \psi(f')\psi(g') \\
& =
(\varphi(f') - \psi(f'))\varphi(g') + \psi(f')(\varphi(g') - \psi(g')) \\
& =
gD(f) + fD(g)
\end{align*}
したがって $D$ は一意な $B_1$-線形写像
$\Omega_{B_1/A_1}\to N_2$ に対応する。逆に、そのような線形写像から
導分 $D$ が得られ、さらに図式に収まる環準同型
$\psi:B'_1\to B'_2$ が与えられれば、$\psi+D$ も図式に収まる
別の環準同型となる。
\end{proof}

\begin{lemma}
\label{defos-lemma-choices}
(\ref{defos-equation-to-solve}) の解が存在するならば、解の同型類全体の集合は
$\Ext^1_B(\NL_{B/A}, N)$
の下で主等質空間をなす。
\end{lemma}

\begin{proof}
まず、(\ref{defos-equation-to-solve}) の二つの解 $B'_1$ と $B'_2$ が与えられると、
補題 \ref{defos-lemma-huge-diagram} により、写像 $B'_1\to B'_2$ の存在に対する
障害元
$o(B'_1,B'_2)\in\Ext^1_B(\NL_{B/A},N)$
が得られる。この元は明らかに同型の存在に対する障害であり、したがって
同型類を区別する。ゆえに証明を完了するには、解 $B'$ と元
$\xi \in \Ext^1_B(\NL_{B/A}, N)$
が与えられたとき、
$o(B', B'_\xi) = \xi$
を満たす別の解 $B'_\xi$ を見いだせることを示せば十分である。

\medskip\noindent
$E=B$ を集合とみなす。核を $J$ とする全射 $A[E]\to B$ を考える。
これは
$$
\NL_{B/A} = (J/J^2 \to \Omega_{A[E]/A} \otimes_{A[E]} B)
$$
という式によって素朴な余接複体を定義するために用いられるもの
である（『代数』第 \ref{algebra-section-netherlander} 節）。
$\Omega_{A[E]/A}\otimes B$ は自由 $B$-加群なので、
$$
\Ext^1_B(\NL_{B/A}, N) =
\frac{\Hom_B(J/J^2, N)}
{\Hom_B(\Omega_{A[E]/A} \otimes B, N)}
$$
したがって $\xi$ は準同型 $\delta:J/J^2\to N$ の類で表せる。

\medskip\noindent
各 $e\in E$ に対し、対応する変数を $x_e\in A[E]$ と書く。$B$ における
$x_e$ の像の持ち上げ $y_e\in B'$ を選ぶ。これらの選択は
$A'$-代数準同型 $\varphi:A'[E]\to B'$ を定める。
$J'=\Ker(A'[E]\to B)$ とおく。$\varphi$ は写像
$\varphi|_{J'}:J'\to N$ を誘導し、次の図式に示すように $B'$ は
押し出しである。
$$
\xymatrix{
0 \ar[r] & N \ar[r] & B' \ar[r] & B \ar[r] & 0 \\
0 \ar[r] & J' \ar[u]^{\varphi|_{J'}} \ar[r] & A'[E] \ar[u] \ar[r] &
B \ar[u]_{=} \ar[r] & 0
}
$$
$\psi:J'\to N$ を、写像 $\varphi|_{J'}$ と次の合成との和とする。
$$
J' \to J'/(J')^2 \to J/J^2 \xrightarrow{\delta} N.
$$
このとき $\psi$ に沿う押し出しは、上と同様の図式に収まる別の環拡大
$B'_\xi$ である。計算により、望みどおり
$o(B',B'_\xi)=\xi$ が分かる。
\end{proof}

\begin{lemma}
\label{defos-lemma-extensions-of-algebras}
$A$ を環、$B$ を $A$-代数、$N$ を $B$-加群とする。$A$-代数の拡大
$$
0 \to N \to B' \to B \to 0
$$
で $N$ が二乗零イデアルであるものの同型類全体の集合は、
$\Ext^1_B(\NL_{B/A}, N)$
と正準的に全単射で対応する。
\end{lemma}

\begin{proof}
これを示すため、(\ref{defos-equation-to-solve}) が次の図式で与えられる場合に
上の結果を適用する。
$$
\xymatrix{
0 \ar[r] & N \ar[r] & {?} \ar[r] & B \ar[r] & 0 \\
0 \ar[r] & 0 \ar[u] \ar[r] & A \ar[u] \ar[r]^{\text{id}} & A \ar[u] \ar[r] & 0
}
$$
したがって本補題は、補題 \ref{defos-lemma-choices} と、解 $N\oplus B$ が
存在することから従う（全単射の直接的な構成については次の注意を見よ）。
\end{proof}

\begin{remark}
\label{defos-remark-extensions-of-algebras}
$A\to B$ と $N$ を補題 \ref{defos-lemma-extensions-of-algebras} のとおりとする。
$\alpha:P\to B$ を $A$ 上の $B$ の表示とする（『代数』第
\ref{algebra-section-netherlander} 節）。$J=\Ker(\alpha)$ とおくと、
$\alpha$ に付随する素朴な余接複体 $\NL(\alpha)$ は複体
$J/J^2\to\Omega_{P/A}\otimes_PB$ である。また、
$$
\Ext^1_B(\NL(\alpha), N) =
\Coker\left(\Hom_B(\Omega_{P/A} \otimes_P B, N) \to \Hom_B(J/J^2, N)\right)
$$
である。これは $\Omega_{P/A}$ が自由加群だからである。補題における
拡大 $0\to N\to B'\to B\to0$ を考える。$P$ は $A$ 上の多項式代数
なので、$\alpha$ を $A$-代数準同型 $\alpha':P'\to B'$ に持ち上げられる。
$N$ は $B'$ の中で二乗零だから、$\alpha'|_J:J\to N$ は
$J\to J/J^2\to N$ と分解する。本補題の対応は、この拡大を、表示された
余核における写像 $J/J^2\to N$ の類へ送る。
\end{remark}

\begin{lemma}
\label{defos-lemma-extensions-of-algebras-functorial}
環準同型 $A\to B\to C$、$B$-加群 $M$、$C$-加群 $N$、
$B$-線形写像 $c:M\to N$、および二乗零な核をもつ次の $A$-代数の
拡大が与えられているとする。
\begin{enumerate}
\item[(a)] $\xi \in \Ext^1_B(\NL_{B/A}, M)$ に対応する
$0 \to M \to B' \to B \to 0$；
\item[(b)] $\zeta \in \Ext^1_C(\NL_{C/A}, N)$ に対応する
$0 \to N \to C' \to C \to 0$。
\end{enumerate}
補題 \ref{defos-lemma-extensions-of-algebras} を参照せよ。このとき、
$B\to C$ および $c$ と両立する $A$-代数準同型 $B'\to C'$ が
存在するための必要十分条件は、$\xi$ と $\zeta$ が
$\Ext^1_B(\NL_{B/A}, N)$
の同じ元へ写ることである。
\end{lemma}

\begin{proof}
この主張が意味をもつことを確認する。写像 $M\to N$ により、
$$
\Ext^1_B(\NL_{B/A}, M) \to \Ext^1_B(\NL_{B/A}, N)
$$
があり、また
$$
\Ext^1_C(\NL_{C/A}, N) \to \Ext^1_B(\NL_{C/A}, N) \to \Ext^1_B(\NL_{B/A}, N)
$$
がある。ここで第一の矢印は制限写像 $D(C)\to D(B)$ を用い、第二の
矢印は複体の正準写像 $\NL_{B/A}\to\NL_{C/A}$ を用いる。本補題の
主張は、補題 \ref{defos-lemma-huge-diagram} を次の図式に適用し、
$$
\xymatrix{
& 0 \ar[r] & N \ar[r] & C' \ar[r] & C \ar[r] & 0 \\
& 0 \ar[r]|\hole & 0 \ar[u] \ar[r] &
A \ar[u] \ar[r]|\hole & A \ar[u] \ar[r] & 0 \\
0 \ar[r] & M \ar[ruu] \ar[r] & B' \ar[r] & B \ar[ruu] \ar[r] & 0 \\
0 \ar[r] & 0 \ar[ruu]|\hole \ar[u] \ar[r] &
A \ar[ruu]|\hole \ar[u] \ar[r] & A \ar[ruu]|\hole \ar[u] \ar[r] & 0
}
$$
補題 \ref{defos-lemma-extensions-of-algebras} と
\ref{defos-lemma-huge-diagram} の証明における構成の両立性を用いれば従う。
その両立性の主張と証明は省略する（直接的な議論については次の注意を見よ）。
\end{proof}

\begin{remark}
\label{defos-remark-extensions-of-algebras-functorial}
$A\to B\to C$、$M$、$N$、$c:M\to N$、
$0\to M\to B'\to B\to0$、$\xi\in\Ext^1_B(\NL_{B/A},M)$、
$0\to N\to C'\to C\to0$、および
$\zeta\in\Ext^1_C(\NL_{C/A},N)$ を補題
\ref{defos-lemma-extensions-of-algebras-functorial} のとおりとする。
$c:M\to N$ に沿う押し出しを用いて、次の拡大を構成できる。
$$
\xymatrix{
0 \ar[r] & N \ar[r] & B'_1 \ar[r] & B \ar[r] & 0 \\
0 \ar[r] & M \ar[u]^c \ar[r] & B' \ar[u] \ar[r] & B \ar[u] \ar[r] & 0
}
$$
ここで $M$ を反対角的に埋め込み、$B'_1=(N\times B')/M$ とおく。
$B\to C$ に沿う引き戻しを用いて、次の拡大を構成できる。
$$
\xymatrix{
0 \ar[r] & N \ar[r] & C' \ar[r] & C \ar[r] & 0 \\
0 \ar[r] & N \ar[u] \ar[r] & B'_2 \ar[u] \ar[r] & B \ar[u] \ar[r] & 0
}
$$
ここで $B'_2=C'\times_CB$（環のファイバー積）とおく。簡単な図式追跡に
より、$B\to C$ および $c$ と両立する $A$-代数準同型 $B'\to C'$ が
存在するための必要十分条件は、$B$ の $N$ による $A$-代数拡大として
$B'_1$ と $B'_2$ が同型であることだと分かる。したがって補題
\ref{defos-lemma-extensions-of-algebras-functorial} を示すには、補題
\ref{defos-lemma-extensions-of-algebras} の全単射により、$B'_1$ が写像
$\Ext^1_B(\NL_{B/A}, M) \to \Ext^1_B(\NL_{B/A}, N)$
による $\xi$ の像に対応し、$B'_2$ が写像
$\Ext^1_C(\NL_{C/A}, N) \to \Ext^1_B(\NL_{B/A}, N)$
による $\zeta$ の像に対応することを示せば十分である。前者は注意
\ref{defos-remark-extensions-of-algebras} における類の構成から直ちに従う。
後者については、$A$-代数の可換図式
$$
\xymatrix{
Q \ar[r]_\beta & C \\
P \ar[u]^\varphi \ar[r]^\alpha & B \ar[u]
}
$$
を選ぶ。ただし $\alpha$ は $A$ 上の $B$ の表示、$\beta$ は $A$ 上の
$C$ の表示である。注意 \ref{defos-remark-extensions-of-algebras} とそこでの
参考文献を参照せよ。$J=\Ker(\alpha)$、$K=\Ker(\beta)$ とおく。
写像 $\varphi$ は複体の写像 $\NL(\alpha)\to\NL(\beta)$、特に
$\bar\varphi:J/J^2\to K/K^2$ を誘導する。$\beta$ の持ち上げである
$A$-代数準同型 $\beta':Q\to C'$ を選ぶ。このとき
$\alpha' = (\beta' \circ \varphi, \alpha) : P \to B'_2 = C' \times_C B$
は $\alpha$ の持ち上げである。これらの選択のもとで、$\beta'$ が誘導する
写像 $K/K^2\to N$ と $\bar\varphi:J/J^2\to K/K^2$ の合成は、
$\alpha'$ の $J/J^2$ への制限である。注意
\ref{defos-remark-extensions-of-algebras} における類の構成を展開すれば、
$B'_2$ が写像
$\Ext^1_C(\NL_{C/A}, N) \to \Ext^1_B(\NL_{B/A}, N)$
による $\zeta$ の像に対応することが確かに分かる。
\end{remark}

\begin{lemma}
\label{defos-lemma-parametrize-solutions}
$0\to I\to A'\to A\to0$、$A\to B$、および $c:I\to N$ を
(\ref{defos-equation-to-solve}) のとおりとする。補題
\ref{defos-lemma-extensions-of-algebras} によって $A$ の $I$ による拡大 $A'$
に対応する元を $\xi\in\Ext^1_A(\NL_{A/A'},I)$ と書く。解の同型類全体の
集合は、写像
$$
\Ext^1_B(\NL_{B/A'}, N) \to \Ext^1_A(\NL_{A/A'}, N)
$$
の $\xi$ の像上のファイバーと正準的に全単射で対応する。
\end{lemma}

\begin{proof}
$A'\to B$ と $B$-加群 $N$ に補題
\ref{defos-lemma-extensions-of-algebras} を適用すると、
$\Ext^1_B(\NL_{B/A'},N)$ の元 $\zeta$ は $A'$-代数の拡大
$0\to N\to B'\to B\to0$ をパラメータ付ける。さらに
$A'\to A\to B$ と $c:I\to N$ に補題
\ref{defos-lemma-extensions-of-algebras-functorial} を適用すると、
$c$ および $A\to B$ と両立する $A'$-代数準同型 $A'\to B'$ が
存在するための必要十分条件は、$\zeta$ が $\xi$ に写ることである。
これはもちろん、$B'$ が (\ref{defos-equation-to-solve}) の解であることと
同じである。
\end{proof}

\begin{remark}
\label{defos-remark-parametrize-solutions}
補題 \ref{defos-lemma-parametrize-solutions} の状況では、
$$
\Ext^1_A(\NL_{A/A'}, N) =
\Ext^1_B(\NL_{A/A'} \otimes_A^\mathbf{L} B, N) =
\Ext^1_B(\NL_{A/A'} \otimes_A B, N)
$$
であることに注意する。第一の等式は『代数詳論』補題
\ref{more-algebra-lemma-tensor-hom-adjoint} により、第二の等式は
『代数詳論』補題 \ref{more-algebra-lemma-tensor-NL} による。
複体の写像
$$
\NL_{A/A'} \otimes_A B \to \NL_{B/A'} \to \NL_{B/A}
$$
があり、これは完全三角形に近い（『代数』補題
\ref{algebra-lemma-exact-sequence-NL} を見よ）。もしこれが完全三角形
であれば、$\Ext^2_B(\NL_{B/A},N)$ における $\xi$ の像が
(\ref{defos-equation-to-solve}) の解の存在に対する障害であると結論できる。
\end{remark}

\noindent
環準同型 $A\to B$ が局所完全交叉ならば解が存在する。これは一種の
持ち上げ定理である。シントミックな環準同型については、『環準同型の
平滑化』命題 \ref{smoothing-proposition-lift-smooth} で、かなり強い
持ち上げ定理をすでに証明していることに注意する。

\begin{lemma}
\label{defos-lemma-existence-lci}
$A\to B$ が局所完全交叉環準同型ならば、
(\ref{defos-equation-to-solve}) の解が存在する。
\end{lemma}

\begin{proof}[第一の証明]
$B=A[x_1,\ldots,x_n]/J$ と書く。『代数詳論』定義
\ref{more-algebra-definition-local-complete-intersection}
により、イデアル $J$ は Koszul 正則である。したがって $J$ は
$H_1$-正則かつ擬正則である（『代数詳論』第
\ref{more-algebra-section-ideals} 節を見よ）。
$J'\subset A'[x_1,\ldots,x_n]$ を $J$ の逆像とする。また、
$A'[x_1,\ldots,x_n]\to A[x_1,\ldots,x_n]$ の核を
$I[x_1,\ldots,x_n]$ と書く。『代数詳論』補題
\ref{more-algebra-lemma-conormal-sequence-H1-regular-ideal} により、
$I[x_1, \ldots, x_n] \cap (J')^2 = J'I[x_1, \ldots, x_n] =
JI[x_1, \ldots, x_n]$ である。したがって短完全列
$$
0 \to I \otimes_A B \to J'/(J')^2 \to J/J^2 \to 0
$$
を得る。$J/J^2$ は射影的なので（『代数詳論』補題
\ref{more-algebra-lemma-quasi-regular-ideal-finite-projective}）、この列の
分裂を選ぶことができる。
$$
J'/(J')^2 = I \otimes_A B \oplus J/J^2 
$$
$(J')^2\subset J''\subset J'$ を、上の分解の第二成分へ写る元全体とする。
このとき
$$
0 \to I \otimes_A B \to A'[x_1, \ldots, x_n]/J'' \to B \to 0
$$
は $N=I\otimes_AB$ の場合の (\ref{defos-equation-to-solve}) の解である。
一般の場合は、与えられた写像 $I\otimes_AB\to N$ に沿って押し出しを
とることにより得られる。
\end{proof}

\begin{proof}[第二の証明]
この証明を読む前に注意 \ref{defos-remark-parametrize-solutions} を読まれたい。
『代数詳論』補題 \ref{more-algebra-lemma-transitive-lci-at-end} により、
写像 $\NL_{A/A'}\otimes_AB\to\NL_{B/A'}\to\NL_{B/A}$ は
$D(B)$ において実際に完全三角形をなす。したがって
$\Ext^2_B(\NL_{B/A},N)$ が消滅することを示せば十分である。
『代数詳論』補題 \ref{more-algebra-lemma-lci-NL} により、複体
$\NL_{B/A}$ は Tor 振幅 $[-1,0]$ をもつ完全複体である。したがって、
例えば『代数詳論』補題 \ref{more-algebra-lemma-splitting-unique} (1)
により、この $\Ext^2$ は消滅する。
\end{proof}







\section{環付き空間の厚化}
\label{defos-section-thickenings-spaces}

\noindent
以下のいくつかの節では、次の概念を用いる。
\begin{enumerate}
\item 環付き空間 $(X',\mathcal{O}_{X'})$ 上のイデアル層
$\mathcal{I}\subset\mathcal{O}_{X'}$ が {\it 局所冪零} であるとは、
$\mathcal{I}$ の任意の局所切断が局所的に冪零であることをいう。
『代数』項 \ref{algebra-item-ideal-locally-nilpotent} と比較せよ。
\item 環付き空間の {\it 厚化} とは、環付き空間の射
$i:(X,\mathcal{O}_X)\to(X',\mathcal{O}_{X'})$ であって、
次を満たすものをいう。
\begin{enumerate}
\item $i$ は同相写像 $X\to X'$ を誘導する；
\item 写像 $i^\sharp:\mathcal{O}_{X'}\to i_*\mathcal{O}_X$ は全射である；
\item $i^\sharp$ の核は局所冪零なイデアル層である。
\end{enumerate}
\item 環付き空間の {\it 一次厚化} とは、$\Ker(i^\sharp)$ が二乗零である
環付き空間の厚化
$i:(X,\mathcal{O}_X)\to(X',\mathcal{O}_{X'})$ をいう。
\item {\it 厚化の射}、{\it 基礎環付き空間上の厚化の射} などの定義は
明らかである。
\end{enumerate}
もし $i:(X,\mathcal{O}_X)\to(X',\mathcal{O}_{X'})$ が環付き空間の
厚化ならば、台となる位相空間を同一視し、$\mathcal{O}_X$、
$\mathcal{O}_{X'}$、および $\mathcal{I}=\Ker(i^\sharp)$ を
$X=X'$ 上の層とみなす。このとき短完全列
$$
0 \to \mathcal{I} \to \mathcal{O}_{X'} \to \mathcal{O}_X \to 0
$$
を $\mathcal{O}_{X'}$-加群の列として得る。『加群』補題
\ref{modules-lemma-i-star-equivalence} により、$\mathcal{O}_X$-加群の
圏は、$\mathcal{I}$ によって零化される $\mathcal{O}_{X'}$-加群の圏と
同値である。特に $i$ が一次厚化ならば、$\mathcal{I}$ は
$\mathcal{O}_X$-加群である。

\begin{situation}
\label{defos-situation-morphism-thickenings}
厚化の射 $(f,f')$ は、次の可換図式で与えられる。
\begin{equation}
\label{defos-equation-morphism-thickenings}
\vcenter{
\xymatrix{
(X, \mathcal{O}_X) \ar[r]_i \ar[d]_f & (X', \mathcal{O}_{X'}) \ar[d]^{f'} \\
(S, \mathcal{O}_S) \ar[r]^t & (S', \mathcal{O}_{S'})
}
}
\end{equation}
ここで横の矢印は環付き空間の厚化である。この状況で
$\mathcal{I}=\Ker(i^\sharp)\subset\mathcal{O}_{X'}$ および
$\mathcal{J}=\Ker(t^\sharp)\subset\mathcal{O}_{S'}$ とおく。
台となる位相空間上では $f=f'$ なので、位相的引き戻し関手
$f^{-1}$ と $(f')^{-1}$ を同一視する。
$(f')^\sharp:f^{-1}\mathcal{O}_{S'}\to\mathcal{O}_{X'}$ は特に写像
$f^{-1}\mathcal{J}\to\mathcal{I}$ を誘導し、したがって
$\mathcal{O}_{X'}$-加群の写像
$$
(f')^*\mathcal{J} \longrightarrow \mathcal{I}
$$
を誘導することに注意する。$i$ と $t$ が一次厚化ならば、
$(f')^*\mathcal{J}=f^*\mathcal{J}$ であり、上の写像は
$f^*\mathcal{J}\to\mathcal{I}$ となる。
\end{situation}

\begin{definition}
\label{defos-definition-strict-morphism-thickenings}
状況 \ref{defos-situation-morphism-thickenings} において、写像
$(f')^*\mathcal{J}\longrightarrow\mathcal{I}$ が全射であるとき、
$(f,f')$ を {\it 厳密な厚化の射} と呼ぶ。
\end{definition}

\noindent
次の補題から特に、スキームの厚化の射
$(f,f'):(X\subset X')\to(S\subset S')$ が厳密であるための
必要十分条件は $X=S\times_{S'}X'$ であることが分かる。

\begin{lemma}
\label{defos-lemma-strict-morphism-thickenings}
状況 \ref{defos-situation-morphism-thickenings} において、射 $(f,f')$ が
厳密な厚化の射であるための必要十分条件は、
(\ref{defos-equation-morphism-thickenings}) が環付き空間の圏で
デカルト的であることである。
\end{lemma}

\begin{proof}
省略する。
\end{proof}




\section{環付き空間の一次厚化上の加群}
\label{defos-section-modules-thickenings}

\noindent
本節では、加群の変形理論に必要ないくつかの準備事項を論じる。
$i : (X, \mathcal{O}_X) \to (X', \mathcal{O}_{X'})$ を環付き空間の
一次厚化とする。節 \ref{defos-section-thickenings-spaces} で導入した記法を
断りなく用い、特に台となる位相空間を同一視する。
本節では、次の短完全列を考える。
\begin{equation}
\label{defos-equation-extension}
0 \to \mathcal{K} \to \mathcal{F}' \to \mathcal{F} \to 0
\end{equation}
これは $\mathcal{O}_{X'}$-加群の列であり、$\mathcal{F}$ と
$\mathcal{K}$ は $\mathcal{O}_X$-加群、$\mathcal{F}'$ は
$\mathcal{O}_{X'}$-加群である。この状況では、正準な
$\mathcal{O}_X$-加群準同型
$$
c_{\mathcal{F}'} :
\mathcal{I} \otimes_{\mathcal{O}_X} \mathcal{F}
\longrightarrow
\mathcal{K}
$$
がある。ここで $\mathcal{I} = \Ker(i^\sharp)$ である。
実際、$\mathcal{I}$ の局所切断 $f$ と $\mathcal{F}$ の局所切断 $s$
に対し、$s$ を持ち上げる $\mathcal{F}'$ の局所切断を $s'$ として
$c_{\mathcal{F}'}(f \otimes s) = fs'$ と定める。

\begin{lemma}
\label{defos-lemma-inf-map}
$i : (X, \mathcal{O}_X) \to (X', \mathcal{O}_{X'})$ を環付き空間の
一次厚化とする。次の拡大が与えられていると仮定する。
$$
0 \to \mathcal{K} \to \mathcal{F}' \to \mathcal{F} \to 0
\quad\text{and}\quad
0 \to \mathcal{L} \to \mathcal{G}' \to \mathcal{G} \to 0
$$
これらは (\ref{defos-equation-extension}) の形であり、さらに写像
$\varphi : \mathcal{F} \to \mathcal{G}$ と
$\psi : \mathcal{K} \to \mathcal{L}$ が与えられているとする。
\begin{enumerate}
\item $\varphi$ および $\psi$ と両立する $\mathcal{O}_{X'}$-加群準同型
$\varphi' : \mathcal{F}' \to \mathcal{G}'$ が存在すれば、次の図式は
可換である。
$$
\xymatrix{
\mathcal{I} \otimes_{\mathcal{O}_X} \mathcal{F}
\ar[r]_-{c_{\mathcal{F}'}} \ar[d]_{1 \otimes \varphi} &
\mathcal{K} \ar[d]^\psi \\
\mathcal{I} \otimes_{\mathcal{O}_X} \mathcal{G}
\ar[r]^-{c_{\mathcal{G}'}} &
\mathcal{L}
}
$$
\item $\varphi$ および $\psi$ と両立する $\mathcal{O}_{X'}$-加群準同型
$\varphi' : \mathcal{F}' \to \mathcal{G}'$ 全体の集合は、空でなければ
$\Hom_{\mathcal{O}_X}(\mathcal{F}, \mathcal{L})$ の下で主等質空間をなす。
\end{enumerate}
\end{lemma}

\begin{proof}
(1) は写像の記述から直ちに従う。
(2) について、$\varphi'$ と $\varphi''$ が $\varphi$ および $\psi$ と
両立する二つの写像 $\mathcal{F}' \to \mathcal{G}'$ ならば、
$\varphi' - \varphi''$ は次のように分解する。
$$
\mathcal{F}' \to \mathcal{F} \to \mathcal{L} \to \mathcal{G}'
$$
『加群』補題 \ref{modules-lemma-i-star-equivalence} により、中間の写像は
$\Hom_{\mathcal{O}_X}(\mathcal{F}, \mathcal{L})$ の一意な元から生じる。
逆に、この群の元 $\alpha$ が与えられれば、上の表示で中間に $\alpha$
を置いた合成を $\varphi'$ に加えることができる。細部は省略する。
\end{proof}

\begin{lemma}
\label{defos-lemma-inf-obs-map}
$i : (X, \mathcal{O}_X) \to (X', \mathcal{O}_{X'})$ を環付き空間の
一次厚化とする。次の拡大が与えられていると仮定する。
$$
0 \to \mathcal{K} \to \mathcal{F}' \to \mathcal{F} \to 0
\quad\text{and}\quad
0 \to \mathcal{L} \to \mathcal{G}' \to \mathcal{G} \to 0
$$
これらは (\ref{defos-equation-extension}) の形であり、さらに写像
$\varphi : \mathcal{F} \to \mathcal{G}$ と
$\psi : \mathcal{K} \to \mathcal{L}$ が与えられているとする。図式
$$
\xymatrix{
\mathcal{I} \otimes_{\mathcal{O}_X} \mathcal{F}
\ar[r]_-{c_{\mathcal{F}'}} \ar[d]_{1 \otimes \varphi} &
\mathcal{K} \ar[d]^\psi \\
\mathcal{I} \otimes_{\mathcal{O}_X} \mathcal{G}
\ar[r]^-{c_{\mathcal{G}'}} &
\mathcal{L}
}
$$
が可換であると仮定する。このとき、元
$$
o(\varphi, \psi) \in
\Ext^1_{\mathcal{O}_X}(\mathcal{F}, \mathcal{L})
$$
が存在し、その消滅は、$\varphi$ および $\psi$ と両立する写像
$\varphi' : \mathcal{F}' \to \mathcal{G}'$ が存在するための
必要十分条件である。
\end{lemma}

\begin{proof}
次の拡大を明示的に構成できる。
$$
0 \to \mathcal{L} \to \mathcal{H} \to \mathcal{F} \to 0
$$
$\mathcal{H}$ を複体
$$
\mathcal{K}
\xrightarrow{1, - \psi}
\mathcal{F}' \oplus \mathcal{G}' \xrightarrow{\varphi, 1}
\mathcal{G}
$$
の中間のコホモロジーとすればよい（記法は明らかであろう）。
補題の図式が可換であるという仮定を用いた局所切断の計算により、
$\mathcal{H}$ は $\mathcal{I}$ によって零化される。したがって
$\mathcal{H}$ は次の群の類を定める。
$$
\Ext^1_{\mathcal{O}_X}(\mathcal{F}, \mathcal{L})
\subset
\Ext^1_{\mathcal{O}_{X'}}(\mathcal{F}, \mathcal{L})
$$
最後に、$\mathcal{H}$ の類は、拡大 $\mathcal{F}'$ の $\psi$ に沿う
押し出しと、拡大 $\mathcal{G}'$ の $\varphi$ に沿う引き戻しとの差である
（計算は省略する）。したがって $\mathcal{H}$ の類が消滅することは、
次の可換図式が存在することと同値である。
$$
\xymatrix{
0 \ar[r] &
\mathcal{K} \ar[r] \ar[d]_{\psi} &
\mathcal{F}' \ar[r] \ar[d]_{\varphi'} &
\mathcal{F} \ar[r] \ar[d]_\varphi & 0\\
0 \ar[r] &
\mathcal{L} \ar[r] &
\mathcal{G}' \ar[r] &
\mathcal{G} \ar[r] & 0
}
$$
これは所望の結論である。
\end{proof}

\begin{lemma}
\label{defos-lemma-inf-ext}
$i : (X, \mathcal{O}_X) \to (X', \mathcal{O}_{X'})$ を環付き空間の
一次厚化とする。$\mathcal{O}_X$-加群 $\mathcal{F}$、$\mathcal{K}$ と
$\mathcal{O}_X$-線形写像
$c : \mathcal{I} \otimes_{\mathcal{O}_X} \mathcal{F} \to \mathcal{K}$
が与えられているとする。$c_{\mathcal{F}'} = c$ を満たす列
(\ref{defos-equation-extension}) が存在するならば、そのような拡大の
同型類全体の集合は
$\Ext^1_{\mathcal{O}_X}(\mathcal{F}, \mathcal{K})$ の下で主等質空間をなす。
\end{lemma}

\begin{proof}
次の拡大が与えられていると仮定する。
$$
0 \to \mathcal{K} \to \mathcal{F}'_1 \to \mathcal{F} \to 0
\quad\text{and}\quad
0 \to \mathcal{K} \to \mathcal{F}'_2 \to \mathcal{F} \to 0
$$
$c_{\mathcal{F}'_1} = c_{\mathcal{F}'_2} = c$ とする。このとき、
（拡大群における）差（『ホモロジー』節
\ref{homology-section-extensions} を見よ）は拡大
$$
0 \to \mathcal{K} \to \mathcal{E} \to \mathcal{F} \to 0
$$
であり、$\mathcal{E}$ は $\mathcal{I}$ によって零化される
（局所計算は省略する）。したがってこの列は $\mathcal{O}_X$-加群の
拡大である。『加群』補題 \ref{modules-lemma-i-star-equivalence} を見よ。
逆に、そのような拡大 $\mathcal{E}$ が与えられれば、写像
$c_{\mathcal{F}'}$ を変えずに、拡大 $\mathcal{E}$ を
$\mathcal{O}_{X'}$-拡大 $\mathcal{F}'$ に加えることができる。
細部は省略する。
\end{proof}

\begin{lemma}
\label{defos-lemma-inf-obs-ext}
$i : (X, \mathcal{O}_X) \to (X', \mathcal{O}_{X'})$ を環付き空間の
一次厚化とする。$\mathcal{O}_X$-加群 $\mathcal{F}$、$\mathcal{K}$ と
$\mathcal{O}_X$-線形写像
$c : \mathcal{I} \otimes_{\mathcal{O}_X} \mathcal{F} \to \mathcal{K}$
が与えられているとする。このとき元
$$
o(\mathcal{F}, \mathcal{K}, c) \in
\Ext^2_{\mathcal{O}_X}(\mathcal{F}, \mathcal{K})
$$
が存在し、その消滅は、$c_{\mathcal{F}'} = c$ を満たす列
(\ref{defos-equation-extension}) が存在するための必要十分条件である。
\end{lemma}

\begin{proof}
まず、$\mathcal{K}$ が入射 $\mathcal{O}_X$-加群ならば、
$c_{\mathcal{F}'} = c$ を満たす列 (\ref{defos-equation-extension}) が実際に
存在することを示す。そのため、平坦 $\mathcal{O}_{X'}$-加群
$\mathcal{H}'$ と全射 $\mathcal{H}' \to \mathcal{F}$ を選ぶ
（『加群』補題 \ref{modules-lemma-module-quotient-flat}）。
その核を $\mathcal{J} \subset \mathcal{H}'$ とする。
$\mathcal{H}'$ は平坦なので、
$$
\mathcal{I} \otimes_{\mathcal{O}_{X'}} \mathcal{H}' =
\mathcal{I}\mathcal{H}'
\subset \mathcal{J} \subset \mathcal{H}'
$$
写像
$$
\mathcal{I}\mathcal{H}' =
\mathcal{I} \otimes_{\mathcal{O}_{X'}} \mathcal{H}'
\longrightarrow
\mathcal{I} \otimes_{\mathcal{O}_{X'}} \mathcal{F} =
\mathcal{I} \otimes_{\mathcal{O}_X} \mathcal{F}
$$
が $\mathcal{I}\mathcal{J}$ を零化することに注意する。実際、$f$ が
$\mathcal{I}$ の局所切断、$s$ が $\mathcal{H}$ の局所切断ならば、
$fs$ は $f \otimes \overline{s}$ に写る。ここで $\overline{s}$ は
$s$ の $\mathcal{F}$ における像である。したがって、
$$
\xymatrix{
\mathcal{I}\mathcal{H}'/\mathcal{I}\mathcal{J}
\ar@{^{(}->}[r] \ar[d] &
\mathcal{J}/\mathcal{I}\mathcal{J} \ar@{..>}[d]_\gamma \\
\mathcal{I} \otimes_{\mathcal{O}_X} \mathcal{F} \ar[r]^-c &
\mathcal{K}
}
$$
という $\mathcal{O}_X$-加群の図式を得る。$\mathcal{K}$ が
$\mathcal{O}_X$-加群として入射ならば、点線の矢印が得られる。
$\gamma$ と $\mathcal{J} \to \mathcal{J}/\mathcal{I}\mathcal{J}$ の合成を
$\gamma' : \mathcal{J} \to \mathcal{K}$ と記す。局所計算により、押し出し
$$
\xymatrix{
0 \ar[r] &
\mathcal{J} \ar[r] \ar[d]_{\gamma'} &
\mathcal{H}' \ar[r] \ar[d] &
\mathcal{F} \ar[r] \ar@{=}[d] &
0 \\
0 \ar[r] &
\mathcal{K} \ar[r] &
\mathcal{F}' \ar[r] &
\mathcal{F} \ar[r] &
0
}
$$
は補題で課された問題の解である。

\medskip\noindent
一般の場合。$\mathcal{K}'$ が入射 $\mathcal{O}_X$-加群となる埋め込み
$\mathcal{K} \subset \mathcal{K}'$ を選ぶ。商を $\mathcal{Q}$ とすると、
完全列
$$
0 \to \mathcal{K} \to \mathcal{K}' \to \mathcal{Q} \to 0
$$
を得る。合成を
$c' : \mathcal{I} \otimes_{\mathcal{O}_X} \mathcal{F} \to \mathcal{K}'$
と記す。上の段落により、(\ref{defos-equation-extension}) の形の列
$$
0 \to \mathcal{K}' \to \mathcal{E}' \to \mathcal{F} \to 0
$$
で $c_{\mathcal{E}'} = c'$ を満たすものが存在する。
$c'$ と写像 $\mathcal{K}' \to \mathcal{Q}$ の合成は零であることに
注意する。したがって $\mathcal{E}'$ の
$\mathcal{K}' \to \mathcal{Q}$ に沿う押し出しは拡大
$$
0 \to \mathcal{Q} \to \mathcal{D}' \to \mathcal{F} \to 0
$$
であり、(\ref{defos-equation-extension}) の形で
$c_{\mathcal{D}'} = 0$ を満たす。これはちょうど $\mathcal{D}'$ が
$\mathcal{I}$ によって零化されることを意味する。言い換えれば、
$\mathcal{D}'$ は $\mathcal{O}_X$-加群の拡大、すなわち元
$$
o(\mathcal{F}, \mathcal{K}, c) \in
\Ext^1_{\mathcal{O}_X}(\mathcal{F}, \mathcal{Q}) =
\Ext^2_{\mathcal{O}_X}(\mathcal{F}, \mathcal{K})
$$
を定める（等号は、上の完全列に付随する長完全コホモロジー列と、
入射加群 $\mathcal{K}'$ を値域とする高次 Ext 群の消滅から従う）。
$o(\mathcal{F}, \mathcal{K}, c) = 0$ ならば、分裂
$s : \mathcal{F} \to \mathcal{D}'$ を選び、
$$
\mathcal{F}' = \Ker(\mathcal{E}' \to \mathcal{D}'/s(\mathcal{F}))
$$
とおくことができる。すると、行が完全な次の図式を得る。
$$
\xymatrix{
0 \ar[r] &
\mathcal{K} \ar[r] \ar[d] &
\mathcal{F}' \ar[r] \ar[d] &
\mathcal{F} \ar[r] \ar@{=}[d] &
0 \\
0 \ar[r] &
\mathcal{K}' \ar[r] &
\mathcal{E}' \ar[r] &
\mathcal{F} \ar[r] & 0
}
$$
これにより $c_{\mathcal{F}'} = c$ が分かる。逆に $\mathcal{F}'$ が
存在するならば、補題 \ref{defos-lemma-inf-ext} と、入射加群 $\mathcal{K}'$
を値域とする高次 Ext 群の消滅により、$\mathcal{F}'$ の写像
$\mathcal{K} \to \mathcal{K}'$ に沿う押し出しは $\mathcal{E}'$ と同型である。
これから上のような図式が得られ、$\mathcal{D}'$ は拡大として分裂する。
すなわち類 $o(\mathcal{F}, \mathcal{K}, c)$ は零である。
\end{proof}

\begin{remark}
\label{defos-remark-trivial-thickening}
$(X, \mathcal{O}_X)$ を環付き空間とする。環付き空間の射
$\pi : (X', \mathcal{O}_{X'}) \to (X, \mathcal{O}_X)$ で $i$ の
左逆となるものが存在するとき、一次厚化
$i : (X, \mathcal{O}_X) \to (X', \mathcal{O}_{X'})$ を
{\it 自明} であるという。そのような射 $\pi$ の選択を、一次厚化の
{\it 自明化} と呼ぶ。$\pi$ が与えられると、分裂
\begin{equation}
\label{defos-equation-splitting}
\mathcal{O}_{X'} = \mathcal{O}_X \oplus \mathcal{I}
\end{equation}
を $X$ 上の代数層として得る。これは $\pi^\sharp$ を用いて全射
$\mathcal{O}_{X'} \to \mathcal{O}_X$ を分裂させたものである。逆に、
そのような分裂は射 $\pi$ を定める。$(X, \mathcal{O}_X)$ の自明化された
一次厚化の圏は、$\mathcal{O}_X$-加群の圏と同値である。
\end{remark}

\begin{remark}
\label{defos-remark-trivial-extension}
$i : (X, \mathcal{O}_X) \to (X', \mathcal{O}_{X'})$ を環付き空間の
自明な一次厚化とし、
$\pi : (X', \mathcal{O}_{X'}) \to (X, \mathcal{O}_X)$ をその自明化とする。
二つの $\mathcal{O}_X$-加群と写像
$c : \mathcal{I} \otimes_{\mathcal{O}_X} \mathcal{F} \to \mathcal{K}$
からなる任意の三つ組 $(\mathcal{F}, \mathcal{K}, c)$ に対して、
$$
\mathcal{F}'_{c, triv} = \mathcal{F} \oplus \mathcal{K}
$$
とおき、$\pi$ に付随する分裂 (\ref{defos-equation-splitting}) と写像 $c$ を
用いて $\mathcal{O}_{X'}$-加群構造を定め、拡大
(\ref{defos-equation-extension}) を得ることができる。
$\mathcal{F}'_{c, triv}$ を、$c$ と自明化 $\pi$ に対応する
$\mathcal{F}$ の $\mathcal{K}$ による {\it 自明な拡大} と呼ぶ。
(\ref{defos-equation-extension}) の形の任意の拡大 $\mathcal{F}'$ に対し、
$\pi^\sharp : \mathcal{O}_X \to \mathcal{O}_{X'}$ を用いて
$\mathcal{F}'$ を $\mathcal{O}_X$-加群の拡大とみなせるので、類
$\xi_{\mathcal{F}'}$ を
$\Ext^1_{\mathcal{O}_X}(\mathcal{F}, \mathcal{K})$ に得る。
補題 \ref{defos-lemma-inf-ext} により、
$\mathcal{F}' \mapsto \xi_{\mathcal{F}'}$
は全単射
$$
\left\{
\begin{matrix}
\text{拡大の同型類}\\
\mathcal{F}'\text{ は (\ref{defos-equation-extension}) の形で、}c = c_{\mathcal{F}'}
\end{matrix}
\right\}
\longrightarrow
\Ext^1_{\mathcal{O}_X}(\mathcal{F}, \mathcal{K})
$$
を誘導する。さらに、自明な拡大 $\mathcal{F}'_{c, triv}$ は零類に写る。
\end{remark}

\begin{remark}
\label{defos-remark-extension-functorial}
$(X, \mathcal{O}_X)$ を環付き空間とする。
$(X, \mathcal{O}_X) \to (X'_i, \mathcal{O}_{X'_i})$、$i = 1, 2$ を、
イデアル層 $\mathcal{I}_i$ をもつ一次厚化とする。また、
$h : (X'_1, \mathcal{O}_{X'_1}) \to (X'_2, \mathcal{O}_{X'_2})$ を
$(X, \mathcal{O}_X)$ の一次厚化の射とする。図示すれば、
$$
\xymatrix{
& (X, \mathcal{O}_X) \ar[ld] \ar[rd] & \\
(X'_1, \mathcal{O}_{X'_1}) \ar[rr]^h & & 
(X'_2, \mathcal{O}_{X'_2})
}
$$
である。$h^\sharp : \mathcal{O}_{X'_2} \to \mathcal{O}_{X'_1}$ は、特に
$\mathcal{O}_X$-加群準同型 $\mathcal{I}_2 \to \mathcal{I}_1$ を誘導する
ことに注意する。$\mathcal{F}$ を $\mathcal{O}_X$-加群とする。
$(\mathcal{K}_i, c_i)$、$i = 1, 2$ を、それぞれ
$\mathcal{O}_X$-加群 $\mathcal{K}_i$ と写像
$c_i : \mathcal{I}_i \otimes_{\mathcal{O}_X} \mathcal{F} \to
\mathcal{K}_i$ からなる組とする。さらに、次の図式を可換にする
$\mathcal{O}_X$-加群準同型 $\mathcal{K}_2 \to \mathcal{K}_1$ が
与えられていると仮定する。
$$
\xymatrix{
\mathcal{I}_2 \otimes_{\mathcal{O}_X} \mathcal{F}
\ar[r]_-{c_2} \ar[d] &
\mathcal{K}_2 \ar[d] \\
\mathcal{I}_1 \otimes_{\mathcal{O}_X} \mathcal{F}
\ar[r]^-{c_1} &
\mathcal{K}_1
}
$$
このとき、正準な関手性
$$
\left\{
\begin{matrix}
\mathcal{F}'_2\text{ は (\ref{defos-equation-extension}) の形で、}\\
c_2 = c_{\mathcal{F}'_2}\text{、かつ }\mathcal{K} = \mathcal{K}_2
\end{matrix}
\right\}
\longrightarrow
\left\{
\begin{matrix}
\mathcal{F}'_1\text{ は (\ref{defos-equation-extension}) の形で、}\\
c_1 = c_{\mathcal{F}'_1}\text{、かつ }\mathcal{K} = \mathcal{K}_1
\end{matrix}
\right\}
$$
がある。実際、すべての層 $\mathcal{O}_X$、$\mathcal{O}_{X'_i}$、
$\mathcal{F}$、$\mathcal{K}_i$ などを $X$ 上の層とみなし、
$\mathcal{F}'_2$ が与えられたとき、$\mathcal{F}'_1$ を押し出し、
すなわち次の拡大の図式に収まる層とする。
$$
\xymatrix{
0 \ar[r] &
\mathcal{K}_2 \ar[r] \ar[d] &
\mathcal{F}'_2 \ar[r] \ar[d] &
\mathcal{F} \ar@{=}[d] \ar[r] & 0 \\
0 \ar[r] &
\mathcal{K}_1 \ar[r] &
\mathcal{F}'_1 \ar[r] &
\mathcal{F} \ar[r] & 0
}
$$
押し出し上の $\mathcal{O}_{X'_1}$-加群構造の構成は省略する
（これには $c_1$ と $c_2$ を含む図式の可換性を用いる）。
\end{remark}

\begin{remark}
\label{defos-remark-trivial-extension-functorial}
$(X, \mathcal{O}_X)$、$(X, \mathcal{O}_X) \to (X'_i, \mathcal{O}_{X'_i})$、
$\mathcal{I}_i$、および
$h : (X'_1, \mathcal{O}_{X'_1}) \to (X'_2, \mathcal{O}_{X'_2})$ は、
注意 \ref{defos-remark-extension-functorial} と同様とする。
$\pi_1 = h \circ \pi_2$ を満たす自明化 $\pi_i : X'_i \to X$ が
与えられていると仮定する。言い換えれば、$h$ は
$(X, \mathcal{O}_X)$ の自明化された一次厚化の射であると仮定する。
$(\mathcal{K}_i, c_i)$、$i = 1, 2$ を、$\mathcal{O}_X$-加群
$\mathcal{K}_i$ と写像
$c_i : \mathcal{I}_i \otimes_{\mathcal{O}_X} \mathcal{F} \to
\mathcal{K}_i$ からなる組とする。さらに、次の図式を可換にする
$\mathcal{O}_X$-加群準同型 $\mathcal{K}_2 \to \mathcal{K}_1$ が
与えられていると仮定する。
$$
\xymatrix{
\mathcal{I}_2 \otimes_{\mathcal{O}_X} \mathcal{F}
\ar[r]_-{c_2} \ar[d] &
\mathcal{K}_2 \ar[d] \\
\mathcal{I}_1 \otimes_{\mathcal{O}_X} \mathcal{F}
\ar[r]^-{c_1} &
\mathcal{K}_1
}
$$
この状況では、注意 \ref{defos-remark-trivial-extension} の構成は
可換図式
$$
\xymatrix{
\{\mathcal{F}'_2\text{ は (\ref{defos-equation-extension}) の形で、}
c_2 = c_{\mathcal{F}'_2}\text{、かつ }\mathcal{K} = \mathcal{K}_2\}
\ar[d] \ar[rr] & &
\Ext^1_{\mathcal{O}_X}(\mathcal{F}, \mathcal{K}_2) \ar[d] \\
\{\mathcal{F}'_1\text{ は (\ref{defos-equation-extension}) の形で、}
c_1 = c_{\mathcal{F}'_1}\text{、かつ }\mathcal{K} = \mathcal{K}_1\}
\ar[rr] & &
\Ext^1_{\mathcal{O}_X}(\mathcal{F}, \mathcal{K}_1)
}
$$
を誘導する。ここで右側の垂直写像は $\Ext$ の関手性と写像
$\mathcal{K}_2 \to \mathcal{K}_1$ によって与えられ、左側の垂直写像は
注意 \ref{defos-remark-extension-functorial} の写像である。
\end{remark}

\begin{remark}
\label{defos-remark-short-exact-sequence-thickenings}
$(X, \mathcal{O}_X)$ を環付き空間とする。$(X, \mathcal{O}_X)$ の
一次厚化の射の列
$$
(X'_1, \mathcal{O}_{X'_1}) \to
(X'_2, \mathcal{O}_{X'_2}) \to
(X'_3, \mathcal{O}_{X'_3})
$$
を考える。対応するイデアル層 $\mathcal{I}_i$ の間の写像が
$\mathcal{O}_X$-加群の複体
$\mathcal{I}_3 \to \mathcal{I}_2 \to \mathcal{I}_1$、すなわち合成が零、
を与えるとき、この列を {\it 複体} と呼ぶ。この場合、合成
$(X'_1, \mathcal{O}_{X'_1}) \to (X_3', \mathcal{O}_{X'_3})$ は
$(X, \mathcal{O}_X) \to (X'_3, \mathcal{O}_{X'_3})$ を経由する。
すなわち、$(X, \mathcal{O}_X)$ の一次厚化
$(X'_1, \mathcal{O}_{X'_1})$ は自明であり、正準な自明化
$\pi : (X'_1, \mathcal{O}_{X'_1}) \to (X, \mathcal{O}_X)$
を伴う。

\medskip\noindent
$(X, \mathcal{O}_X)$ の一次厚化の射の列
$$
(X'_1, \mathcal{O}_{X'_1}) \to
(X'_2, \mathcal{O}_{X'_2}) \to
(X'_3, \mathcal{O}_{X'_3})
$$
について、対応するイデアル層の間の写像が $\mathcal{O}_X$-加群の
短完全列
$$
0 \to \mathcal{I}_3 \to \mathcal{I}_2 \to \mathcal{I}_1 \to 0
$$
をなすとき、これを {\it 短完全列} と呼ぶ。
\end{remark}

\begin{remark}
\label{defos-remark-complex-thickenings-and-ses-modules}
$(X, \mathcal{O}_X)$ を環付き空間とし、$\mathcal{F}$ を
$\mathcal{O}_X$-加群とする。
$$
(X'_1, \mathcal{O}_{X'_1}) \to
(X'_2, \mathcal{O}_{X'_2}) \to
(X'_3, \mathcal{O}_{X'_3})
$$
を $(X, \mathcal{O}_X)$ の一次厚化の複体とする。注意
\ref{defos-remark-short-exact-sequence-thickenings} を見よ。
$(\mathcal{K}_i, c_i)$、$i = 1, 2, 3$ を、それぞれ
$\mathcal{O}_X$-加群 $\mathcal{K}_i$ と写像
$c_i : \mathcal{I}_i \otimes_{\mathcal{O}_X} \mathcal{F} \to
\mathcal{K}_i$ からなる組とする。$\mathcal{O}_X$-加群の短完全列
$$
0 \to \mathcal{K}_3 \to \mathcal{K}_2 \to \mathcal{K}_1 \to 0
$$
が与えられ、次の二つの図式が可換であると仮定する。
$$
\vcenter{
\xymatrix{
\mathcal{I}_2 \otimes_{\mathcal{O}_X} \mathcal{F}
\ar[r]_-{c_2} \ar[d] &
\mathcal{K}_2 \ar[d] \\
\mathcal{I}_1 \otimes_{\mathcal{O}_X} \mathcal{F}
\ar[r]^-{c_1} &
\mathcal{K}_1
}
}
\quad\text{and}\quad
\vcenter{
\xymatrix{
\mathcal{I}_3 \otimes_{\mathcal{O}_X} \mathcal{F}
\ar[r]_-{c_3} \ar[d] &
\mathcal{K}_3 \ar[d] \\
\mathcal{I}_2 \otimes_{\mathcal{O}_X} \mathcal{F}
\ar[r]^-{c_2} &
\mathcal{K}_2
}
}
$$
最後に、$\mathcal{O}_{X'_2}$-加群の拡大
$$
0 \to \mathcal{K}_2 \to \mathcal{F}'_2 \to \mathcal{F} \to 0
$$
が与えられていると仮定する。これは (\ref{defos-equation-extension}) の形で
$\mathcal{K} = \mathcal{K}_2$ および $c_{\mathcal{F}'_2} = c_2$ を満たす。
この状況で、注意 \ref{defos-remark-extension-functorial} の関手性を適用して
$X'_1$ 上の拡大 $\mathcal{F}'_1$ を得る（この特別な場合の
$\mathcal{F}'_1$ は後で記述する）。また注意
\ref{defos-remark-trivial-extension} により、注意
\ref{defos-remark-short-exact-sequence-thickenings} の正準な分裂
$\pi : (X'_1, \mathcal{O}_{X'_1}) \to (X, \mathcal{O}_X)$ を用いて、
$\xi_{\mathcal{F}'_1} \in
\Ext^1_{\mathcal{O}_X}(\mathcal{F}, \mathcal{K}_1)$ を得る。
最後に、障害
$$
o(\mathcal{F}, \mathcal{K}_3, c_3) \in
\Ext^2_{\mathcal{O}_X}(\mathcal{F}, \mathcal{K}_3)
$$
がある。補題 \ref{defos-lemma-inf-obs-ext} を見よ。この状況では、短完全列
$0 \to \mathcal{K}_3 \to \mathcal{K}_2 \to \mathcal{K}_1 \to 0$
から生じる正準写像
$$
\partial :
\Ext^1_{\mathcal{O}_X}(\mathcal{F}, \mathcal{K}_1)
\longrightarrow
\Ext^2_{\mathcal{O}_X}(\mathcal{F}, \mathcal{K}_3)
$$
が $\xi_{\mathcal{F}'_1}$ を障害類
$o(\mathcal{F}, \mathcal{K}_3, c_3)$ に写すと {\bf 主張する}。

\medskip\noindent
この主張を証明するため、$\mathcal{K}$ が入射 $\mathcal{O}_X$-加群となる
埋め込み $j : \mathcal{K}_3 \to \mathcal{K}$ を選ぶ。
$j$ を写像 $j' : \mathcal{K}_2 \to \mathcal{K}$ に持ち上げることができる。
$\mathcal{E}'_2 = j'_*\mathcal{F}'_2$ を、$\mathcal{F}'_2$ の $j'$ に沿う
押し出しとする。このとき $c_{\mathcal{E}'_2} = j' \circ c_2$ である。
図示すれば、
$$
\xymatrix{
0 \ar[r] &
\mathcal{K}_2 \ar[r] \ar[d]_{j'} &
\mathcal{F}'_2 \ar[r] \ar[d] &
\mathcal{F} \ar[r] \ar[d] & 0 \\
0 \ar[r] &
\mathcal{K} \ar[r] &
\mathcal{E}'_2 \ar[r] &
\mathcal{F} \ar[r] & 0
}
$$
$\mathcal{E}'_3 = \mathcal{E}'_2$ とおく。ただし、
$\mathcal{O}_{X'_3} \to \mathcal{O}_{X'_2}$ を通じて
$\mathcal{O}_{X'_3}$-加群とみなす。このとき
$c_{\mathcal{E}'_3} = j \circ c_3$ である。
補題 \ref{defos-lemma-inf-obs-ext} の証明では、
$o(\mathcal{F}, \mathcal{K}_3, c_3)$ を、$\mathcal{O}_X$-加群の拡大
$$
0 \to
\mathcal{K}/\mathcal{K}_3 \to
\mathcal{E}'_3/\mathcal{K}_3 \to
\mathcal{F} \to 0
$$
の類の境界として構成している。一方、
$\mathcal{F}'_1 = \mathcal{F}'_2/\mathcal{K}_3$ であることに注意する。
したがって類 $\xi_{\mathcal{F}'_1}$ は拡大
$$
0 \to \mathcal{K}_2/\mathcal{K}_3 \to \mathcal{F}'_2/\mathcal{K}_3
\to \mathcal{F} \to 0
$$
の類である。ここでは、正準な分裂
$\pi : (X'_1, \mathcal{O}_{X'_1}) \to (X, \mathcal{O}_X)$ の
$\pi^\sharp$ を用いて、この列を $\mathcal{O}_X$-加群の列とみなしている。
最後に、次の可換図式が存在することから主張が従う。
$$
\xymatrix{
0 \ar[r] &
\mathcal{K}_2/\mathcal{K}_3 \ar[r] \ar[d] &
\mathcal{F}'_2/\mathcal{K}_3 \ar[r] \ar[d] &
\mathcal{F} \ar[r] \ar[d] & 0 \\
0 \ar[r] &
\mathcal{K}/\mathcal{K}_3 \ar[r] &
\mathcal{E}'_3/\mathcal{K}_3 \ar[r] &
\mathcal{F} \ar[r] & 0
}
$$
この図式は、上で与えた $\mathcal{O}_X$-加群構造に関して
$\mathcal{O}_X$-線形である。
\end{remark}








\section{環付き空間上の加群の無限小変形}
\label{defos-section-deformation-modules}

\noindent
$i : (X, \mathcal{O}_X) \to (X', \mathcal{O}_{X'})$ を環付き空間の
一次厚化とする。節 \ref{defos-section-thickenings-spaces} で導入した記法を
断りなく用いる。$\mathcal{F}'$ を $\mathcal{O}_{X'}$-加群とし、
$\mathcal{F} = i^*\mathcal{F}'$ とおく。この状況では、
$\mathcal{O}_{X'}$-加群の短完全列
$$
0 \to \mathcal{I}\mathcal{F}' \to \mathcal{F}' \to \mathcal{F} \to 0
$$
を得る。$\mathcal{I}^2 = 0$ なので、$\mathcal{I}\mathcal{F}'$ 上の
$\mathcal{O}_{X'}$-加群構造は一意な $\mathcal{O}_X$-加群構造から生じる。
したがって上の列は (\ref{defos-equation-extension}) の形の拡大である。
特別な場合として、$\mathcal{F}' = \mathcal{O}_{X'}$ ならば
$i^*\mathcal{O}_{X'} = \mathcal{O}_X$ かつ
$\mathcal{I}\mathcal{O}_{X'} = \mathcal{I}$ であり、構造層の列
$$
0 \to \mathcal{I} \to \mathcal{O}_{X'} \to \mathcal{O}_X \to 0
$$
を再び得る。

\begin{lemma}
\label{defos-lemma-inf-map-special}
$i : (X, \mathcal{O}_X) \to (X', \mathcal{O}_{X'})$ を環付き空間の
一次厚化とする。$\mathcal{F}'$、$\mathcal{G}'$ を
$\mathcal{O}_{X'}$-加群とし、$\mathcal{F} = i^*\mathcal{F}'$、
$\mathcal{G} = i^*\mathcal{G}'$ とおく。
$\varphi : \mathcal{F} \to \mathcal{G}$ を $\mathcal{O}_X$-線形写像とする。
$\varphi$ を $\mathcal{O}_{X'}$-線形写像
$\varphi' : \mathcal{F}' \to \mathcal{G}'$ に持ち上げる方法全体の集合は、
空でなければ
$\Hom_{\mathcal{O}_X}(\mathcal{F}, \mathcal{I}\mathcal{G}')$
の下で主等質空間をなす。
\end{lemma}

\begin{proof}
これは補題 \ref{defos-lemma-inf-map} の特別な場合であるが、直接証明も与える。
加群の短完全列
$$
0 \to \mathcal{I} \to \mathcal{O}_{X'} \to \mathcal{O}_X \to 0
\quad\text{and}\quad
0 \to \mathcal{I}\mathcal{G}' \to \mathcal{G}' \to \mathcal{G} \to 0
$$
があり、$\mathcal{F}'$ についても同様である。$\mathcal{I}$ は二乗零なので、
$\mathcal{I}$ および $\mathcal{I}\mathcal{G}'$ 上の
$\mathcal{O}_{X'}$-加群構造は一意な $\mathcal{O}_X$-加群構造から生じる。
したがって、
$$
\Hom_{\mathcal{O}_{X'}}(\mathcal{F}', \mathcal{I}\mathcal{G}') =
\Hom_{\mathcal{O}_X}(\mathcal{F}, \mathcal{I}\mathcal{G}')
\quad\text{and}\quad
\Hom_{\mathcal{O}_{X'}}(\mathcal{F}', \mathcal{G}) =
\Hom_{\mathcal{O}_X}(\mathcal{F}, \mathcal{G})
$$
補題は完全列
$$
0 \to \Hom_{\mathcal{O}_{X'}}(\mathcal{F}', \mathcal{I}\mathcal{G}') \to
\Hom_{\mathcal{O}_{X'}}(\mathcal{F}', \mathcal{G}') \to
\Hom_{\mathcal{O}_{X'}}(\mathcal{F}', \mathcal{G})
$$
から従う。『ホモロジー』補題 \ref{homology-lemma-check-exactness} を見よ。
\end{proof}

\begin{lemma}
\label{defos-lemma-deform-module}
$(f, f')$ を状況 \ref{defos-situation-morphism-thickenings} のような
環付き空間の一次厚化の射とする。$\mathcal{F}'$ を
$\mathcal{O}_{X'}$-加群とし、$\mathcal{F} = i^*\mathcal{F}'$ とおく。
$\mathcal{F}$ は $S$ 上平坦であり、$(f, f')$ は厳密な厚化の射であると
仮定する（定義 \ref{defos-definition-strict-morphism-thickenings}）。
このとき次は同値である。
\begin{enumerate}
\item $\mathcal{F}'$ は $S'$ 上平坦である。
\item 正準写像
$f^*\mathcal{J} \otimes_{\mathcal{O}_X} \mathcal{F} \to
\mathcal{I}\mathcal{F}'$
は同型である。
\end{enumerate}
さらに、この場合、写像
$$
f^*\mathcal{J} \otimes_{\mathcal{O}_X} \mathcal{F} \to
\mathcal{I} \otimes_{\mathcal{O}_X} \mathcal{F} \to
\mathcal{I}\mathcal{F}'
$$
はいずれも同型である。
\end{lemma}

\begin{proof}
$(f, f')$ は厳密な厚化の射なので、写像
$f^*\mathcal{J} \to \mathcal{I}$ は全射である。したがって最後の主張は
(2) から従う。

\medskip\noindent
(1) と (2) の同値を証明する。これらの条件は茎で確認できる。
$x \in X \subset X'$ を点とし、その像を
$s = f(x) \in S \subset S'$ とする。
$A' = \mathcal{O}_{S', s}$、$B' = \mathcal{O}_{X', x}$、
$A = \mathcal{O}_{S, s}$、$B = \mathcal{O}_{X, x}$ とおく。
ある二乗零イデアルについて $A = A'/J$、$B = B'/I$ である。
$(f, f')$ は厳密な厚化の射なので $I = JB'$ である。
$M' = \mathcal{F}'_x$、$M = \mathcal{F}_x$ とおく。
このとき $M'$ は $B'$-加群、$M$ は $B$-加群である。
$\mathcal{F} = i^*\mathcal{F}'$ なので、全射 $M' \to M$ の核は
$IM' = JM'$ である。したがって短完全列
$$
0 \to JM' \to M' \to M \to 0
$$
を得る。『層』補題 \ref{sheaves-lemma-stalk-pullback-modules} および
『加群』補題 \ref{modules-lemma-stalk-tensor-product} を用いて
引き戻しとテンソル積の茎を同一視すると、補題の正準写像の $x$ における
茎は写像
$$
(J \otimes_A B) \otimes_B M = J \otimes_A M = J \otimes_{A'} M'
\longrightarrow JM'
$$
である。$\mathcal{F}$ が $S$ 上平坦であるという仮定は、$M$ が
平坦 $A$-加群であることを意味する。

\medskip\noindent
(1) を仮定する。平坦性と『代数』補題
\ref{algebra-lemma-characterize-flat} により
$\text{Tor}_1^{A'}(M', A) = 0$ である。『代数』注意
\ref{algebra-remark-Tor-ring-mod-ideal} により、これは
$J \otimes_{A'} M' \to M'$ が単射であることを意味する。
したがって $J \otimes_A M \to JM'$ は同型である。

\medskip\noindent
(2) を仮定する。このとき $J \otimes_{A'} M' \to M'$ は単射である。
したがって『代数』注意 \ref{algebra-remark-Tor-ring-mod-ideal} により
$\text{Tor}_1^{A'}(M', A) = 0$ である。ゆえに『代数』補題
\ref{algebra-lemma-what-does-it-mean} により $M'$ は $A'$ 上平坦である。
\end{proof}

\begin{lemma}
\label{defos-lemma-inf-map-rel}
$(f, f')$ を状況 \ref{defos-situation-morphism-thickenings} のような
一次厚化の射とする。$\mathcal{F}'$、$\mathcal{G}'$ を
$\mathcal{O}_{X'}$-加群とし、$\mathcal{F} = i^*\mathcal{F}'$、
$\mathcal{G} = i^*\mathcal{G}'$ とおく。
$\varphi : \mathcal{F} \to \mathcal{G}$ を $\mathcal{O}_X$-線形写像とする。
$\mathcal{G}'$ は $S'$ 上平坦であり、$(f, f')$ は厳密な厚化の射であると
仮定する。$\varphi$ を $\mathcal{O}_{X'}$-線形写像
$\varphi' : \mathcal{F}' \to \mathcal{G}'$ に持ち上げる方法全体の集合は、
空でなければ次の群の下で主等質空間をなす。
$$
\Hom_{\mathcal{O}_X}(\mathcal{F},
\mathcal{G} \otimes_{\mathcal{O}_X} f^*\mathcal{J})
$$
\end{lemma}

\begin{proof}
補題 \ref{defos-lemma-inf-map-special} と補題 \ref{defos-lemma-deform-module} を
組み合わせればよい。
\end{proof}

\begin{lemma}
\label{defos-lemma-inf-obs-map-special}
$i : (X, \mathcal{O}_X) \to (X', \mathcal{O}_{X'})$ を環付き空間の
一次厚化とする。$\mathcal{F}'$、$\mathcal{G}'$ を
$\mathcal{O}_{X'}$-加群とし、$\mathcal{F} = i^*\mathcal{F}'$、
$\mathcal{G} = i^*\mathcal{G}'$ とおく。
$\varphi : \mathcal{F} \to \mathcal{G}$ を $\mathcal{O}_X$-線形写像とする。
元
$$
o(\varphi) \in
\Ext^1_{\mathcal{O}_X}(Li^*\mathcal{F}',
\mathcal{I}\mathcal{G}')
$$
が存在し、その消滅は、$\varphi$ を $\mathcal{O}_{X'}$-線形写像
$\varphi' : \mathcal{F}' \to \mathcal{G}'$ に持ち上げられるための
必要十分条件である。
\end{lemma}

\begin{proof}
補題 \ref{defos-lemma-inf-map-special} の証明から、写像
$$
\Hom_{\mathcal{O}_X}(\mathcal{F}, \mathcal{G}) =
\Hom_{\mathcal{O}_{X'}}(\mathcal{F}', \mathcal{G}) \longrightarrow
\Ext^1_{\mathcal{O}_{X'}}(\mathcal{F}', \mathcal{I}\mathcal{G}')
$$
による $\varphi$ の境界が消滅することは、持ち上げが存在するための
必要十分条件であることが分かる。導来圏における $i_* = Ri_*$ と
$Li^*$ の随伴性（『コホモロジー』補題
\ref{cohomology-lemma-adjoint}）により、
$$
\Ext^1_{\mathcal{O}_{X'}}(\mathcal{F}', \mathcal{I}\mathcal{G}') =
\Ext^1_{\mathcal{O}_X}(Li^*\mathcal{F}', \mathcal{I}\mathcal{G}')
$$
であるから、結論を得る。
\end{proof}

\begin{lemma}
\label{defos-lemma-inf-obs-map-rel}
$(f, f')$ を状況 \ref{defos-situation-morphism-thickenings} のような
一次厚化の射とする。$\mathcal{F}'$、$\mathcal{G}'$ を
$\mathcal{O}_{X'}$-加群とし、$\mathcal{F} = i^*\mathcal{F}'$、
$\mathcal{G} = i^*\mathcal{G}'$ とおく。
$\varphi : \mathcal{F} \to \mathcal{G}$ を $\mathcal{O}_X$-線形写像とする。
$\mathcal{F}'$ と $\mathcal{G}'$ は $S'$ 上平坦であり、$(f, f')$ は
厳密な厚化の射であると仮定する。このとき元
$$
o(\varphi) \in  \Ext^1_{\mathcal{O}_X}(\mathcal{F},
\mathcal{G} \otimes_{\mathcal{O}_X} f^*\mathcal{J})
$$
が存在し、その消滅は、$\varphi$ を $\mathcal{O}_{X'}$-線形写像
$\varphi' : \mathcal{F}' \to \mathcal{G}'$ に持ち上げられるための
必要十分条件である。
\end{lemma}

\begin{proof}[第一の証明]
補題 \ref{defos-lemma-inf-obs-map-special} から従う。実際、補題の仮定の下で
$$
\Ext^1_{\mathcal{O}_X}(Li^*\mathcal{F}',
\mathcal{I}\mathcal{G}') =
\Ext^1_{\mathcal{O}_X}(\mathcal{F},
\mathcal{G} \otimes_{\mathcal{O}_X} f^*\mathcal{J})
$$
であると主張する。まず、補題 \ref{defos-lemma-deform-module} により
$\mathcal{I}\mathcal{G}' =
\mathcal{G} \otimes_{\mathcal{O}_X} f^*\mathcal{J}$
である。一方、
$$
H^{-1}(Li^*\mathcal{F}') =
\text{Tor}_1^{\mathcal{O}_{X'}}(\mathcal{F}', \mathcal{O}_X)
$$
である（局所計算は省略する）。短完全列
$$
0 \to \mathcal{I} \to \mathcal{O}_{X'} \to \mathcal{O}_X \to 0
$$
を用いると、この $\text{Tor}_1$ は写像
$\mathcal{I} \otimes_{\mathcal{O}_X} \mathcal{F} \to \mathcal{I}\mathcal{F}'$
の核で計算されることが分かる。この核は補題
\ref{defos-lemma-deform-module} の最後の主張により零である。したがって
$\tau_{\geq -1}Li^*\mathcal{F}' = \mathcal{F}$ である。一方、
$$
\Ext^1_{\mathcal{O}_X}(Li^*\mathcal{F}',
\mathcal{I}\mathcal{G}') =
\Ext^1_{\mathcal{O}_X}(\tau_{\geq -1}Li^*\mathcal{F}',
\mathcal{I}\mathcal{G}')
$$
である。これは『導来圏』補題
\ref{derived-lemma-negative-vanishing} の双対から従う。
\end{proof}

\begin{proof}[第二の証明]
補題 \ref{defos-lemma-inf-obs-map} を次のように適用できる。補題
\ref{defos-lemma-deform-module} により
$\mathcal{K} = \mathcal{I} \otimes_{\mathcal{O}_X} \mathcal{F}$ かつ
$\mathcal{L} = \mathcal{I} \otimes_{\mathcal{O}_X} \mathcal{G}$
であることに注意する。また
$c_{\mathcal{F}'} = 1 \otimes 1$ かつ $c_{\mathcal{G}'} = 1 \otimes 1$ であり、
$\psi = 1 \otimes \varphi$ とおけば補題の図式は可換である。
したがって $o(\varphi) = o(\varphi, 1 \otimes \varphi)$ とすればよい。
\end{proof}

\begin{lemma}
\label{defos-lemma-inf-ext-rel}
$(f, f')$ を状況 \ref{defos-situation-morphism-thickenings} のような
一次厚化の射とし、$\mathcal{F}$ を $\mathcal{O}_X$-加群とする。
$(f, f')$ は厳密な厚化の射であり、$\mathcal{F}$ は $S$ 上平坦であると
仮定する。$S'$ 上平坦な $\mathcal{O}_{X'}$-加群 $\mathcal{F}'$ と同型
$\alpha : i^*\mathcal{F}' \to \mathcal{F}$ からなる組
$(\mathcal{F}', \alpha)$ が存在するならば、そのような組の同型類全体の
集合は次の群の下で主等質空間をなす。
$\Ext^1_{\mathcal{O}_X}(
\mathcal{F}, \mathcal{I} \otimes_{\mathcal{O}_X} \mathcal{F})$.
\end{lemma}

\begin{proof}
そのような加群が一つ存在すると仮定すると、補題
\ref{defos-lemma-deform-module} により正準写像
$$
f^*\mathcal{J} \otimes_{\mathcal{O}_X} \mathcal{F} \to
\mathcal{I} \otimes_{\mathcal{O}_X} \mathcal{F}
$$
は同型である。補題 \ref{defos-lemma-inf-ext} を $\mathcal{K} = 
\mathcal{I} \otimes_{\mathcal{O}_X} \mathcal{F}$
および $c = 1$ として適用する。補題 \ref{defos-lemma-deform-module} により、
対応する拡大 $\mathcal{F}'$ はすべて $S'$ 上平坦である。
\end{proof}

\begin{lemma}
\label{defos-lemma-inf-obs-ext-rel}
$(f, f')$ を状況 \ref{defos-situation-morphism-thickenings} のような
一次厚化の射とし、$\mathcal{F}$ を $\mathcal{O}_X$-加群とする。
$(f, f')$ は厳密な厚化の射であり、$\mathcal{F}$ は $S$ 上平坦であると
仮定する。$i^*\mathcal{F}' \cong \mathcal{F}$ を満たす $S'$ 上平坦な
$\mathcal{O}_{X'}$-加群 $\mathcal{F}'$ が存在するための必要十分条件は、
次の二条件である。
\begin{enumerate}
\item 正準写像 $
f^*\mathcal{J} \otimes_{\mathcal{O}_X} \mathcal{F} \to
\mathcal{I} \otimes_{\mathcal{O}_X} \mathcal{F}$
は同型である。
\item 補題 \ref{defos-lemma-inf-obs-ext} の類
$o(\mathcal{F}, \mathcal{I} \otimes_{\mathcal{O}_X} \mathcal{F}, 1)
\in \Ext^2_{\mathcal{O}_X}(
\mathcal{F}, \mathcal{I} \otimes_{\mathcal{O}_X} \mathcal{F})$
は零である。
\end{enumerate}
\end{lemma}

\begin{proof}
補題 \ref{defos-lemma-deform-module} による $S'$ 上平坦な
$\mathcal{O}_{X'}$-加群の特徴づけと、補題
\ref{defos-lemma-inf-obs-ext} から直ちに従う。
\end{proof}






\section{環付き空間の平坦な厚化上の平坦加群への応用}
\label{defos-section-flat}

\noindent
次の環付き空間の可換図式を考える。
$$
\xymatrix{
(X, \mathcal{O}_X) \ar[r]_i \ar[d]_f & (X', \mathcal{O}_{X'}) \ar[d]^{f'} \\
(S, \mathcal{O}_S) \ar[r]^t & (S', \mathcal{O}_{S'})
}
$$
横の矢印は状況 \ref{defos-situation-morphism-thickenings} のような一次厚化である。
$\mathcal{I} = \Ker(i^\sharp) \subset \mathcal{O}_{X'}$ および
$\mathcal{J} = \Ker(t^\sharp) \subset \mathcal{O}_{S'}$ とおく。
$\mathcal{F}$ を $\mathcal{O}_X$-加群とし、次を仮定する。
\begin{enumerate}
\item $(f, f')$ は厳密な厚化の射である。
\item $f'$ は平坦である。
\item $\mathcal{F}$ は $S$ 上平坦である。
\end{enumerate}
(1) $+$ (2) から $\mathcal{I} = f^*\mathcal{J}$ が従うことに注意する
（補題 \ref{defos-lemma-deform-module} を $\mathcal{O}_{X'}$ に適用せよ）。
これらの仮定の下では、前節の理論は特に簡明になる。すでに得られた結果を
次の補題にまとめる。

\begin{lemma}
\label{defos-lemma-flat}
上の状況において、次が成り立つ。
\begin{enumerate}
\item $i^*\mathcal{F}' \cong \mathcal{F}$ を満たす $S'$ 上平坦な
$\mathcal{O}_{X'}$-加群 $\mathcal{F}'$ が存在するための必要十分条件は、
補題 \ref{defos-lemma-inf-obs-ext} の類
$o(\mathcal{F}, f^*\mathcal{J} \otimes_{\mathcal{O}_X} \mathcal{F}, 1)
\in \Ext^2_{\mathcal{O}_X}(
\mathcal{F}, f^*\mathcal{J} \otimes_{\mathcal{O}_X} \mathcal{F})$
が零であることである。
\item そのような加群が存在するならば、持ち上げの同型類全体の集合は
$\Ext^1_{\mathcal{O}_X}(
\mathcal{F}, f^*\mathcal{J} \otimes_{\mathcal{O}_X} \mathcal{F})$.
の下で主等質空間をなす。
\item 持ち上げ $\mathcal{F}'$ が与えられたとき、引き戻すと
$\text{id}_\mathcal{F}$ になる $\mathcal{F}'$ の自己同型全体の集合は、
正準に $\Ext^0_{\mathcal{O}_X}(
\mathcal{F}, f^*\mathcal{J} \otimes_{\mathcal{O}_X} \mathcal{F})$
と同型である。
\end{enumerate}
\end{lemma}

\begin{proof}
(1) は、上で $\mathcal{I} = f^*\mathcal{J}$ を見たので、補題
\ref{defos-lemma-inf-obs-ext-rel} から従う。(2) は補題
\ref{defos-lemma-inf-ext-rel} から従い、(3) は補題
\ref{defos-lemma-inf-map-rel} から従う。
\end{proof}

\begin{situation}
\label{defos-situation-ses-flat-thickenings}
$f : (X, \mathcal{O}_X) \to (S, \mathcal{O}_S)$ を環付き空間の射とする。
次の可換図式を考える。
$$
\xymatrix{
(X'_1, \mathcal{O}'_1) \ar[r]_h \ar[d]_{f'_1} &
(X'_2, \mathcal{O}'_2) \ar[r] \ar[d]_{f'_2} &
(X'_3, \mathcal{O}'_3) \ar[d]_{f'_3} \\
(S'_1, \mathcal{O}_{S'_1}) \ar[r] &
(S'_2, \mathcal{O}_{S'_2}) \ar[r] &
(S'_3, \mathcal{O}_{S'_3})
}
$$
ここで、(a) 上段は $X$ の一次厚化の短完全列、(b) 下段は $S$ の
一次厚化の短完全列、(c) 各 $f'_i$ の制限は $f$、(d) 各組
$(f, f_i')$ は厳密な厚化の射、(e) 各 $f'_i$ は平坦であるとする。
最後に、$\mathcal{F}'_2$ を $S'_2$ 上平坦な $\mathcal{O}'_2$-加群とし、
$\mathcal{F} = \mathcal{F}'_2|_X$ とおく。
$\pi : X'_1 \to X$ を正準な分裂とする（注意
\ref{defos-remark-short-exact-sequence-thickenings}）。
\end{situation}

\begin{lemma}
\label{defos-lemma-verify-iv}
状況 \ref{defos-situation-ses-flat-thickenings} において、加群
$\pi^*\mathcal{F}$ と $h^*\mathcal{F}'_2$ は $S'_1$ 上平坦な
$\mathcal{O}'_1$-加群であり、$X$ 上での制限は $\mathcal{F}$ である。
その差（補題 \ref{defos-lemma-flat}）は
$\Ext^1_{\mathcal{O}_X}(
\mathcal{F}, f^*\mathcal{J}_1 \otimes_{\mathcal{O}_X} \mathcal{F})$
の元 $\theta$ であり、その
$\Ext^2_{\mathcal{O}_X}(
\mathcal{F}, f^*\mathcal{J}_3 \otimes_{\mathcal{O}_X} \mathcal{F})$
における境界は、$\mathcal{F}$ を $S'_3$ 上平坦な
$\mathcal{O}'_3$-加群に持ち上げることへの障害（補題
\ref{defos-lemma-flat}）に等しい。
\end{lemma}

\begin{proof}
$\pi^*\mathcal{F}$ と $h^*\mathcal{F}'_2$ はいずれも $X$ 上で
$\mathcal{F}$ に制限され、$\pi^*\mathcal{F} \to \mathcal{F}$ と
$h^*\mathcal{F}'_2 \to \mathcal{F}$ の核はいずれも
$f^*\mathcal{J}_1 \otimes_{\mathcal{O}_X} \mathcal{F}$ で与えられる。
したがって補題 \ref{defos-lemma-deform-module} により平坦である。
加群の列
$$
0 \to f^*\mathcal{J}_3 \otimes_{\mathcal{O}_X} \mathcal{F} \to
f^*\mathcal{J}_2 \otimes_{\mathcal{O}_X} \mathcal{F} \to
f^*\mathcal{J}_1 \otimes_{\mathcal{O}_X} \mathcal{F} \to 0
$$
は、状況 \ref{defos-situation-ses-flat-thickenings} の仮定と
$\mathcal{F}$ が $S$ 上平坦であることにより短完全である。したがって
境界をとることには意味がある。障害類についての主張は、注意
\ref{defos-remark-complex-thickenings-and-ses-modules} の結果をこの特別な
状況に移したものにほかならない。
\end{proof}






\section{環付き空間の変形と素朴な余接複体}
\label{defos-section-deformations-ringed-spaces}

\noindent
本節では、素朴な余接複体を用いて変形理論を少し展開する。まず、
環付き空間の一次厚化
$t : (S, \mathcal{O}_S) \to (S', \mathcal{O}_{S'})$ から始める。
$\mathcal{J} = \Ker(t^\sharp)$ と記し、$S$ と $S'$ の台となる位相空間を
同一視する。さらに、環付き空間の射
$f : (X, \mathcal{O}_X) \to (S, \mathcal{O}_S)$、
$\mathcal{O}_X$-加群 $\mathcal{G}$、および加群層の $f$-写像
$c : \mathcal{J} \to \mathcal{G}$ が与えられていると仮定する
（『層』定義 \ref{sheaves-definition-f-map} および節
\ref{sheaves-section-ringed-spaces-functoriality-modules}）。本節では、
次の図式の疑問符に入るものを見いだせるかを問う。
\begin{equation}
\label{defos-equation-to-solve-ringed-spaces}
\vcenter{
\xymatrix{
0 \ar[r] & \mathcal{G} \ar[r] & {?} \ar[r] & \mathcal{O}_X \ar[r] & 0 \\
0 \ar[r] & \mathcal{J} \ar[u]^c \ar[r] & \mathcal{O}_{S'} \ar[u] \ar[r] &
\mathcal{O}_S \ar[u] \ar[r] & 0
}
}
\end{equation}
ここで垂直の矢印は $f$-写像である。また、解が存在する場合にそれが
どの程度一意であるかも問う。より正確には、一次厚化
$i : (X, \mathcal{O}_X) \to (X', \mathcal{O}_{X'})$
と、(\ref{defos-equation-morphism-thickenings}) のような厚化の射 $(f, f')$
であって、$\Ker(i^\sharp)$ が $\mathcal{G}$ と同一視され、
$(f')^\sharp$ が与えられた写像 $c$ を誘導するものを求める。
$X'$ を (\ref{defos-equation-to-solve-ringed-spaces}) の {\it 解} と呼ぶ。

\begin{lemma}
\label{defos-lemma-huge-diagram-ringed-spaces}
環付き空間の射の可換図式
\begin{equation}
\label{defos-equation-huge-1}
\vcenter{
\xymatrix{
& (X_2, \mathcal{O}_{X_2}) \ar[r]_{i_2} \ar[d]_{f_2} \ar[ddl]_g &
(X'_2, \mathcal{O}_{X'_2}) \ar[d]^{f'_2} \\
& (S_2, \mathcal{O}_{S_2}) \ar[r]^{t_2} \ar[ddl]|\hole &
(S'_2, \mathcal{O}_{S'_2}) \ar[ddl] \\
(X_1, \mathcal{O}_{X_1}) \ar[r]_{i_1} \ar[d]_{f_1} &
(X'_1, \mathcal{O}_{X'_1}) \ar[d]^{f'_1} \\
(S_1, \mathcal{O}_{S_1}) \ar[r]^{t_1} &
(S'_1, \mathcal{O}_{S'_1})
}
}
\end{equation}
が与えられ、横の矢印は一次厚化であると仮定する。
$\mathcal{G}_j = \Ker(i_j^\sharp)$ とおき、加群の $g$-写像
$\nu : \mathcal{G}_1 \to \mathcal{G}_2$ が与えられ、それが可換図式
\begin{equation}
\label{defos-equation-huge-2}
\vcenter{
\xymatrix{
& 0 \ar[r] & \mathcal{G}_2 \ar[r] &
\mathcal{O}_{X'_2} \ar[r] &
\mathcal{O}_{X_2} \ar[r] & 0 \\
& 0 \ar[r]|\hole &
\mathcal{J}_2 \ar[u]_{c_2} \ar[r] &
\mathcal{O}_{S'_2} \ar[u] \ar[r]|\hole &
\mathcal{O}_{S_2} \ar[u] \ar[r] & 0 \\
0 \ar[r] & \mathcal{G}_1 \ar[ruu] \ar[r] &
\mathcal{O}_{X'_1} \ar[r] &
\mathcal{O}_{X_1} \ar[ruu] \ar[r] & 0 \\
0 \ar[r] & \mathcal{J}_1 \ar[ruu]|\hole \ar[u]^{c_1} \ar[r] &
\mathcal{O}_{S'_1} \ar[ruu]|\hole \ar[u] \ar[r] &
\mathcal{O}_{S_1} \ar[ruu]|\hole \ar[u] \ar[r] & 0
}
}
\end{equation}
を誘導すると仮定する。その手前と奥は
(\ref{defos-equation-to-solve-ringed-spaces}) の解である。
\begin{enumerate}
\item
$\Ext^1_{\mathcal{O}_{X_2}}(Lg^*\NL_{X_1/S_1}, \mathcal{G}_2)$
には正準元が存在し、その消滅は、$\nu$ と両立しながら
(\ref{defos-equation-huge-1}) に収まる環付き空間の射
$X'_2 \to X'_1$ が存在するための必要十分条件である。
\item $\nu$ と両立しながら (\ref{defos-equation-huge-1}) に収まる射
$X'_2 \to X'_1$ が存在するならば、そのような射全体の集合は
$$
\Hom_{\mathcal{O}_{X_1}}(\Omega_{X_1/S_1}, g_*\mathcal{G}_2) =
\Hom_{\mathcal{O}_{X_2}}(g^*\Omega_{X_1/S_1}, \mathcal{G}_2) =
\Ext^0_{\mathcal{O}_{X_2}}(Lg^*\NL_{X_1/S_1}, \mathcal{G}_2).
$$
の下で主等質空間をなす。
\end{enumerate}
\end{lemma}

\begin{proof}
素朴な余接複体 $\NL_{X_1/S_1}$ は『加群』定義
\ref{modules-definition-cotangent-complex-morphism-ringed-topoi} で定義される。
補題の最後の主張の等号は、$g^*$ が $g_*$ の左随伴であること、
$H^0(\NL_{X_1/S_1}) = \Omega_{X_1/S_1}$ であること
（素朴な余接複体の構成による）、および $Lg^*$ が $g^*$ の左導来関手で
あることから従う。したがって証明の残りでは、群
$\Ext^k_{\mathcal{O}_{X_2}}(Lg^*\NL_{X_1/S_1}, \mathcal{G}_2)$,
$k = 0, 1$ を用いる。まず、補題に現れるすべての環付き空間の台となる
位相空間が同じ場合に帰着できることを示す。

\medskip\noindent
そのため、$g^{-1}\NL_{X_1/S_1}$ は環の層の準同型
$g^{-1}f_1^{-1}\mathcal{O}_{S_1} \to g^{-1}\mathcal{O}_{X_1}$ の
素朴な余接複体に等しいことに注意する。『加群』補題
\ref{modules-lemma-pullback-NL} を見よ。さらに、$\NL_{X_1/S_1}$ の
次数 $0$ の項は平坦 $\mathcal{O}_{X_1}$-加群なので、正準写像
$$
Lg^*\NL_{X_1/S_1}
\longrightarrow
g^{-1}\NL_{X_1/S_1} \otimes_{g^{-1}\mathcal{O}_{X_1}} \mathcal{O}_{X_2}
$$
は次数 $0$ および $-1$ のコホモロジー層上で同型を誘導する。
したがって補題の Ext 群を
$$
\Ext^k_{g^{-1}\mathcal{O}_{X_1}}(g^{-1}\NL_{X_1/S_1}, \mathcal{G}_2) =
\Ext^k_{g^{-1}\mathcal{O}_{X_1}}(
\NL_{g^{-1}\mathcal{O}_{X_1}/g^{-1}f_1^{-1}\mathcal{O}_{S_1}}, \mathcal{G}_2)
$$
で置き換えられる。$\nu$ と両立しながら
(\ref{defos-equation-huge-1}) に収まる環付き空間の射
$X'_2 \to X'_1$ 全体の集合は、$f^\sharp$ および $\nu$ と両立する
$g^{-1}f_1^{-1}\mathcal{O}_{S'_1}$-代数準同型
$g^{-1}\mathcal{O}_{X'_1} \to \mathcal{O}_{X'_2}$ 全体の集合と全単射する。
このようにして、$X$ 上の層の図式 (\ref{defos-equation-huge-2}) があり、
それに収まる環の層の準同型
$\mathcal{O}_{X'_1} \to \mathcal{O}_{X'_2}$ を求めるものと仮定できる。

\medskip\noindent
補題の証明の残りでは、台となる位相空間はすべて同じであると仮定する。
すなわち、空間 $X$ 上の層の図式 (\ref{defos-equation-huge-2}) があり、
それに収まる環の層の準同型
$\mathcal{O}_{X'_1} \to \mathcal{O}_{X'_2}$ を求める。
Ext 群としては
$\Ext^k_{\mathcal{O}_{X_1}}(
\NL_{\mathcal{O}_{X_1}/\mathcal{O}_{S_1}}, \mathcal{G}_2)$、$k = 0, 1$ を用いる。

\medskip\noindent
第1段階。障害類の構成。集合の層
$$
\mathcal{E} = \mathcal{O}_{X'_1} \times_{\mathcal{O}_{X_2}} \mathcal{O}_{X'_2}
$$
を考える。これには全射 $\alpha : \mathcal{E} \to \mathcal{O}_{X_1}$ が
付随するので、$\NL_{\mathcal{O}_{X_1}/\mathcal{O}_{S_1}}$ の代わりに
$\NL(\alpha)$ を用いることができる。『加群』補題
\ref{modules-lemma-NL-up-to-qis} を見よ。次のようにおく。
$$
\mathcal{I}' =
\Ker(\mathcal{O}_{S'_1}[\mathcal{E}] \to \mathcal{O}_{X_1})
\quad\text{and}\quad
\mathcal{I} =
\Ker(\mathcal{O}_{S_1}[\mathcal{E}] \to \mathcal{O}_{X_1})
$$
全射 $\mathcal{I}' \to \mathcal{I}$ があり、その核は
$\mathcal{J}_1\mathcal{O}_{S'_1}[\mathcal{E}]$ である。
二つの $\mathcal{O}_{S'_1}$-代数準同型
$$
a : \mathcal{O}_{S'_1}[\mathcal{E}] \to \mathcal{O}_{X'_1}
\quad\text{and}\quad
b : \mathcal{O}_{S'_1}[\mathcal{E}] \to \mathcal{O}_{X'_2}
$$
を得る。これらは写像
$a|_{\mathcal{I}'} : \mathcal{I}' \to \mathcal{G}_1$ および
$b|_{\mathcal{I}'} : \mathcal{I}' \to \mathcal{G}_2$ を誘導する。
$a$ と $b$ はともに $(\mathcal{I}')^2$ を零化する。さらに、
(\ref{defos-equation-huge-2}) の左側の正方形は可換なので、$a$ と $b$ は
$\mathcal{G}_2$ への写像として
$\mathcal{J}_1\mathcal{O}_{S'_1}[\mathcal{E}]$ 上で一致する。したがって差
$b|_{\mathcal{I}'} - \nu \circ a|_{\mathcal{I}'}$
は、良定義な $\mathcal{O}_{X_1}$-線形写像
$$
\xi : \mathcal{I}/\mathcal{I}^2 \longrightarrow \mathcal{G}_2
$$
を誘導する。この写像は、$\mathcal{I}$ の局所切断 $f$ の類を
$\nu(a(f')) - b(f')$ に写す。ここで $f'$ は $f$ を
$\mathcal{I}'$ の局所切断に持ち上げたものである。その像（下記参照）を
$[\xi] \in \Ext^1_{\mathcal{O}_{X_1}}(\NL(\alpha), \mathcal{G}_2)$ と記す。

\medskip\noindent
第2段階。$[\xi]$ の消滅が必要であること。次のように記す。
$\Omega = \Omega_{\mathcal{O}_{S_1}[\mathcal{E}]/\mathcal{O}_{S_1}}
\otimes_{\mathcal{O}_{S_1}[\mathcal{E}]} \mathcal{O}_{X_1}$.
$\NL(\alpha) = (\mathcal{I}/\mathcal{I}^2 \to \Omega)$ は完全三角形
$$
\Omega[0] \to
\NL(\alpha) \to
\mathcal{I}/\mathcal{I}^2[1] \to
\Omega[1]
$$
に収まる。したがって $[\xi]$ が零であるための必要十分条件は、ある写像
$\Omega \to \mathcal{G}_2$ に対して $\xi$ が合成
$\mathcal{I}/\mathcal{I}^2 \to \Omega \to \mathcal{G}_2$ となることである。
(\ref{defos-equation-huge-2}) に収まる環の層の準同型
$\varphi : \mathcal{O}_{X'_1} \to \mathcal{O}_{X'_2}$
が存在すると仮定する。この場合、写像
$\mathcal{O}_{S'_1}[\mathcal{E}] \to \mathcal{G}_2$,
$f' \mapsto b(f') - \varphi(a(f'))$ を考える。計算により、これは
$\mathcal{J}_1\mathcal{O}_{S'_1}[\mathcal{E}]$ を零化し、導分
$\mathcal{O}_{S_1}[\mathcal{E}] \to \mathcal{G}_2$ を誘導する。
こうして得られる線形写像 $\Omega \to \mathcal{G}_2$ により、この場合
$[\xi] = 0$ であることが分かる。

\medskip\noindent
第3段階。$[\xi]$ の消滅が十分であること。$\theta : \Omega \to \mathcal{G}_2$
を、$\xi$ が $\theta \circ (\mathcal{I}/\mathcal{I}^2 \to \Omega)$ に
等しくなるような $\mathcal{O}_{X_1}$-線形写像とする。計算により、
$$
b + \theta \circ d : \mathcal{O}_{S'_1}[\mathcal{E}] \to \mathcal{O}_{X'_2}
$$
の $\mathcal{I}'$ への制限は
$\nu \circ a : \mathcal{I}' \to \mathcal{G}_2$ と一致する。
$\mathcal{O}_{X'_1}$ は $\mathcal{I}' \to
\mathcal{O}_{S'_1}[\mathcal{E}]$ と $\mathcal{I}' \to \mathcal{G}_1$ の
押し出しなので、写像 $b + \theta \circ d$ と $a$ は
(\ref{defos-equation-huge-2}) に収まる写像
$\mathcal{O}_{X'_1} \to \mathcal{O}_{X'_2}$ を定める。

\medskip\noindent
上の特別な場合における (2) の証明。省略する。ヒント：補題
\ref{defos-lemma-huge-diagram} の (2) の証明とまったく同じである。
\end{proof}

\begin{lemma}
\label{defos-lemma-NL-represent-ext-class}
$X$ を位相空間とし、$\mathcal{A} \to \mathcal{B}$ を環の層の準同型、
$\mathcal{G}$ を $\mathcal{B}$-加群とする。
$\xi \in \Ext^1_\mathcal{B}(\NL_{\mathcal{B}/\mathcal{A}}, \mathcal{G})$
とする。このとき集合の層の写像 $\alpha : \mathcal{E} \to \mathcal{B}$
であって、$\xi \in \Ext^1_\mathcal{B}(\NL(\alpha), \mathcal{G})$ が写像
$\mathcal{I}/\mathcal{I}^2 \to \mathcal{G}$ の類となるものが存在する
（記法については証明を見よ）。
\end{lemma}

\begin{proof}
写像 $\alpha : \mathcal{E} \to \mathcal{B}$ が与えられ、
$\mathcal{A}[\mathcal{E}] \to \mathcal{B}$ が全射で核を $\mathcal{I}$ と
すると、複体
$\NL(\alpha) = (\mathcal{I}/\mathcal{I}^2 \to 
\Omega_{\mathcal{A}[\mathcal{E}]/\mathcal{A}}
\otimes_{\mathcal{A}[\mathcal{E}]} \mathcal{B})$ は正準に
$\NL_{\mathcal{B}/\mathcal{A}}$ と同型であることを思い出そう。
『加群』補題 \ref{modules-lemma-NL-up-to-qis} を見よ。さらに、
$\Omega = \Omega_{\mathcal{A}[\mathcal{E}]/\mathcal{A}}
\otimes_{\mathcal{A}[\mathcal{E}]} \mathcal{B}$ は前層
$U \mapsto \bigoplus_{e \in \mathcal{E}(U)} \mathcal{B}(U)$ に付随する層である。
言い換えれば、$\Omega$ は集合の層 $\mathcal{E}$ 上の自由
$\mathcal{B}$-加群であり、特に正準写像 $\mathcal{E} \to \Omega$ がある。

\medskip\noindent
以上を踏まえ、ある $\mathcal{E}$ を選ぶ（例えば、素朴な余接複体の
定義のように $\mathcal{E} = \mathcal{B}$ としてよい）。$\xi$ を写像
$\mathcal{I}/\mathcal{I}^2 \to \mathcal{G}$ の類として表すことへの障害は
$\Ext^1_\mathcal{B}(\Omega, \mathcal{G})$ の元である。これが
$\mathcal{B}$-加群の拡大
$0 \to \mathcal{G} \to \mathcal{H} \to \Omega \to 0$ で表されるとする。
集合の層 $\mathcal{E}' = \mathcal{E} \times_\Omega \mathcal{H}$ を考える。
これには誘導された写像 $\alpha' : \mathcal{E}' \to \mathcal{B}$ がある。
$\mathcal{I}' = \Ker(\mathcal{A}[\mathcal{E}'] \to \mathcal{B})$ および
$\Omega' = \Omega_{\mathcal{A}[\mathcal{E}']/\mathcal{A}}
\otimes_{\mathcal{A}[\mathcal{E}']} \mathcal{B}$ とおく。
擬同型 $\NL(\alpha') \to \NL(\alpha)$ による $\xi$ の引き戻しは
$\Ext^1_\mathcal{B}(\Omega', \mathcal{G})$ で零に写る。実際、
$\Omega'$ は集合の層 $\mathcal{E}'$ 上の自由 $\mathcal{B}$-加群であり、
構成により可換図式
$$
\xymatrix{
\mathcal{E}' \ar[r] \ar[d] & \mathcal{E} \ar[d] \\
\mathcal{H} \ar[r] & \Omega
}
$$
があるので、写像 $\Omega' \to \Omega$ による拡大 $\mathcal{H}$ の
引き戻しは分裂する。これで証明が完了する。
\end{proof}

\begin{lemma}
\label{defos-lemma-choices-ringed-spaces}
(\ref{defos-equation-to-solve-ringed-spaces}) に解が存在するならば、解の
同型類全体の集合は
$\Ext^1_{\mathcal{O}_X}(\NL_{X/S}, \mathcal{G})$ の下で主等質空間をなす。
\end{lemma}

\begin{proof}
(\ref{defos-equation-to-solve-ringed-spaces}) の二つの解 $X'_1$、$X'_2$ が
与えられると、補題 \ref{defos-lemma-huge-diagram-ringed-spaces} により、
写像 $X'_1 \to X'_2$ の存在に対する障害元
$o(X'_1, X'_2) \in \Ext^1_{\mathcal{O}_X}(\NL_{X/S}, \mathcal{G})$ を得る。
明らかに、この元は同型の存在に対する障害であり、したがって同型類を
区別する。ゆえに証明を終えるには、解 $X'$ と元
$\xi \in \Ext^1_{\mathcal{O}_X}(\NL_{X/S}, \mathcal{G})$ が与えられたとき、
$o(X', X'_\xi) = \xi$ を満たす第二の解 $X'_\xi$ を見いだせることを
示せば十分である。

\medskip\noindent
類 $\xi$ に対して補題 \ref{defos-lemma-NL-represent-ext-class} のような
$\alpha : \mathcal{E} \to \mathcal{O}_X$ を選ぶ。核が $\mathcal{I}$ である
全射 $f^{-1}\mathcal{O}_S[\mathcal{E}] \to \mathcal{O}_X$ と、対応する
素朴な余接複体
$\NL(\alpha) = (\mathcal{I}/\mathcal{I}^2 \to
\Omega_{f^{-1}\mathcal{O}_S[\mathcal{E}]/f^{-1}\mathcal{O}_S}
\otimes_{f^{-1}\mathcal{O}_S[\mathcal{E}]} \mathcal{O}_X)$ を考える。
補題により、$\xi$ は射
$\delta : \mathcal{I}/\mathcal{I}^2 \to \mathcal{G}$ の類である。
$\mathcal{E}$ を $\mathcal{E} \times_{\mathcal{O}_X} \mathcal{O}_{X'}$ で
置き換えることにより、$\alpha$ が写像
$\alpha' : \mathcal{E} \to \mathcal{O}_{X'}$ を経由すると仮定してもよい。

\medskip\noindent
これらの選択は $f^{-1}\mathcal{O}_{S'}$-代数準同型
$\varphi : f^{-1}\mathcal{O}_{S'}[\mathcal{E}] \to \mathcal{O}_{X'}$ を定める。
$\mathcal{I}' = \Ker(f^{-1}\mathcal{O}_{S'}[\mathcal{E}] \to
\mathcal{O}_X)$ とおく。$\varphi$ は写像
$\varphi|_{\mathcal{I}'} : \mathcal{I}' \to \mathcal{G}$ を誘導し、
$\mathcal{O}_{X'}$ は次の図式のような押し出しであることに注意する。
$$
\xymatrix{
0 \ar[r] & \mathcal{G} \ar[r] & \mathcal{O}_{X'} \ar[r] &
\mathcal{O}_X \ar[r] & 0 \\
0 \ar[r] & \mathcal{I}' \ar[u]^{\varphi|_{\mathcal{I}'}} \ar[r] &
f^{-1}\mathcal{O}_{S'}[\mathcal{E}] \ar[u] \ar[r] &
\mathcal{O}_X \ar[u]_{=} \ar[r] & 0
}
$$
$\psi : \mathcal{I}' \to \mathcal{G}$ を、写像
$\varphi|_{\mathcal{I}'}$ と合成
$$
\mathcal{I}' \to \mathcal{I}'/(\mathcal{I}')^2 \to
\mathcal{I}/\mathcal{I}^2 \xrightarrow{\delta} \mathcal{G}.
$$
との和とする。このとき $\psi$ に沿う押し出しは、上のような図式に
収まる別の環拡大 $\mathcal{O}_{X'_\xi}$ である。計算（省略）により、
所望のとおり $o(X', X'_\xi) = \xi$ であることが分かる。
\end{proof}

\begin{lemma}
\label{defos-lemma-extensions-of-relative-ringed-spaces}
$f : (X, \mathcal{O}_X) \to (S, \mathcal{O}_S)$ を環付き空間の射とし、
$\mathcal{G}$ を $\mathcal{O}_X$-加群とする。
$f^{-1}\mathcal{O}_S$-代数の拡大
$$
0 \to \mathcal{G} \to \mathcal{O}_{X'} \to \mathcal{O}_X \to 0
$$
で $\mathcal{G}$ が二乗零イデアルであるもの\footnote{言い換えれば、
$\mathcal{O}_X$-加群の同型 $\mathcal{G} \to \Ker(i^\sharp)$ を備えた、
$S$ 上の一次厚化 $i : X \to X'$ の同型類全体の集合である。} の
同型類全体の集合は、正準に
$\Ext^1_{\mathcal{O}_X}(\NL_{X/S}, \mathcal{G})$ と全単射する。
\end{lemma}

\begin{proof}
これを証明するため、(\ref{defos-equation-to-solve-ringed-spaces}) が図式
$$
\xymatrix{
0 \ar[r] &
\mathcal{G} \ar[r] &
{?} \ar[r] &
\mathcal{O}_X \ar[r] & 0 \\
0 \ar[r] &
0 \ar[u] \ar[r] &
\mathcal{O}_S \ar[u] \ar[r]^{\text{id}} &
\mathcal{O}_S \ar[u] \ar[r] & 0
}
$$
で与えられる場合に先の結果を適用する。補題
\ref{defos-lemma-choices-ringed-spaces} と、解
$\mathcal{G} \oplus \mathcal{O}_X$ が存在することから結論が従う。
（全単射の直接の構成については次の注意を見よ。）
\end{proof}

\begin{remark}
\label{defos-remark-extensions-of-relative-ringed-spaces}
補題 \ref{defos-lemma-extensions-of-relative-ringed-spaces} と同様の
$f : (X, \mathcal{O}_X) \to (S, \mathcal{O}_S)$ と $\mathcal{G}$ をとる。
補題のような拡大
$0 \to \mathcal{G} \to \mathcal{O}_{X'} \to \mathcal{O}_X \to 0$
を考える。集合の層 $\mathcal{E}$ と可換図式
$$
\xymatrix{
\mathcal{E} \ar[d]_{\alpha'} \ar[rd]^\alpha \\
\mathcal{O}_{X'} \ar[r] & \mathcal{O}_X
}
$$
を、$f^{-1}\mathcal{O}_S[\mathcal{E}] \to \mathcal{O}_X$ が核
$\mathcal{J}$ をもつ全射となるように選べる。
（例えば、$\mathcal{O}_{X'}$ に全射する任意の集合の層をとればよい。）
このとき
$$
\NL_{X/S} \cong \NL(\alpha) = 
\left(
\mathcal{J}/\mathcal{J}^2
\longrightarrow
\Omega_{f^{-1}\mathcal{O}_S[\mathcal{E}]/f^{-1}\mathcal{O}_S}
\otimes_{f^{-1}\mathcal{O}_S[\mathcal{E}]} \mathcal{O}_X\right)
$$
である。『加群』節 \ref{modules-section-netherlander}、特に補題
\ref{modules-lemma-NL-up-to-qis} を見よ。もちろん $\alpha'$ は写像
$f^{-1}\mathcal{O}_S[\mathcal{E}] \to \mathcal{O}_{X'}$ を定め、これはさらに写像
$$
\mathcal{J}/\mathcal{J}^2 \longrightarrow \mathcal{G}
$$
を定める。この写像はさらに、補題の全単射により
$\mathcal{O}_{X'}$ に対応する
$\Ext^1_{\mathcal{O}_X}(\NL(\alpha), \mathcal{G}) =
\Ext^1_{\mathcal{O}_X}(\NL_{X/S}, \mathcal{G})$ の元を定める。
\end{remark}

\begin{lemma}
\label{defos-lemma-extensions-of-relative-ringed-spaces-functorial}
$f : (X, \mathcal{O}_X) \to (S, \mathcal{O}_S)$ と
$g : (Y, \mathcal{O}_Y) \to (X, \mathcal{O}_X)$ を環付き空間の射とする。
$\mathcal{F}$ を $\mathcal{O}_X$-加群、$\mathcal{G}$ を
$\mathcal{O}_Y$-加群とし、$c : \mathcal{F} \to \mathcal{G}$ を
$g$-写像とする。最後に次を考える。
\begin{enumerate}
\item[(a)] $0 \to \mathcal{F} \to \mathcal{O}_{X'} \to \mathcal{O}_X \to 0$
$\xi \in \Ext^1_{\mathcal{O}_X}(\NL_{X/S}, \mathcal{F})$ に対応する
$f^{-1}\mathcal{O}_S$-代数の拡大。
\item[(b)] $0 \to \mathcal{G} \to \mathcal{O}_{Y'} \to \mathcal{O}_Y \to 0$
$\zeta \in \Ext^1_{\mathcal{O}_Y}(\NL_{Y/S}, \mathcal{G})$ に対応する
$g^{-1}f^{-1}\mathcal{O}_S$-代数の拡大。
\end{enumerate}
補題 \ref{defos-lemma-extensions-of-relative-ringed-spaces} を見よ。
$g$ および $c$ と両立する $S$-射 $g' : Y' \to X'$ が存在するための
必要十分条件は、$\xi$ と $\zeta$ が
$\Ext^1_{\mathcal{O}_Y}(Lg^*\NL_{X/S}, \mathcal{G})$.
の同じ元に写ることである。
\end{lemma}

\begin{proof}
この主張には意味がある。実際、写像
$$
\Ext^1_{\mathcal{O}_X}(\NL_{X/S}, \mathcal{F}) \to
\Ext^1_{\mathcal{O}_Y}(Lg^*\NL_{X/S}, Lg^*\mathcal{F}) \to
\Ext^1_{\mathcal{O}_Y}(Lg^*\NL_{X/S}, \mathcal{G})
$$
がある。ここでは写像
$Lg^*\mathcal{F} \to g^*\mathcal{F} \xrightarrow{c} \mathcal{G}$ を用いる。
また、写像
$$
\Ext^1_{\mathcal{O}_Y}(\NL_{Y/S}, \mathcal{G}) \to
\Ext^1_{\mathcal{O}_Y}(Lg^*\NL_{X/S}, \mathcal{G})
$$
があり、ここでは $Lg^*\NL_{X/S} \to \NL_{Y/S}$ を用いる。
図式
$$
\xymatrix{
& 0 \ar[r] &
\mathcal{G} \ar[r] &
\mathcal{O}_{Y'} \ar[r] &
\mathcal{O}_Y \ar[r] & 0 \\
& 0 \ar[r]|\hole & 0 \ar[u] \ar[r] &
\mathcal{O}_S \ar[u] \ar[r]|\hole &
\mathcal{O}_S \ar[u] \ar[r] & 0 \\
0 \ar[r] &
\mathcal{F} \ar[ruu] \ar[r] &
\mathcal{O}_{X'} \ar[r] &
\mathcal{O}_X \ar[ruu] \ar[r] & 0 \\
0 \ar[r] & 0 \ar[ruu]|\hole \ar[u] \ar[r] &
\mathcal{O}_S \ar[ruu]|\hole \ar[u] \ar[r] &
\mathcal{O}_S \ar[ruu]|\hole \ar[u] \ar[r] & 0
}
$$
に補題 \ref{defos-lemma-huge-diagram-ringed-spaces} を適用し、さらに補題
\ref{defos-lemma-extensions-of-relative-ringed-spaces} と補題
\ref{defos-lemma-huge-diagram-ringed-spaces} の証明における構成の両立性を
用いれば、本補題の主張が従う。この両立性の主張と証明は省略する。
（直接の議論については次の注意を見よ。）
\end{proof}

\begin{remark}
\label{defos-remark-extensions-of-relative-ringed-spaces-functorial}
$f : (X, \mathcal{O}_X) \to (S, \mathcal{O}_S)$、
$g : (Y, \mathcal{O}_Y) \to (X, \mathcal{O}_X)$、$\mathcal{F}$、
$\mathcal{G}$、$c : \mathcal{F} \to \mathcal{G}$、
$0 \to \mathcal{F} \to \mathcal{O}_{X'} \to \mathcal{O}_X \to 0$、
$\xi \in \Ext^1_{\mathcal{O}_X}(\NL_{X/S}, \mathcal{F})$、
$0 \to \mathcal{G} \to \mathcal{O}_{Y'} \to \mathcal{O}_Y \to 0$、および
$\zeta \in \Ext^1_{\mathcal{O}_Y}(\NL_{Y/S}, \mathcal{G})$ は、補題
\ref{defos-lemma-extensions-of-relative-ringed-spaces-functorial} と同様とする。
$c : g^{-1}\mathcal{F} \to \mathcal{G}$ に沿う押し出しを用いて、拡大
$$
\xymatrix{
0 \ar[r] &
\mathcal{G} \ar[r] &
\mathcal{O}'_1 \ar[r] &
g^{-1}\mathcal{O}_X \ar[r] & 0 \\
0 \ar[r] &
g^{-1}\mathcal{F} \ar[u]^c \ar[r] &
g^{-1}\mathcal{O}_{X'} \ar[u] \ar[r] &
g^{-1}\mathcal{O}_X \ar@{=}[u] \ar[r] & 0
}
$$
を構成できる。また
$g^\sharp : g^{-1}\mathcal{O}_X \to \mathcal{O}_Y$ に沿う引き戻しを用いて、
拡大
$$
\xymatrix{
0 \ar[r] &
\mathcal{G} \ar[r] &
\mathcal{O}_{Y'} \ar[r] &
\mathcal{O}_Y \ar[r] & 0 \\
0 \ar[r] &
\mathcal{G} \ar@{=}[u] \ar[r] &
\mathcal{O}'_2 \ar[u] \ar[r] &
g^{-1}\mathcal{O}_X \ar[u] \ar[r] & 0
}
$$
を構成できる。図式追跡により、$g$ および $c$ と両立する $S$-射
$Y' \to X'$ が存在するための必要十分条件は、$\mathcal{O}'_1$ と
$\mathcal{O}'_2$ が、$g^{-1}\mathcal{O}_X$ の $\mathcal{G}$ による
$g^{-1}f^{-1}\mathcal{O}_S$-代数拡大として同型であることだと分かる。
補題 \ref{defos-lemma-extensions-of-relative-ringed-spaces} により、これらの拡大は
次の左辺によって分類される。
$$
\Ext^1_{g^{-1}\mathcal{O}_X}(
\NL_{g^{-1}\mathcal{O}_X/g^{-1}f^{-1}\mathcal{O}_S}, \mathcal{G}) =
\Ext^1_{\mathcal{O}_Y}(Lg^*\NL_{X/S}, \mathcal{G})
$$
ここで等号は、テンソル--Hom 随伴と等式
$$
\NL_{g^{-1}\mathcal{O}_X/g^{-1}f^{-1}\mathcal{O}_S} = g^{-1}\NL_{X/S}
\quad\text{and}\quad
Lg^*\NL_{X/S} =
g^{-1}\NL_{X/S} \otimes_{g^{-1}\mathcal{O}_X}^\mathbf{L} \mathcal{O}_Y
$$
から従う。第一の等式については『加群』補題
\ref{modules-lemma-pullback-NL} を見よ。第二の等式は導来引き戻しの
定義から従う。
したがって、補題
\ref{defos-lemma-extensions-of-relative-ringed-spaces-functorial} を示すには、
$\mathcal{O}'_1$ が $\xi$ の像に対応し、$\mathcal{O}'_2$ が $\zeta$ の像に
対応することを示せば十分である。$\xi$ と $\mathcal{O}'_1$ の対応は、
注意 \ref{defos-remark-extensions-of-relative-ringed-spaces} における類 $\xi$ の
構成から直ちに従う。$\zeta$ と $\mathcal{O}'_2$ の対応については、まず
可換図式
$$
\xymatrix{
\mathcal{E} \ar[d]_{\beta'} \ar[rd]^\beta \\
\mathcal{O}_{Y'} \ar[r] & \mathcal{O}_Y
}
$$
を、$g^{-1}f^{-1}\mathcal{O}_S[\mathcal{E}] \to \mathcal{O}_Y$ が核
$\mathcal{K}$ をもつ全射となるように選ぶ。次に可換図式
$$
\xymatrix{
\mathcal{E} \ar[d]_{\beta'} &
\mathcal{E}' \ar[l]^\varphi \ar[d]_{\alpha'} \ar[rd]^\alpha \\
\mathcal{O}_{Y'} &
\mathcal{O}'_2 \ar[l] \ar[r] &
g^{-1}\mathcal{O}_X
}
$$
を、$g^{-1}f^{-1}\mathcal{O}_S[\mathcal{E}'] \to
g^{-1}\mathcal{O}_X$ が核 $\mathcal{J}$ をもつ全射となるように選ぶ。
（例えば集合の層として
$\mathcal{E}' = \mathcal{E} \amalg \mathcal{O}'_2$ とすればよい。）
写像 $\varphi$ は複体の写像 $\NL(\alpha) \to \NL(\beta)$ を誘導し
（記法は『加群』節 \ref{modules-section-netherlander} と同様）、特に
$\bar\varphi : \mathcal{J}/\mathcal{J}^2 \to \mathcal{K}/\mathcal{K}^2$ を
誘導する。このとき $\NL(\alpha) \cong \NL_{Y/S}$ および
$\NL(\beta) \cong \NL_{g^{-1}\mathcal{O}_X/g^{-1}f^{-1}\mathcal{O}_S}$ であり、
複体の写像 $\NL(\alpha) \to \NL(\beta)$ は、補題
\ref{defos-lemma-extensions-of-relative-ringed-spaces-functorial} の主張で用いた写像
$Lg^*\NL_{X/S} \to \NL_{Y/S}$ を表す（その証明の前半を見よ）。
ここで $\zeta$ は、$\beta'$ が誘導する写像
$\mathcal{K}/\mathcal{K}^2 \to \mathcal{G}$ の類に対応する。注意
\ref{defos-remark-extensions-of-relative-ringed-spaces} を見よ。同様に、拡大
$\mathcal{O}'_2$ は $\alpha'$ が誘導する写像
$\mathcal{J}/\mathcal{J}^2 \to \mathcal{G}$ に対応する。上の可換図式から、
この写像は、$\beta'$ が誘導する写像
$\mathcal{K}/\mathcal{K}^2 \to \mathcal{G}$ と写像
$\bar\varphi : \mathcal{J}/\mathcal{J}^2 \to \mathcal{K}/\mathcal{K}^2$
との合成である。これで求めていた両立性が示された。
\end{remark}

\begin{lemma}
\label{defos-lemma-parametrize-solutions-ringed-spaces}
$t : (S, \mathcal{O}_S) \to (S', \mathcal{O}_{S'})$、
$\mathcal{J} = \Ker(t^\sharp)$、
$f : (X, \mathcal{O}_X) \to (S, \mathcal{O}_S)$、$\mathcal{G}$、および
$c : \mathcal{J} \to \mathcal{G}$ は
(\ref{defos-equation-to-solve-ringed-spaces}) と同様とする。補題
\ref{defos-lemma-extensions-of-relative-ringed-spaces} により、
$\mathcal{O}_S$ の $\mathcal{J}$ による拡大 $\mathcal{O}_{S'}$ に対応する元を
$\xi \in \Ext^1_{\mathcal{O}_S}(\NL_{S/S'}, \mathcal{J})$ と記す。
解の同型類全体の集合は、正準に写像
$$
\Ext^1_{\mathcal{O}_X}(\NL_{X/S'}, \mathcal{G}) \to
\Ext^1_{\mathcal{O}_X}(Lf^*\NL_{S/S'}, \mathcal{G})
$$
の $\xi$ の像上のファイバーと全単射する。
\end{lemma}

\begin{proof}
補題 \ref{defos-lemma-extensions-of-relative-ringed-spaces} を $X \to S'$ と
$\mathcal{O}_X$-加群 $\mathcal{G}$ に適用すると、
$\Ext^1_{\mathcal{O}_X}(\NL_{X/S'}, \mathcal{G})$ の元 $\zeta$ は
$f^{-1}\mathcal{O}_{S'}$-代数の拡大
$0 \to \mathcal{G} \to \mathcal{O}_{X'} \to \mathcal{O}_X \to 0$ を
パラメータ付けることが分かる。補題
\ref{defos-lemma-extensions-of-relative-ringed-spaces-functorial} を
$X \to S \to S'$ と $c : \mathcal{J} \to \mathcal{G}$ に適用すると、
$c$ および $f : X \to S$ と両立する $S'$-射 $X' \to S'$ が存在する
ための必要十分条件は、$\zeta$ が $\xi$ に写ることである。
もちろん、これは $\mathcal{O}_{X'}$ が
(\ref{defos-equation-to-solve-ringed-spaces}) の解であることと同じである。
\end{proof}

\begin{remark}
\label{defos-remark-parametrize-solutions-ringed-spaces}
補題 \ref{defos-lemma-parametrize-solutions-ringed-spaces} の状況では、複体の写像
$$
Lf^*\NL_{S'/S} \to \NL_{X/S'} \to \NL_{X/S}
$$
がある。これらの写像は完全三角形をなすことに近い。『加群』補題
\ref{modules-lemma-exact-sequence-NL-ringed-topoi} を見よ。もしこれが
完全三角形であれば、$\Ext^2_{\mathcal{O}_X}(\NL_{X/S}, \mathcal{G})$
における $\xi$ の像が、(\ref{defos-equation-to-solve-ringed-spaces}) の解の
存在に対する障害であると結論できるだろう。
\end{remark}













\section{スキームの変形}
\label{defos-section-deformations-schemes}

\noindent
本節では、節 \ref{defos-section-deformations-ringed-spaces} の結果が
スキームの変形について何を意味するかを明記する。

\begin{lemma}
\label{defos-lemma-deform}
$S \subset S'$ をスキームの一次厚化とし、$f : X \to S$ をスキームの
平坦射とする。スキームの平坦射 $f' : X' \to S'$ と $S$ 上の同型
$a : X \to X' \times_{S'} S$ が存在するならば、次が成り立つ。
\begin{enumerate}
\item 組 $(f' : X' \to S', a)$ の同型類全体の集合は
$\Ext^1_{\mathcal{O}_X}(\NL_{X/S}, f^*\mathcal{C}_{S/S'})$
の下で主等質空間をなす。
\item $X' \times_{S'} S$ 上で恒等写像に帰着する $S'$ 上の自己同型
$\varphi : X' \to X'$ 全体の集合は
$\Ext^0_{\mathcal{O}_X}(\NL_{X/S}, f^*\mathcal{C}_{S/S'})$ である。
\end{enumerate}
\end{lemma}

\begin{proof}
まず、『射詳論』節 \ref{more-morphisms-section-thickenings} で定義される
スキームの厚化は、節 \ref{defos-section-thickenings-spaces} の意味で厚化となる
スキームの射と同じものであることに注意する。$X$ を $X'$ の閉部分スキームと
みなせるので、$(f, f') : (X \subset X') \to (S \subset S')$ は
一次厚化の射となる。このとき『射詳論』補題
\ref{more-morphisms-lemma-deform}（または、より一般的な補題
\ref{defos-lemma-deform-module}）から、$X'$ における $X$ のイデアル層は
$f^*\mathcal{C}_{S/S'}$ に等しいことが分かる。したがって可換図式
$$
\xymatrix{
0 \ar[r] & f^*\mathcal{C}_{S/S'} \ar[r] &
\mathcal{O}_{X'} \ar[r] &
\mathcal{O}_X \ar[r] & 0 \\
0 \ar[r] & \mathcal{C}_{S/S'} \ar[u] \ar[r] &
\mathcal{O}_{S'} \ar[u] \ar[r] &
\mathcal{O}_S \ar[u] \ar[r] & 0
}
$$
を得る。ここで垂直の矢印は $f$-写像である。
(\ref{defos-equation-to-solve-ringed-spaces}) と比較せよ。したがって (1) は
補題 \ref{defos-lemma-choices-ringed-spaces} から、(2) は補題
\ref{defos-lemma-huge-diagram-ringed-spaces} の (2) から従う。
（『射詳論』節 \ref{more-morphisms-section-netherlander} でスキームの射に
対して定義される $\NL_{X/S}$ は、節
\ref{defos-section-deformations-ringed-spaces} で用いた $\NL_{X/S}$ と一致する
ことに注意せよ。）
\end{proof}










\section{環付きトポスの厚化}
\label{defos-section-thickenings-ringed-topoi}

\noindent
本節は、環付きトポスについての節 \ref{defos-section-thickenings-spaces} の類似である。
以下のいくつかの節では、次の概念を用いる。
\begin{enumerate}
\item 環付きトポス $(\Sh(\mathcal{D}), \mathcal{O}')$ 上のイデアル層
$\mathcal{I} \subset \mathcal{O}'$ が {\it 局所冪零} であるとは、
$\mathcal{I}$ の任意の局所切断が局所的に冪零であることをいう。
\item 環付きトポスの {\it 厚化} とは、環付きトポスの射
$i : (\Sh(\mathcal{C}), \mathcal{O}) \to (\Sh(\mathcal{D}), \mathcal{O}')$
であって次を満たすものをいう。
\begin{enumerate}
\item $i_*$ は同値 $\Sh(\mathcal{C}) \to \Sh(\mathcal{D})$ である。
\item 写像 $i^\sharp : \mathcal{O}' \to i_*\mathcal{O}$ は全射である。
\item $i^\sharp$ の核は局所冪零なイデアル層である。
\end{enumerate}
\item 環付きトポスの {\it 一次厚化} とは、
$\Ker(i^\sharp)$ が二乗零である環付きトポスの厚化
$i : (\Sh(\mathcal{C}), \mathcal{O}) \to (\Sh(\mathcal{D}), \mathcal{O}')$
をいう。
\item {\it 環付きトポスの厚化の射}、
{\it 基礎環付きトポス上の環付きトポスの厚化の射} などの定義は明らかである。
\end{enumerate}
もし
$i : (\Sh(\mathcal{C}), \mathcal{O}) \to (\Sh(\mathcal{D}), \mathcal{O}')$
が環付きトポスの厚化ならば、台となるトポスを同一視し、
$\mathcal{O}$、$\mathcal{O}'$、および $\mathcal{I} = \Ker(i^\sharp)$ を
$\mathcal{C}$ 上の層とみなす。このとき $\mathcal{O}'$-加群の短完全列
$$
0 \to \mathcal{I} \to \mathcal{O}' \to \mathcal{O} \to 0
$$
を得る。『サイト上の加群』補題
\ref{sites-modules-lemma-i-star-equivalence} により、$\mathcal{O}$-加群の圏は、
$\mathcal{I}$ によって零化される $\mathcal{O}'$-加群の圏と同値である。
特に $i$ が一次厚化ならば、$\mathcal{I}$ は $\mathcal{O}$-加群である。

\begin{situation}
\label{defos-situation-morphism-thickenings-ringed-topoi}
環付きトポスの厚化の射 $(f, f')$ は可換図式
\begin{equation}
\label{defos-equation-morphism-thickenings-ringed-topoi}
\vcenter{
\xymatrix{
(\Sh(\mathcal{C}), \mathcal{O}) \ar[r]_i \ar[d]_f &
(\Sh(\mathcal{D}), \mathcal{O}') \ar[d]^{f'} \\
(\Sh(\mathcal{B}), \mathcal{O}_\mathcal{B}) \ar[r]^t &
(\Sh(\mathcal{B}'), \mathcal{O}_{\mathcal{B}'})
}
}
\end{equation}
で与えられ、横の矢印は厚化である。この状況で
$\mathcal{I} = \Ker(i^\sharp) \subset \mathcal{O}'$ および
$\mathcal{J} = \Ker(t^\sharp) \subset \mathcal{O}_{\mathcal{B}'}$
とおく。台となるトポス上では $f = f'$ なので、引き戻し関手
$f^{-1}$ と $(f')^{-1}$ を同一視する。
$(f')^\sharp : f^{-1}\mathcal{O}_{\mathcal{B}'} \to \mathcal{O}'$
は特に写像 $f^{-1}\mathcal{J} \to \mathcal{I}$ を誘導し、したがって
$\mathcal{O}'$-加群の写像
$$
(f')^*\mathcal{J} \longrightarrow \mathcal{I}
$$
を誘導することに注意する。$i$ と $t$ が一次厚化ならば、
$(f')^*\mathcal{J} = f^*\mathcal{J}$ であり、上の写像は
$f^*\mathcal{J} \to \mathcal{I}$ となる。
\end{situation}

\begin{definition}
\label{defos-definition-strict-morphism-thickenings-ringed-topoi}
状況 \ref{defos-situation-morphism-thickenings-ringed-topoi} において、写像
$(f')^*\mathcal{J} \longrightarrow \mathcal{I}$ が全射であるとき、
$(f, f')$ を {\it 厳密な厚化の射} と呼ぶ。
\end{definition}








\section{環付きトポスの一次厚化上の加群}
\label{defos-section-modules-thickenings-ringed-topoi}

\noindent
本節では、加群の変形理論に必要ないくつかの準備事項を論じる。
$i : (\Sh(\mathcal{C}, \mathcal{O}) \to (\Sh(\mathcal{D}), \mathcal{O}')$
を環付きトポスの一次厚化とする。節
\ref{defos-section-thickenings-ringed-topoi} で導入した記法を断りなく用い、
特に台となるトポスを同一視する。本節では、短完全列
\begin{equation}
\label{defos-equation-extension-ringed-topoi}
0 \to \mathcal{K} \to \mathcal{F}' \to \mathcal{F} \to 0
\end{equation}
を考える。これは $\mathcal{O}'$-加群の列であり、$\mathcal{F}$ と
$\mathcal{K}$ は $\mathcal{O}$-加群、$\mathcal{F}'$ は
$\mathcal{O}'$-加群である。この状況では、正準な $\mathcal{O}$-加群準同型
$$
c_{\mathcal{F}'} :
\mathcal{I} \otimes_\mathcal{O} \mathcal{F}
\longrightarrow
\mathcal{K}
$$
がある。ここで $\mathcal{I} = \Ker(i^\sharp)$ である。
実際、$\mathcal{I}$ の局所切断 $f$ と $\mathcal{F}$ の局所切断 $s$ に対し、
$s$ を持ち上げる $\mathcal{F}'$ の局所切断を $s'$ として
$c_{\mathcal{F}'}(f \otimes s) = fs'$ と定める。

\begin{lemma}
\label{defos-lemma-inf-map-ringed-topoi}
$i : (\Sh(\mathcal{C}), \mathcal{O}) \to (\Sh(\mathcal{D}), \mathcal{O}')$
を環付きトポスの一次厚化とする。次の拡大が与えられていると仮定する。
$$
0 \to \mathcal{K} \to \mathcal{F}' \to \mathcal{F} \to 0
\quad\text{and}\quad
0 \to \mathcal{L} \to \mathcal{G}' \to \mathcal{G} \to 0
$$
これらは (\ref{defos-equation-extension-ringed-topoi}) の形であり、さらに写像
$\varphi : \mathcal{F} \to \mathcal{G}$ と
$\psi : \mathcal{K} \to \mathcal{L}$ が与えられているとする。
\begin{enumerate}
\item $\varphi$ および $\psi$ と両立する $\mathcal{O}'$-加群準同型
$\varphi' : \mathcal{F}' \to \mathcal{G}'$ が存在すれば、次の図式は
可換である。
$$
\xymatrix{
\mathcal{I} \otimes_\mathcal{O} \mathcal{F}
\ar[r]_-{c_{\mathcal{F}'}} \ar[d]_{1 \otimes \varphi} &
\mathcal{K} \ar[d]^\psi \\
\mathcal{I} \otimes_\mathcal{O} \mathcal{G}
\ar[r]^-{c_{\mathcal{G}'}} &
\mathcal{L}
}
$$
\item $\varphi$ および $\psi$ と両立する $\mathcal{O}'$-加群準同型
$\varphi' : \mathcal{F}' \to \mathcal{G}'$ 全体の集合は、空でなければ
$\Hom_\mathcal{O}(\mathcal{F}, \mathcal{L})$ の下で主等質空間をなす。
\end{enumerate}
\end{lemma}

\begin{proof}
(1) は写像の記述から直ちに従う。(2) について、$\varphi'$ と
$\varphi''$ が $\varphi$ および $\psi$ と両立する二つの写像
$\mathcal{F}' \to \mathcal{G}'$ ならば、$\varphi' - \varphi''$ は
次のように分解する。
$$
\mathcal{F}' \to \mathcal{F} \to \mathcal{L} \to \mathcal{G}'
$$
『サイト上の加群』補題 \ref{sites-modules-lemma-i-star-equivalence} により、
中間の写像は $\Hom_\mathcal{O}(\mathcal{F}, \mathcal{L})$ の一意な元から
生じる。逆に、この群の元 $\alpha$ が与えられれば、上の表示で中間に
$\alpha$ を置いた合成を $\varphi'$ に加えることができる。細部は省略する。
\end{proof}

\begin{lemma}
\label{defos-lemma-inf-obs-map-ringed-topoi}
$i : (\Sh(\mathcal{C}), \mathcal{O}) \to (\Sh(\mathcal{D}), \mathcal{O}')$
を環付きトポスの一次厚化とする。次の拡大が与えられていると仮定する。
$$
0 \to \mathcal{K} \to \mathcal{F}' \to \mathcal{F} \to 0
\quad\text{and}\quad
0 \to \mathcal{L} \to \mathcal{G}' \to \mathcal{G} \to 0
$$
これらは (\ref{defos-equation-extension-ringed-topoi}) の形であり、さらに写像
$\varphi : \mathcal{F} \to \mathcal{G}$ と
$\psi : \mathcal{K} \to \mathcal{L}$ が与えられているとする。図式
$$
\xymatrix{
\mathcal{I} \otimes_\mathcal{O} \mathcal{F}
\ar[r]_-{c_{\mathcal{F}'}} \ar[d]_{1 \otimes \varphi} &
\mathcal{K} \ar[d]^\psi \\
\mathcal{I} \otimes_\mathcal{O} \mathcal{G}
\ar[r]^-{c_{\mathcal{G}'}} &
\mathcal{L}
}
$$
が可換であると仮定する。このとき元
$$
o(\varphi, \psi) \in
\Ext^1_\mathcal{O}(\mathcal{F}, \mathcal{L})
$$
が存在し、その消滅は、$\varphi$ および $\psi$ と両立する写像
$\varphi' : \mathcal{F}' \to \mathcal{G}'$ が存在するための
必要十分条件である。
\end{lemma}

\begin{proof}
次の拡大を明示的に構成できる。
$$
0 \to \mathcal{L} \to \mathcal{H} \to \mathcal{F} \to 0
$$
$\mathcal{H}$ を複体
$$
\mathcal{K}
\xrightarrow{1, - \psi}
\mathcal{F}' \oplus \mathcal{G}' \xrightarrow{\varphi, 1}
\mathcal{G}
$$
の中間のコホモロジーとすればよい（記法は明らかであろう）。
補題の図式が可換であるという仮定を用いた局所切断の計算により、
$\mathcal{H}$ は $\mathcal{I}$ によって零化される。したがって
$\mathcal{H}$ は次の群の類を定める。
$$
\Ext^1_\mathcal{O}(\mathcal{F}, \mathcal{L})
\subset
\Ext^1_{\mathcal{O}'}(\mathcal{F}, \mathcal{L})
$$
最後に、$\mathcal{H}$ の類は、拡大 $\mathcal{F}'$ の $\psi$ に沿う
押し出しと、拡大 $\mathcal{G}'$ の $\varphi$ に沿う引き戻しとの差である
（計算は省略する）。したがって $\mathcal{H}$ の類が消滅することは、
次の可換図式が存在することと同値である。
$$
\xymatrix{
0 \ar[r] &
\mathcal{K} \ar[r] \ar[d]_{\psi} &
\mathcal{F}' \ar[r] \ar[d]_{\varphi'} &
\mathcal{F} \ar[r] \ar[d]_\varphi & 0\\
0 \ar[r] &
\mathcal{L} \ar[r] &
\mathcal{G}' \ar[r] &
\mathcal{G} \ar[r] & 0
}
$$
これは所望の結論である。
\end{proof}

\begin{lemma}
\label{defos-lemma-inf-ext-ringed-topoi}
$i : (\Sh(\mathcal{C}), \mathcal{O}) \to (\Sh(\mathcal{D}), \mathcal{O}')$
を環付きトポスの一次厚化とする。$\mathcal{O}$-加群
$\mathcal{F}$、$\mathcal{K}$ と $\mathcal{O}$-線形写像
$c : \mathcal{I} \otimes_\mathcal{O} \mathcal{F} \to \mathcal{K}$
が与えられているとする。$c_{\mathcal{F}'} = c$ を満たす列
(\ref{defos-equation-extension-ringed-topoi}) が存在するならば、そのような拡大の
同型類全体の集合は $\Ext^1_\mathcal{O}(\mathcal{F}, \mathcal{K})$
の下で主等質空間をなす。
\end{lemma}

\begin{proof}
次の拡大が与えられていると仮定する。
$$
0 \to \mathcal{K} \to \mathcal{F}'_1 \to \mathcal{F} \to 0
\quad\text{and}\quad
0 \to \mathcal{K} \to \mathcal{F}'_2 \to \mathcal{F} \to 0
$$
$c_{\mathcal{F}'_1} = c_{\mathcal{F}'_2} = c$ とする。このとき、拡大群に
おける差（『ホモロジー』節 \ref{homology-section-extensions} を見よ）は
拡大
$$
0 \to \mathcal{K} \to \mathcal{E} \to \mathcal{F} \to 0
$$
であり、$\mathcal{E}$ は $\mathcal{I}$ によって零化される
（局所計算は省略する）。したがってこの列は $\mathcal{O}$-加群の拡大である。
『サイト上の加群』補題 \ref{sites-modules-lemma-i-star-equivalence} を見よ。
逆に、そのような拡大 $\mathcal{E}$ が与えられれば、写像
$c_{\mathcal{F}'}$ を変えずに、拡大 $\mathcal{E}$ を
$\mathcal{O}'$-拡大 $\mathcal{F}'$ に加えることができる。細部は省略する。
\end{proof}

\begin{lemma}
\label{defos-lemma-inf-obs-ext-ringed-topoi}
$i : (\Sh(\mathcal{C}), \mathcal{O}) \to (\Sh(\mathcal{D}), \mathcal{O}')$
を環付きトポスの一次厚化とする。$\mathcal{O}$-加群
$\mathcal{F}$、$\mathcal{K}$ と $\mathcal{O}$-線形写像
$c : \mathcal{I} \otimes_\mathcal{O} \mathcal{F} \to \mathcal{K}$
が与えられているとする。このとき元
$$
o(\mathcal{F}, \mathcal{K}, c) \in
\Ext^2_\mathcal{O}(\mathcal{F}, \mathcal{K})
$$
が存在し、その消滅は、$c_{\mathcal{F}'} = c$ を満たす列
(\ref{defos-equation-extension-ringed-topoi}) が存在するための必要十分条件である。
\end{lemma}

\begin{proof}
まず、$\mathcal{K}$ が入射 $\mathcal{O}$-加群ならば、
$c_{\mathcal{F}'} = c$ を満たす列 (\ref{defos-equation-extension-ringed-topoi}) が
実際に存在することを示す。そのため、平坦 $\mathcal{O}'$-加群
$\mathcal{H}'$ と全射 $\mathcal{H}' \to \mathcal{F}$ を選ぶ
（『サイト上の加群』補題
\ref{sites-modules-lemma-module-quotient-flat}）。その核を
$\mathcal{J} \subset \mathcal{H}'$ とする。$\mathcal{H}'$ は平坦なので、
$$
\mathcal{I} \otimes_{\mathcal{O}'} \mathcal{H}' =
\mathcal{I}\mathcal{H}'
\subset \mathcal{J} \subset \mathcal{H}'
$$
写像
$$
\mathcal{I}\mathcal{H}' =
\mathcal{I} \otimes_{\mathcal{O}'} \mathcal{H}'
\longrightarrow
\mathcal{I} \otimes_{\mathcal{O}'} \mathcal{F} =
\mathcal{I} \otimes_\mathcal{O} \mathcal{F}
$$
が $\mathcal{I}\mathcal{J}$ を零化することに注意する。実際、$f$ が
$\mathcal{I}$ の局所切断、$s$ が $\mathcal{H}$ の局所切断ならば、
$fs$ は $f \otimes \overline{s}$ に写る。ここで $\overline{s}$ は
$s$ の $\mathcal{F}$ における像である。したがって、
$$
\xymatrix{
\mathcal{I}\mathcal{H}'/\mathcal{I}\mathcal{J}
\ar@{^{(}->}[r] \ar[d] &
\mathcal{J}/\mathcal{I}\mathcal{J} \ar@{..>}[d]_\gamma \\
\mathcal{I} \otimes_\mathcal{O} \mathcal{F} \ar[r]^-c &
\mathcal{K}
}
$$
という $\mathcal{O}$-加群の図式を得る。$\mathcal{K}$ が
$\mathcal{O}$-加群として入射ならば、点線の矢印が得られる。
$\gamma$ と $\mathcal{J} \to \mathcal{J}/\mathcal{I}\mathcal{J}$ の合成を
$\gamma' : \mathcal{J} \to \mathcal{K}$ と記す。局所計算により、押し出し
$$
\xymatrix{
0 \ar[r] &
\mathcal{J} \ar[r] \ar[d]_{\gamma'} &
\mathcal{H}' \ar[r] \ar[d] &
\mathcal{F} \ar[r] \ar@{=}[d] &
0 \\
0 \ar[r] &
\mathcal{K} \ar[r] &
\mathcal{F}' \ar[r] &
\mathcal{F} \ar[r] &
0
}
$$
は補題で提起された問題の解を与える。

\medskip\noindent
一般の場合。$\mathcal{K}'$ が入射 $\mathcal{O}$-加群となる埋め込み
$\mathcal{K} \subset \mathcal{K}'$ を選ぶ。$\mathcal{Q}$ をその商とする。このとき完全列
$$
0 \to \mathcal{K} \to \mathcal{K}' \to \mathcal{Q} \to 0
$$
を得る。合成を
$c' : \mathcal{I} \otimes_\mathcal{O} \mathcal{F} \to \mathcal{K}'$
と書く。上の段落により、$c_{\mathcal{E}'} = c'$ を満たす
(
\ref{defos-equation-extension-ringed-topoi}) の形の列
$$
0 \to \mathcal{K}' \to \mathcal{E}' \to \mathcal{F} \to 0
$$
が存在する。$c'$ と写像 $\mathcal{K}' \to \mathcal{Q}$ の合成は
零であることに注意する。したがって、$\mathcal{K}' \to \mathcal{Q}$
による $\mathcal{E}'$ の押し出しは、$c_{\mathcal{D}'} = 0$ を満たす
(\ref{defos-equation-extension-ringed-topoi}) の形の拡大
$$
0 \to \mathcal{Q} \to \mathcal{D}' \to \mathcal{F} \to 0
$$
である。これはちょうど $\mathcal{D}'$ が $\mathcal{I}$ によって零化される
ことを意味する。言い換えれば、$\mathcal{D}'$ は $\mathcal{O}$-加群の拡大であり、
したがって元
$$
o(\mathcal{F}, \mathcal{K}, c) \in
\Ext^1_\mathcal{O}(\mathcal{F}, \mathcal{Q}) =
\Ext^2_\mathcal{O}(\mathcal{F}, \mathcal{K})
$$
を定める（等号は、上の完全列に付随する長完全コホモロジー列と、
入射加群 $\mathcal{K}'$ を値にとる高次 Ext 群の消滅から従う）。
$o(\mathcal{F}, \mathcal{K}, c) = 0$ ならば、分裂
$s : \mathcal{F} \to \mathcal{D}'$ を選び、
$$
\mathcal{F}' = \Ker(\mathcal{E}' \to \mathcal{D}'/s(\mathcal{F}))
$$
と置くことができる。すると次の図式
$$
\xymatrix{
0 \ar[r] &
\mathcal{K} \ar[r] \ar[d] &
\mathcal{F}' \ar[r] \ar[d] &
\mathcal{F} \ar[r] \ar@{=}[d] &
0 \\
0 \ar[r] &
\mathcal{K}' \ar[r] &
\mathcal{E}' \ar[r] &
\mathcal{F} \ar[r] & 0
}
$$
を得、その各行は完全である。これにより $c_{\mathcal{F}'} = c$ が分かる。
逆に $\mathcal{F}'$ が存在すると仮定する。補題
\ref{defos-lemma-inf-ext-ringed-topoi} と、入射加群 $\mathcal{K}'$ を値にとる
高次 Ext 群の消滅により、写像 $\mathcal{K} \to \mathcal{K}'$ による
$\mathcal{F}'$ の押し出しは $\mathcal{E}'$ と同型である。
これから上と同じ図式が得られ、$\mathcal{D}'$ は拡大として分裂する。
すなわち類 $o(\mathcal{F}, \mathcal{K}, c)$ は零である。
\end{proof}

\begin{remark}
\label{defos-remark-trivial-thickening-ringed-topoi}
$(\Sh(\mathcal{C}), \mathcal{O})$ を環付きトポスとする。一次厚化
$i : (\Sh(\mathcal{C}), \mathcal{O}) \to
(\Sh(\mathcal{D}), \mathcal{O}')$ に対し、環付きトポスの射
$\pi : (\Sh(\mathcal{D}), \mathcal{O}') \to (\Sh(\mathcal{C}), \mathcal{O})$
で $i$ の左逆射となるものが存在するとき、この厚化を {\it 自明} という。
このような射 $\pi$ の選択を、一次厚化の {\it 自明化} という。
$\pi$ が与えられると、分裂
\begin{equation}
\label{defos-equation-splitting-ringed-topoi}
\mathcal{O}' = \mathcal{O} \oplus \mathcal{I}
\end{equation}
が $\mathcal{C}$ 上の代数の層として得られる。ここでは $\pi^\sharp$ を用いて
全射 $\mathcal{O}' \to \mathcal{O}$ を分裂させた。逆に、このような分裂は
射 $\pi$ を定める。$(\Sh(\mathcal{C}), \mathcal{O})$ の自明化された
一次厚化の圏は、$\mathcal{O}$-加群の圏と同値である。
\end{remark}

\begin{remark}
\label{defos-remark-trivial-extension-ringed-topoi}
環付きトポスの自明な一次厚化
$i : (\Sh(\mathcal{C}), \mathcal{O}) \to (\Sh(\mathcal{D}), \mathcal{O}')$
と、その自明化 $\pi : (\Sh(\mathcal{D}), \mathcal{O}') \to
(\Sh(\mathcal{C}), \mathcal{O})$ をとる。二つの $\mathcal{O}$-加群と写像
$c : \mathcal{I} \otimes_\mathcal{O} \mathcal{F} \to \mathcal{K}$
からなる任意の三つ組 $(\mathcal{F}, \mathcal{K}, c)$ に対し、
$$
\mathcal{F}'_{c, triv} = \mathcal{F} \oplus \mathcal{K}
$$
$\pi$ に付随する分裂 (\ref{defos-equation-splitting-ringed-topoi}) と写像 $c$
を用いて $\mathcal{O}'$-加群構造を定めると、拡大
(\ref{defos-equation-extension-ringed-topoi}) が得られる。
$\mathcal{F}'_{c, triv}$ を、$c$ と自明化 $\pi$ に対応する
$\mathcal{F}$ の $\mathcal{K}$ による {\it 自明な拡大} と呼ぶ。
(\ref{defos-equation-extension-ringed-topoi}) の形の任意の拡大 $\mathcal{F}'$ に対し、
$\pi^\sharp : \mathcal{O} \to \mathcal{O}'$ を用いて $\mathcal{F}'$ を
$\mathcal{O}$-加群の拡大とみなせる。したがって
$\Ext^1_\mathcal{O}(\mathcal{F}, \mathcal{K})$ の類
$\xi_{\mathcal{F}'}$ を得る。補題 \ref{defos-lemma-inf-ext-ringed-topoi} により、
$\mathcal{F}' \mapsto \xi_{\mathcal{F}'}$ は全単射
$$
\left\{
\begin{matrix}
\text{拡大 }\mathcal{F}'\text{ の同型類で、}\\
\text{(\ref{defos-equation-extension-ringed-topoi}) の形かつ }
c = c_{\mathcal{F}'}\text{ を満たすもの}
\end{matrix}
\right\}
\longrightarrow
\Ext^1_\mathcal{O}(\mathcal{F}, \mathcal{K})
$$
を誘導する。さらに、自明な拡大 $\mathcal{F}'_{c, triv}$ は零類に写る。
\end{remark}

\begin{remark}
\label{defos-remark-extension-functorial-ringed-topoi}
$(\Sh(\mathcal{C}), \mathcal{O})$ を環付きトポスとする。
$(\Sh(\mathcal{C}), \mathcal{O}) \to (\Sh(\mathcal{D}_i), \mathcal{O}'_i)$,
$i = 1, 2$ を、イデアル層 $\mathcal{I}_i$ をもつ一次厚化とする。
$h : (\Sh(\mathcal{D}_1), \mathcal{O}'_1) \to
(\Sh(\mathcal{D}_2), \mathcal{O}'_2)$
を $(\Sh(\mathcal{C}), \mathcal{O})$ の一次厚化の射とする。図式は
$$
\xymatrix{
& (\Sh(\mathcal{C}), \mathcal{O}) \ar[ld] \ar[rd] & \\
(\Sh(\mathcal{D}_1), \mathcal{O}'_1) \ar[rr]^h & & 
(\Sh(\mathcal{D}_2), \mathcal{O}'_2)
}
$$
である。特に $h^\sharp : \mathcal{O}'_2 \to \mathcal{O}'_1$ は
$\mathcal{O}$-加群の写像 $\mathcal{I}_2 \to \mathcal{I}_1$ を誘導する。
$\mathcal{F}$ を $\mathcal{O}$-加群とする。
各 $i = 1, 2$ に対し、$\mathcal{O}$-加群 $\mathcal{K}_i$ と写像
$c_i : \mathcal{I}_i \otimes_\mathcal{O} \mathcal{F} \to
\mathcal{K}_i$ からなる組 $(\mathcal{K}_i, c_i)$ をとる。さらに、
次の図式を可換にする $\mathcal{O}$-加群の写像
$\mathcal{K}_2 \to \mathcal{K}_1$ が与えられていると仮定する。
$$
\xymatrix{
\mathcal{I}_2 \otimes_\mathcal{O} \mathcal{F}
\ar[r]_-{c_2} \ar[d] &
\mathcal{K}_2 \ar[d] \\
\mathcal{I}_1 \otimes_\mathcal{O} \mathcal{F}
\ar[r]^-{c_1} &
\mathcal{K}_1
}
$$
このとき、標準的な関手的写像
$$
\left\{
\begin{matrix}
\mathcal{F}'_2\text{ は (\ref{defos-equation-extension-ringed-topoi}) の形で、}\\
c_2 = c_{\mathcal{F}'_2}\text{ かつ }\mathcal{K} = \mathcal{K}_2
\end{matrix}
\right\}
\longrightarrow
\left\{
\begin{matrix}
\mathcal{F}'_1\text{ は (\ref{defos-equation-extension-ringed-topoi}) の形で、}\\
c_1 = c_{\mathcal{F}'_1}\text{ かつ }\mathcal{K} = \mathcal{K}_1
\end{matrix}
\right\}
$$
が存在する。実際、すべての層 $\mathcal{O}$、$\mathcal{O}'_i$、
$\mathcal{F}$、$\mathcal{K}_i$ などを $\mathcal{C}$ 上の層とみなし、
$\mathcal{F}'_2$ が与えられたとき $\mathcal{F}'_1$ を押し出し、
すなわち次の拡大の図式に入る層として定める。
$$
\xymatrix{
0 \ar[r] &
\mathcal{K}_2 \ar[r] \ar[d] &
\mathcal{F}'_2 \ar[r] \ar[d] &
\mathcal{F} \ar@{=}[d] \ar[r] & 0 \\
0 \ar[r] &
\mathcal{K}_1 \ar[r] &
\mathcal{F}'_1 \ar[r] &
\mathcal{F} \ar[r] & 0
}
$$
押し出し上の $\mathcal{O}'_1$-加群構造の構成は省略する
（この構成では $c_1$ と $c_2$ を含む図式の可換性を用いる）。
\end{remark}

\begin{remark}
\label{defos-remark-trivial-extension-functorial-ringed-topoi}
$(\Sh(\mathcal{C}), \mathcal{O})$、
$(\Sh(\mathcal{C}), \mathcal{O}) \to (\Sh(\mathcal{D}_i), \mathcal{O}'_i)$,
$\mathcal{I}_i$、および $h : (\Sh(\mathcal{D}_1), \mathcal{O}'_1) \to
(\Sh(\mathcal{D}_2), \mathcal{O}'_2)$ は注意
\ref{defos-remark-extension-functorial-ringed-topoi} のとおりとする。
自明化
$\pi_i : (\Sh(\mathcal{D}_i), \mathcal{O}'_i) \to
(\Sh(\mathcal{C}), \mathcal{O})$ が与えられ、
$\pi_1 = h \circ \pi_2$ を満たすと仮定する。言い換えれば、$h$ は
$(\Sh(\mathcal{C}), \mathcal{O})$ の自明化された一次厚化の射である。
各 $i = 1, 2$ に対し、$\mathcal{O}$-加群 $\mathcal{K}_i$ と写像
$c_i : \mathcal{I}_i \otimes_\mathcal{O} \mathcal{F} \to
\mathcal{K}_i$ からなる組 $(\mathcal{K}_i, c_i)$ をとる。さらに、
次の図式を可換にする $\mathcal{O}$-加群の写像
$\mathcal{K}_2 \to \mathcal{K}_1$ が与えられていると仮定する。
$$
\xymatrix{
\mathcal{I}_2 \otimes_\mathcal{O} \mathcal{F}
\ar[r]_-{c_2} \ar[d] &
\mathcal{K}_2 \ar[d] \\
\mathcal{I}_1 \otimes_\mathcal{O} \mathcal{F}
\ar[r]^-{c_1} &
\mathcal{K}_1
}
$$
この状況では、注意 \ref{defos-remark-trivial-extension-ringed-topoi}
の構成により、可換図式
$$
\xymatrix{
\{\mathcal{F}'_2\text{ は (\ref{defos-equation-extension-ringed-topoi}) の形で、}
c_2 = c_{\mathcal{F}'_2}\text{ かつ }\mathcal{K} = \mathcal{K}_2\}
\ar[d] \ar[rr] & &
\Ext^1_\mathcal{O}(\mathcal{F}, \mathcal{K}_2) \ar[d] \\
\{\mathcal{F}'_1\text{ は (\ref{defos-equation-extension-ringed-topoi}) の形で、}
c_1 = c_{\mathcal{F}'_1}\text{ かつ }\mathcal{K} = \mathcal{K}_1\}
\ar[rr] & &
\Ext^1_\mathcal{O}(\mathcal{F}, \mathcal{K}_1)
}
$$
が誘導される。右側の縦写像は $\Ext$ の関手性と写像
$\mathcal{K}_2 \to \mathcal{K}_1$ によって与えられ、左側の縦写像は
注意 \ref{defos-remark-extension-functorial-ringed-topoi} の写像である。
\end{remark}

\begin{remark}
\label{defos-remark-obstruction-extension-functorial-ringed-topoi}
$(\Sh(\mathcal{C}), \mathcal{O})$、
$(\Sh(\mathcal{C}), \mathcal{O}) \to (\Sh(\mathcal{D}_i), \mathcal{O}'_i)$,
$\mathcal{I}_i$、および $h : (\Sh(\mathcal{D}_1), \mathcal{O}'_1) \to
(\Sh(\mathcal{D}_2), \mathcal{O}'_2)$ は注意
\ref{defos-remark-extension-functorial-ringed-topoi} のとおりとする。
$h^\sharp : \mathcal{O}'_2 \to \mathcal{O}'_1$ は特に
$\mathcal{O}$-加群の写像 $\mathcal{I}_2 \to \mathcal{I}_1$ を誘導する。
$\mathcal{F}$ を $\mathcal{O}$-加群とする。
各 $i = 1, 2$ に対し、$\mathcal{O}$-加群 $\mathcal{K}_i$ と写像
$c_i : \mathcal{I}_i \otimes_\mathcal{O} \mathcal{F} \to
\mathcal{K}_i$ からなる組 $(\mathcal{K}_i, c_i)$ をとる。さらに、
次の図式を可換にする $\mathcal{O}$-加群の写像
$\mathcal{K}_2 \to \mathcal{K}_1$ が与えられていると仮定する。
$$
\xymatrix{
\mathcal{I}_2 \otimes_\mathcal{O} \mathcal{F}
\ar[r]_-{c_2} \ar[d] &
\mathcal{K}_2 \ar[d] \\
\mathcal{I}_1 \otimes_\mathcal{O} \mathcal{F}
\ar[r]^-{c_1} &
\mathcal{K}_1
}
$$
このとき、写像
$$
\Ext^2_\mathcal{O}(\mathcal{F}, \mathcal{K}_2)
\longrightarrow
\Ext^2_\mathcal{O}(\mathcal{F}, \mathcal{K}_1)
$$
は $o(\mathcal{F}, \mathcal{K}_2, c_2)$ を
$o(\mathcal{F}, \mathcal{K}_1, c_1)$ に写すと {\bf 主張} する。

\medskip\noindent
この主張を証明するため、$\mathcal{K}_2'$ が入射 $\mathcal{O}$-加群となる埋め込み
$j_2 : \mathcal{K}_2 \to \mathcal{K}_2'$
を選ぶ。補題 \ref{defos-lemma-inf-obs-ext-ringed-topoi} の証明と同様に、
$c_{\mathcal{E}_2} = j_2 \circ c_2$ を満たす $\mathcal{O}_2$-加群の拡大
$$
0 \to \mathcal{K}_2' \to \mathcal{E}_2 \to \mathcal{F} \to 0
$$
を選べる。補題 \ref{defos-lemma-inf-obs-ext-ringed-topoi} の証明では、
$o(\mathcal{F}, \mathcal{K}_2, c_2)$
を、$\mathcal{O}$-加群の完全列
$$
0 \to
\mathcal{K}_2 \to \mathcal{K}_2' \to
\mathcal{E}_2/\mathcal{K}_2 \to
\mathcal{F} \to 0
$$
の Yoneda 拡大類（『導来圏』第\ref{derived-section-ext}節の意味で）
として構成する。$\mathcal{K}_1'$ を
$\mathcal{K}_2 \to \mathcal{K}_1 \oplus \mathcal{K}_2'$ の余核とする。
可換正方形をなす単射 $j_1 : \mathcal{K}_1 \to \mathcal{K}_1'$ と
写像 $\mathcal{K}_2' \to \mathcal{K}_1'$ が存在する。押し出しを作ると、
$$
\xymatrix{
0 \ar[r] &
\mathcal{K}_2' \ar[r] \ar[d] &
\mathcal{E}_2 \ar[r] \ar[d] &
\mathcal{F} \ar[r] \ar[d] & 0 \\
0 \ar[r] &
\mathcal{K}_1' \ar[r] &
\mathcal{E}_1 \ar[r] &
\mathcal{F} \ar[r] & 0
}
$$
$\mathcal{E}_1$ 上には標準的な $\mathcal{O}_1$-加群構造があり、
この構造について $c_{\mathcal{E}_1} = j_1 \circ c_1$ である
（ここでは上の $c_1$ と $c_2$ を含む図式の可換性を用いる）。
補題 \ref{defos-lemma-inf-obs-ext-ringed-topoi} の手順によれば、
$o(\mathcal{F}, \mathcal{K}_1, c_1)$ は $\mathcal{O}$-加群の完全列
$$
0 \to
\mathcal{K}_1 \to
\mathcal{K}_1' \to
\mathcal{E}_1/\mathcal{K}_1 \to
\mathcal{F} \to 0
$$
の Yoneda 拡大類である。完全列の写像
$$
\xymatrix{
0 \ar[r] &
\mathcal{K}_2 \ar[d] \ar[r] &
\mathcal{K}_2' \ar[d] \ar[r] &
\mathcal{E}_2/\mathcal{K}_2 \ar[r] \ar[d] &
\mathcal{F} \ar[r] \ar@{=}[d] &
0 \\
0 \ar[r] &
\mathcal{K}_2 \ar[r] &
\mathcal{K}_2' \ar[r] &
\mathcal{E}_2/\mathcal{K}_2 \ar[r] &
\mathcal{F} \ar[r] &
0
}
$$
があるので、主張が従う。
\end{remark}

\begin{remark}
\label{defos-remark-short-exact-sequence-thickenings-ringed-topoi}
$(\Sh(\mathcal{C}), \mathcal{O})$ を環付きトポスとする。
$(\Sh(\mathcal{C}), \mathcal{O})$ の一次厚化の射の列
$$
(\Sh(\mathcal{D}_1), \mathcal{O}'_1) \to
(\Sh(\mathcal{D}_2), \mathcal{O}'_2) \to
(\Sh(\mathcal{D}_3), \mathcal{O}'_3)
$$
について、対応するイデアル層 $\mathcal{I}_i$ の間の写像が
$\mathcal{O}$-加群の複体
$\mathcal{I}_3 \to \mathcal{I}_2 \to \mathcal{I}_1$
を与える（すなわち合成が零である）とき、この列を {\it 複体} と呼ぶ。
この場合、合成
$(\Sh(\mathcal{D}_1), \mathcal{O}'_1) \to
(\Sh(\mathcal{D}_3), \mathcal{O}'_3)$ は
$(\Sh(\mathcal{C}), \mathcal{O}) \to
(\Sh(\mathcal{D}_3), \mathcal{O}'_3)$ を経由する。すなわち、
$(\Sh(\mathcal{D}_1), \mathcal{O}'_1)$ は
$(\Sh(\mathcal{C}), \mathcal{O})$ の自明な一次厚化であり、標準的な自明化
$\pi : (\Sh(\mathcal{D}_1), \mathcal{O}'_1) \to
(\Sh(\mathcal{C}), \mathcal{O})$
を伴う。

\medskip\noindent
$(\Sh(\mathcal{C}), \mathcal{O})$ の一次厚化の射の列
$$
(\Sh(\mathcal{D}_1), \mathcal{O}'_1) \to
(\Sh(\mathcal{D}_2), \mathcal{O}'_2) \to
(\Sh(\mathcal{D}_3), \mathcal{O}'_3)
$$
について、対応するイデアル層の間の写像が $\mathcal{O}$-加群の短完全列
$$
0 \to \mathcal{I}_3 \to \mathcal{I}_2 \to \mathcal{I}_1 \to 0
$$
をなすとき、この列を {\it 短完全列} という。
\end{remark}

\begin{remark}
\label{defos-remark-complex-thickenings-and-ses-modules-ringed-topoi}
$(\Sh(\mathcal{C}), \mathcal{O})$ を環付きトポスとし、
$\mathcal{F}$ を $\mathcal{O}$-加群とする。
$$
(\Sh(\mathcal{D}_1), \mathcal{O}'_1) \to
(\Sh(\mathcal{D}_2), \mathcal{O}'_2) \to
(\Sh(\mathcal{D}_3), \mathcal{O}'_3)
$$
を $(\Sh(\mathcal{C}), \mathcal{O})$ の一次厚化の複体とする。注意
\ref{defos-remark-short-exact-sequence-thickenings-ringed-topoi} を参照せよ。
各 $i = 1, 2, 3$ に対し、$\mathcal{O}$-加群 $\mathcal{K}_i$ と写像
$c_i : \mathcal{I}_i \otimes_\mathcal{O} \mathcal{F} \to
\mathcal{K}_i$ からなる組 $(\mathcal{K}_i, c_i)$ をとる。
$\mathcal{O}$-加群の短完全列
$$
0 \to \mathcal{K}_3 \to \mathcal{K}_2 \to \mathcal{K}_1 \to 0
$$
が与えられ、次の二つの図式が可換であると仮定する。
$$
\vcenter{
\xymatrix{
\mathcal{I}_2 \otimes_\mathcal{O} \mathcal{F}
\ar[r]_-{c_2} \ar[d] &
\mathcal{K}_2 \ar[d] \\
\mathcal{I}_1 \otimes_\mathcal{O} \mathcal{F}
\ar[r]^-{c_1} &
\mathcal{K}_1
}
}
\quad\text{および}\quad
\vcenter{
\xymatrix{
\mathcal{I}_3 \otimes_\mathcal{O} \mathcal{F}
\ar[r]_-{c_3} \ar[d] &
\mathcal{K}_3 \ar[d] \\
\mathcal{I}_2 \otimes_\mathcal{O} \mathcal{F}
\ar[r]^-{c_2} &
\mathcal{K}_2
}
}
$$
最後に、(\ref{defos-equation-extension-ringed-topoi}) の形で
$\mathcal{K} = \mathcal{K}_2$ かつ $c_{\mathcal{F}'_2} = c_2$ を満たす
$\mathcal{O}'_2$-加群の拡大
$$
0 \to \mathcal{K}_2 \to \mathcal{F}'_2 \to \mathcal{F} \to 0
$$
が与えられていると仮定する。この状況で、注意
\ref{defos-remark-extension-functorial-ringed-topoi} の関手性を適用すると、
$\mathcal{O}'_1$-加群の拡大 $\mathcal{F}'_1$ を得る
（この特別な場合の $\mathcal{F}'_1$ は後で記述する）。注意
\ref{defos-remark-trivial-extension-ringed-topoi} により、注意
\ref{defos-remark-short-exact-sequence-thickenings-ringed-topoi} の標準的な分裂
$\pi : (\Sh(\mathcal{D}_1), \mathcal{O}'_1) \to
(\Sh(\mathcal{C}), \mathcal{O})$ を用いて
$\xi_{\mathcal{F}'_1} \in
\Ext^1_\mathcal{O}(\mathcal{F}, \mathcal{K}_1)$.
を得る。最後に、補題 \ref{defos-lemma-inf-obs-ext-ringed-topoi} により障害
$$
o(\mathcal{F}, \mathcal{K}_3, c_3) \in
\Ext^2_\mathcal{O}(\mathcal{F}, \mathcal{K}_3)
$$
を得る。この状況で、短完全列
$0 \to \mathcal{K}_3 \to \mathcal{K}_2 \to \mathcal{K}_1 \to 0$
から生じる標準的な写像
$$
\partial :
\Ext^1_\mathcal{O}(\mathcal{F}, \mathcal{K}_1)
\longrightarrow
\Ext^2_\mathcal{O}(\mathcal{F}, \mathcal{K}_3)
$$
は $\xi_{\mathcal{F}'_1}$ を障害類
$o(\mathcal{F}, \mathcal{K}_3, c_3)$ に写すと {\bf 主張} する。

\medskip\noindent
この主張を証明するため、$\mathcal{K}$ が入射 $\mathcal{O}$-加群となる
埋め込み $j : \mathcal{K}_3 \to \mathcal{K}$ を選ぶ。
$j$ を写像 $j' : \mathcal{K}_2 \to \mathcal{K}$ に持ち上げられる。
$\mathcal{E}'_2 = j'_*\mathcal{F}'_2$ を $j'$ による
$\mathcal{F}'_2$ の押し出しと置く。このとき
$c_{\mathcal{E}'_2} = j' \circ c_2$ である。図式は
$$
\xymatrix{
0 \ar[r] &
\mathcal{K}_2 \ar[r] \ar[d]_{j'} &
\mathcal{F}'_2 \ar[r] \ar[d] &
\mathcal{F} \ar[r] \ar[d] & 0 \\
0 \ar[r] &
\mathcal{K} \ar[r] &
\mathcal{E}'_2 \ar[r] &
\mathcal{F} \ar[r] & 0
}
$$
$\mathcal{E}'_3 = \mathcal{E}'_2$ と置くが、写像
$\mathcal{O}'_3 \to \mathcal{O}'_2$ を通じて $\mathcal{O}'_3$-加群とみなす。
すると $c_{\mathcal{E}'_3} = j \circ c_3$ である。
補題 \ref{defos-lemma-inf-obs-ext-ringed-topoi} の証明では、
$o(\mathcal{F}, \mathcal{K}_3, c_3)$ を $\mathcal{O}$-加群の拡大
$$
0 \to
\mathcal{K}/\mathcal{K}_3 \to
\mathcal{E}'_3/\mathcal{K}_3 \to
\mathcal{F} \to 0
$$
の類の境界として構成する。一方、
$\mathcal{F}'_1 = \mathcal{F}'_2/\mathcal{K}_3$ であるから、
類 $\xi_{\mathcal{F}'_1}$ は拡大
$$
0 \to \mathcal{K}_2/\mathcal{K}_3 \to \mathcal{F}'_2/\mathcal{K}_3
\to \mathcal{F} \to 0
$$
の類である。ここでは標準的な分裂
$\pi : (\Sh(\mathcal{D}_1), \mathcal{O}'_1) \to
(\Sh(\mathcal{C}), \mathcal{O})$ による $\pi^\sharp$ を用いて、
これを $\mathcal{O}$-加群の列とみなしている。したがって、可換図式
$$
\xymatrix{
0 \ar[r] &
\mathcal{K}_2/\mathcal{K}_3 \ar[r] \ar[d] &
\mathcal{F}'_2/\mathcal{K}_3 \ar[r] \ar[d] &
\mathcal{F} \ar[r] \ar[d] & 0 \\
0 \ar[r] &
\mathcal{K}/\mathcal{K}_3 \ar[r] &
\mathcal{E}'_3/\mathcal{K}_3 \ar[r] &
\mathcal{F} \ar[r] & 0
}
$$
があり、これは（上で与えた $\mathcal{O}$-加群構造について）
$\mathcal{O}$-線形であることから、主張が従う。
\end{remark}







\section{環付きトポス上の加群の無限小変形}
\label{defos-section-deformation-modules-ringed-topoi}

\noindent
環付きトポスの一次厚化
$i : (\Sh(\mathcal{C}), \mathcal{O}) \to (\Sh(\mathcal{D}), \mathcal{O}')$
をとる。第\ref{defos-section-thickenings-ringed-topoi}節で導入した記法を
自由に用いる。$\mathcal{F}'$ を $\mathcal{O}'$-加群とし、
$\mathcal{F} = i^*\mathcal{F}'$ と置く。この状況では、
$\mathcal{O}'$-加群の短完全列
$$
0 \to \mathcal{I}\mathcal{F}' \to \mathcal{F}' \to \mathcal{F} \to 0
$$
を得る。$\mathcal{I}^2 = 0$ なので、$\mathcal{I}\mathcal{F}'$ 上の
$\mathcal{O}'$-加群構造は一意な $\mathcal{O}$-加群構造から来る。
したがって上の列は (\ref{defos-equation-extension-ringed-topoi}) の形の拡大である。
特別な場合として $\mathcal{F}' = \mathcal{O}'$ ならば、
$i^*\mathcal{O}' = \mathcal{O}$ かつ
$\mathcal{I}\mathcal{O}' = \mathcal{I}$ であり、構造層の列
$$
0 \to \mathcal{I} \to \mathcal{O}' \to \mathcal{O} \to 0
$$
を回収する。

\begin{lemma}
\label{defos-lemma-inf-map-special-ringed-topoi}
環付きトポスの一次厚化
$i : (\Sh(\mathcal{C}), \mathcal{O}) \to (\Sh(\mathcal{D}), \mathcal{O}')$
をとる。$\mathcal{F}'$、$\mathcal{G}'$ を $\mathcal{O}'$-加群とし、
$\mathcal{F} = i^*\mathcal{F}'$ および $\mathcal{G} = i^*\mathcal{G}'$ と置く。
$\varphi : \mathcal{F} \to \mathcal{G}$ を $\mathcal{O}$-線形写像とする。
$\varphi$ の $\mathcal{O}'$-線形写像
$\varphi' : \mathcal{F}' \to \mathcal{G}'$ への持ち上げの集合は、空でなければ、
$\Hom_\mathcal{O}(\mathcal{F}, \mathcal{I}\mathcal{G}')$
を作用群とする主等質空間である。
\end{lemma}

\begin{proof}
これは補題 \ref{defos-lemma-inf-map-ringed-topoi} の特別な場合であるが、
直接証明も与える。加群の短完全列
$$
0 \to \mathcal{I} \to \mathcal{O}' \to \mathcal{O} \to 0
\quad\text{および}\quad
0 \to \mathcal{I}\mathcal{G}' \to \mathcal{G}' \to \mathcal{G} \to 0
$$
があり、$\mathcal{F}'$ についても同様である。
$\mathcal{I}$ の平方は零なので、$\mathcal{I}$ および
$\mathcal{I}\mathcal{G}'$ 上の $\mathcal{O}'$-加群構造は
一意な $\mathcal{O}$-加群構造から来る。したがって
$$
\Hom_{\mathcal{O}'}(\mathcal{F}', \mathcal{I}\mathcal{G}') =
\Hom_\mathcal{O}(\mathcal{F}, \mathcal{I}\mathcal{G}')
\quad\text{および}\quad
\Hom_{\mathcal{O}'}(\mathcal{F}', \mathcal{G}) =
\Hom_\mathcal{O}(\mathcal{F}, \mathcal{G})
$$
である。よって完全列
$$
0 \to \Hom_{\mathcal{O}'}(\mathcal{F}', \mathcal{I}\mathcal{G}') \to
\Hom_{\mathcal{O}'}(\mathcal{F}', \mathcal{G}') \to
\Hom_{\mathcal{O}'}(\mathcal{F}', \mathcal{G})
$$
から補題が従う。『ホモロジー』、補題
\ref{homology-lemma-check-exactness} を参照せよ。
\end{proof}

\begin{lemma}
\label{defos-lemma-deform-module-ringed-topoi}
$(f, f')$ を状況\ref{defos-situation-morphism-thickenings-ringed-topoi}
のような環付きトポスの一次厚化の射とする。
$\mathcal{F}'$ を $\mathcal{O}'$-加群とし、
$\mathcal{F} = i^*\mathcal{F}'$ と置く。$\mathcal{F}$ は
$\mathcal{O}_\mathcal{B}$ 上平坦であり、$(f, f')$ は厚化の厳密な射
（定義\ref{defos-definition-strict-morphism-thickenings-ringed-topoi}）
であると仮定する。このとき次は同値である。
\begin{enumerate}
\item $\mathcal{F}'$ は $\mathcal{O}_{\mathcal{B}'}$ 上平坦である。
\item 標準的な写像
$f^*\mathcal{J} \otimes_\mathcal{O} \mathcal{F} \to
\mathcal{I}\mathcal{F}'$
は同型である。
\end{enumerate}
さらに、この場合、写像
$$
f^*\mathcal{J} \otimes_\mathcal{O} \mathcal{F} \to
\mathcal{I} \otimes_\mathcal{O} \mathcal{F} \to
\mathcal{I}\mathcal{F}'
$$
はいずれも同型である。
\end{lemma}

\begin{proof}
$(f, f')$ は厚化の厳密な射なので、写像
$f^*\mathcal{J} \to \mathcal{I}$ は全射である。
したがって最後の主張は (2) から従う。

\medskip\noindent
(1) と (2) の同値性を証明する。定義により、
$\mathcal{O}_\mathcal{B}$ 上の平坦性は
$f^{-1}\mathcal{O}_\mathcal{B}$ 上の平坦性を意味する。
$f^{-1}\mathcal{O}_{\mathcal{B}'}$ 上の平坦性についても同様である。
$(f, f')$ の厳密性と仮定 $\mathcal{F} = i^*\mathcal{F}'$ から、
$$
\mathcal{F} = \mathcal{F}'/(f^{-1}\mathcal{J})\mathcal{F}'
$$
が $\mathcal{C}$ 上の層として成り立つ。さらに、
$f^*\mathcal{J} \otimes_\mathcal{O} \mathcal{F} =
f^{-1}\mathcal{J} \otimes_{f^{-1}\mathcal{O}_\mathcal{B}} \mathcal{F}$.
である。したがって (1) と (2) の同値性は『サイト上の加群』、補題
\ref{sites-modules-lemma-flat-over-thickening} から従う。
\end{proof}

\begin{lemma}
\label{defos-lemma-deform-fp-module-ringed-topoi}
$(f, f')$ を状況\ref{defos-situation-morphism-thickenings-ringed-topoi}
のような環付きトポスの一次厚化の射とする。
$\mathcal{F}'$ を $\mathcal{O}'$-加群とし、
$\mathcal{F} = i^*\mathcal{F}'$ と置く。$\mathcal{F}'$ は
$\mathcal{O}_{\mathcal{B}'}$ 上平坦であり、$(f, f')$ は
厚化の厳密な射であると仮定する。このとき次は同値である。
\begin{enumerate}
\item $\mathcal{F}'$ は有限表示の $\mathcal{O}'$-加群である。
\item $\mathcal{F}$ は有限表示の $\mathcal{O}$-加群である。
\end{enumerate}
\end{lemma}

\begin{proof}
(1) $\Rightarrow$ (2) は『サイト上の加群』、補題
\ref{sites-modules-lemma-local-pullback} から従う。
逆に、$\mathcal{F}$ が有限表示であると仮定する。
$\mathcal{C} = \mathcal{C}'$ と仮定してよいし、そう仮定する。
補題 \ref{defos-lemma-deform-module-ringed-topoi} により短完全列
$$
0 \to \mathcal{I} \otimes_{\mathcal{O}_X} \mathcal{F} \to
\mathcal{F}' \to \mathcal{F} \to 0
$$
を得る。$U$ を $\mathcal{C}$ の対象で、$\mathcal{F}|_U$ が表示
$$
\mathcal{O}_U^{\oplus m} \to \mathcal{O}_U^{\oplus n} \to \mathcal{F}|_U \to 0
$$
をもつものとする。$U$ を被覆の各要素で置き換えることで、写像
$\mathcal{O}_U^{\oplus n} \to \mathcal{F}|_U$ が写像
$(\mathcal{O}'_U)^{\oplus n} \to \mathcal{F}'|_U$ に持ち上がると
仮定してよい。誘導された写像
$\mathcal{I}^{\oplus n} \to \mathcal{I} \otimes \mathcal{F}$ は
$\otimes$ の右完全性により全射である。したがって、再び $U$ を
被覆の各要素で置き換えることで、与えられた写像
$\mathcal{O}_U^{\oplus m} \to \mathcal{O}_U^{\oplus n}$ の持ち上げ
$(\mathcal{O}'|_U)^{\oplus m} \to (\mathcal{O}'|_U)^{\oplus n}$
で、
$$
(\mathcal{O}'_U)^{\oplus m} \to (\mathcal{O}'_U)^{\oplus n} \to
\mathcal{F}'|_U \to 0
$$
が複体となるものを見いだせる。もう一度 $\otimes$ の右完全性を用いると、
この複体が完全であることが分かる。
\end{proof}

\begin{lemma}
\label{defos-lemma-inf-map-rel-ringed-topoi}
$(f, f')$ を状況\ref{defos-situation-morphism-thickenings-ringed-topoi}
のような一次厚化の射とする。$\mathcal{F}'$、$\mathcal{G}'$ を
$\mathcal{O}'$-加群とし、$\mathcal{F} = i^*\mathcal{F}'$ および
$\mathcal{G} = i^*\mathcal{G}'$ と置く。
$\varphi : \mathcal{F} \to \mathcal{G}$ を $\mathcal{O}$-線形写像とする。
$\mathcal{G}'$ は $\mathcal{O}_{\mathcal{B}'}$ 上平坦であり、
$(f, f')$ は厚化の厳密な射であると仮定する。
$\varphi$ の $\mathcal{O}'$-線形写像
$\varphi' : \mathcal{F}' \to \mathcal{G}'$ への持ち上げの集合は、
空でなければ、
$$
\Hom_\mathcal{O}(\mathcal{F},
\mathcal{G} \otimes_\mathcal{O} f^*\mathcal{J})
$$
を作用群とする主等質空間である。
\end{lemma}

\begin{proof}
補題 \ref{defos-lemma-inf-map-special-ringed-topoi} と
\ref{defos-lemma-deform-module-ringed-topoi} を組み合わせればよい。
\end{proof}

\begin{lemma}
\label{defos-lemma-inf-obs-map-special-ringed-topoi}
環付きトポスの一次厚化
$i : (\Sh(\mathcal{C}), \mathcal{O}) \to (\Sh(\mathcal{D}), \mathcal{O}')$
をとる。$\mathcal{F}'$、$\mathcal{G}'$ を $\mathcal{O}'$-加群とし、
$\mathcal{F} = i^*\mathcal{F}'$ および $\mathcal{G} = i^*\mathcal{G}'$ と置く。
$\varphi : \mathcal{F} \to \mathcal{G}$ を $\mathcal{O}$-線形写像とする。
元
$$
o(\varphi) \in
\Ext^1_\mathcal{O}(Li^*\mathcal{F}', \mathcal{I}\mathcal{G}')
$$
が存在し、その消滅は $\varphi$ の $\mathcal{O}'$-線形写像
$\varphi' : \mathcal{F}' \to \mathcal{G}'$ への持ち上げが存在するための
必要十分条件である。
\end{lemma}

\begin{proof}
補題 \ref{defos-lemma-inf-map-special-ringed-topoi} の証明から、写像
$$
\Hom_\mathcal{O}(\mathcal{F}, \mathcal{G}) =
\Hom_{\mathcal{O}'}(\mathcal{F}', \mathcal{G}) \longrightarrow
\Ext^1_{\mathcal{O}'}(\mathcal{F}', \mathcal{I}\mathcal{G}')
$$
による $\varphi$ の境界が消滅することが、持ち上げの存在の必要十分条件
であることは明らかである。ここで
$$
\Ext^1_{\mathcal{O}'}(\mathcal{F}', \mathcal{I}\mathcal{G}') =
\Ext^1_\mathcal{O}(Li^*\mathcal{F}', \mathcal{I}\mathcal{G}')
$$
である。これは導来圏上の $i_* = Ri_*$ と $Li^*$ の随伴性
（『サイト上のコホモロジー』、補題
\ref{sites-cohomology-lemma-adjoint}）による。以上で結論を得る。
\end{proof}

\begin{lemma}
\label{defos-lemma-inf-obs-map-rel-ringed-topoi}
$(f, f')$ を状況\ref{defos-situation-morphism-thickenings-ringed-topoi}
のような一次厚化の射とする。$\mathcal{F}'$、$\mathcal{G}'$ を
$\mathcal{O}'$-加群とし、$\mathcal{F} = i^*\mathcal{F}'$ および
$\mathcal{G} = i^*\mathcal{G}'$ と置く。
$\varphi : \mathcal{F} \to \mathcal{G}$ を $\mathcal{O}$-線形写像とする。
$\mathcal{F}'$ と $\mathcal{G}'$ は $\mathcal{O}_{\mathcal{B}'}$ 上平坦であり、
$(f, f')$ は厚化の厳密な射であると仮定する。元
$$
o(\varphi) \in
\Ext^1_\mathcal{O}(\mathcal{F},
\mathcal{G} \otimes_\mathcal{O} f^*\mathcal{J})
$$
が存在し、その消滅は $\varphi$ の $\mathcal{O}'$-線形写像
$\varphi' : \mathcal{F}' \to \mathcal{G}'$ への持ち上げが存在するための
必要十分条件である。
\end{lemma}

\begin{proof}[第一の証明]
これは補題 \ref{defos-lemma-inf-obs-map-special-ringed-topoi} から従う。
実際、補題の仮定の下で
$$
\Ext^1_\mathcal{O}(Li^*\mathcal{F}', \mathcal{I}\mathcal{G}') =
\Ext^1_\mathcal{O}(\mathcal{F},
\mathcal{G} \otimes_\mathcal{O} f^*\mathcal{J})
$$
が成り立つと主張する。実際、補題 \ref{defos-lemma-deform-module-ringed-topoi}
により
$\mathcal{I}\mathcal{G}' =
\mathcal{G} \otimes_\mathcal{O} f^*\mathcal{J}$
である。一方、
$$
H^{-1}(Li^*\mathcal{F}') =
\text{Tor}_1^{\mathcal{O}'}(\mathcal{F}', \mathcal{O})
$$
である（局所計算は省略する）。短完全列
$$
0 \to \mathcal{I} \to \mathcal{O}' \to \mathcal{O} \to 0
$$
を用いると、この $\text{Tor}_1$ は写像
$\mathcal{I} \otimes_\mathcal{O} \mathcal{F} \to \mathcal{I}\mathcal{F}'$
の核として計算される。この写像は補題
\ref{defos-lemma-deform-module-ringed-topoi} の最後の主張により零である。
したがって $\tau_{\geq -1}Li^*\mathcal{F}' = \mathcal{F}$ である。
また、
$$
\Ext^1_\mathcal{O}(Li^*\mathcal{F}',
\mathcal{I}\mathcal{G}') =
\Ext^1_\mathcal{O}(\tau_{\geq -1}Li^*\mathcal{F}',
\mathcal{I}\mathcal{G}')
$$
である。これは『導来圏』、補題
\ref{derived-lemma-negative-vanishing} の双対による。
\end{proof}

\begin{proof}[第二の証明]
補題 \ref{defos-lemma-inf-obs-map-ringed-topoi} を次のように適用できる。
$\mathcal{K} = \mathcal{I} \otimes_\mathcal{O} \mathcal{F}$ および
$\mathcal{L} = \mathcal{I} \otimes_\mathcal{O} \mathcal{G}$
は補題 \ref{defos-lemma-deform-module-ringed-topoi} により成り立ち、
$c_{\mathcal{F}'} = 1 \otimes 1$ および $c_{\mathcal{G}'} = 1 \otimes 1$
である。また $\psi = 1 \otimes \varphi$ とすれば補題の図式は可換である。
したがって $o(\varphi) = o(\varphi, 1 \otimes \varphi)$ が求める元である。
\end{proof}

\begin{lemma}
\label{defos-lemma-inf-ext-rel-ringed-topoi}
$(f, f')$ を状況\ref{defos-situation-morphism-thickenings-ringed-topoi}
のような一次厚化の射とする。$\mathcal{F}$ を $\mathcal{O}$-加群とする。
$(f, f')$ は厚化の厳密な射であり、$\mathcal{F}$ は
$\mathcal{O}_\mathcal{B}$ 上平坦であると仮定する。
$\mathcal{O}_{\mathcal{B}'}$ 上平坦な $\mathcal{O}'$-加群 $\mathcal{F}'$ と同型
$\alpha : i^*\mathcal{F}' \to \mathcal{F}$ からなる組
$(\mathcal{F}', \alpha)$ が一つ存在するならば、そのような組の同型類の集合は
$\Ext^1_\mathcal{O}(
\mathcal{F}, \mathcal{I} \otimes_\mathcal{O} \mathcal{F})$
を作用群とする主等質空間である。
\end{lemma}

\begin{proof}
そのような加群が一つ存在すると仮定すると、標準的な写像
$$
f^*\mathcal{J} \otimes_\mathcal{O} \mathcal{F} \to
\mathcal{I} \otimes_\mathcal{O} \mathcal{F}
$$
は補題 \ref{defos-lemma-deform-module-ringed-topoi} により同型である。
補題 \ref{defos-lemma-inf-ext-ringed-topoi} を
$\mathcal{K} = \mathcal{I} \otimes_\mathcal{O} \mathcal{F}$ および
$c = 1$ として適用する。補題 \ref{defos-lemma-deform-module-ringed-topoi}
により、対応する拡大 $\mathcal{F}'$ はすべて
$\mathcal{O}_{\mathcal{B}'}$ 上平坦である。
\end{proof}

\begin{lemma}
\label{defos-lemma-inf-obs-ext-rel-ringed-topoi}
$(f, f')$ を状況\ref{defos-situation-morphism-thickenings-ringed-topoi}
のような一次厚化の射とする。$\mathcal{F}$ を $\mathcal{O}$-加群とする。
$(f, f')$ は厚化の厳密な射であり、$\mathcal{F}$ は
$\mathcal{O}_\mathcal{B}$ 上平坦であると仮定する。
$i^*\mathcal{F}' \cong \mathcal{F}$ を満たす
$\mathcal{O}_{\mathcal{B}'}$ 上平坦な $\mathcal{O}'$-加群
$\mathcal{F}'$ が存在するための必要十分条件は、次の二条件である。
\begin{enumerate}
\item 標準的な写像
$f^*\mathcal{J} \otimes_\mathcal{O} \mathcal{F} \to
\mathcal{I} \otimes_\mathcal{O} \mathcal{F}$
が同型である。
\item 補題 \ref{defos-lemma-inf-obs-ext-ringed-topoi} の類
$o(\mathcal{F}, \mathcal{I} \otimes_\mathcal{O} \mathcal{F}, 1)
\in \Ext^2_\mathcal{O}(
\mathcal{F}, \mathcal{I} \otimes_\mathcal{O} \mathcal{F})$
が零である。
\end{enumerate}
\end{lemma}

\begin{proof}
補題 \ref{defos-lemma-deform-module-ringed-topoi} による
$\mathcal{O}_{\mathcal{B}'}$ 上平坦な $\mathcal{O}'$-加群の特徴づけと、
補題 \ref{defos-lemma-inf-obs-ext-ringed-topoi} から直ちに従う。
\end{proof}






\section{環付きトポスの平坦な厚化上の平坦加群への応用}
\label{defos-section-flat-ringed-topoi}

\noindent
環付きトポスの可換図式
$$
\xymatrix{
(\Sh(\mathcal{C}), \mathcal{O}) \ar[r]_i \ar[d]_f &
(\Sh(\mathcal{D}), \mathcal{O}') \ar[d]^{f'} \\
(\Sh(\mathcal{B}), \mathcal{O}_\mathcal{B}) \ar[r]^t &
(\Sh(\mathcal{B}'), \mathcal{O}_{\mathcal{B}'})
}
$$
で、その横向きの矢印が状況
\ref{defos-situation-morphism-thickenings-ringed-topoi} のような一次厚化であるものを考える。
$\mathcal{I} = \Ker(i^\sharp) \subset \mathcal{O}'$ および
$\mathcal{J} = \Ker(t^\sharp) \subset \mathcal{O}_{\mathcal{B}'}$ と置く。
$\mathcal{F}$ を $\mathcal{O}$-加群とする。次を仮定する。
\begin{enumerate}
\item $(f, f')$ は厚化の厳密な射である。
\item $f'$ は平坦である。
\item $\mathcal{F}$ は $\mathcal{O}_\mathcal{B}$ 上平坦である。
\end{enumerate}
(1) $+$ (2) から $\mathcal{I} = f^*\mathcal{J}$ が従うことに注意せよ
（補題 \ref{defos-lemma-deform-module-ringed-topoi} を $\mathcal{O}'$ に適用する）。
これらの仮定の下では、前節の理論は特に簡明になる。
すでに得られた結果を次の補題にまとめる。

\begin{lemma}
\label{defos-lemma-flat-ringed-topoi}
上の状況において、次が成り立つ。
\begin{enumerate}
\item $i^*\mathcal{F}' \cong \mathcal{F}$ を満たす
$\mathcal{O}_{\mathcal{B}'}$ 上平坦な $\mathcal{O}'$-加群 $\mathcal{F}'$
が存在するための必要十分条件は、補題
\ref{defos-lemma-inf-obs-ext-ringed-topoi} の類
$o(\mathcal{F}, f^*\mathcal{J} \otimes_\mathcal{O} \mathcal{F}, 1)
\in \Ext^2_\mathcal{O}(
\mathcal{F}, f^*\mathcal{J} \otimes_\mathcal{O} \mathcal{F})$
が零であることである。
\item そのような加群が存在するならば、持ち上げの同型類の集合は
$\Ext^1_\mathcal{O}(
\mathcal{F}, f^*\mathcal{J} \otimes_\mathcal{O} \mathcal{F})$
を作用群とする主等質空間である。
\item 持ち上げ $\mathcal{F}'$ が与えられたとき、引き戻しが
$\text{id}_\mathcal{F}$ となる $\mathcal{F}'$ の自己同型の集合は、標準的に
$\Ext^0_\mathcal{O}(
\mathcal{F}, f^*\mathcal{J} \otimes_\mathcal{O} \mathcal{F})$
と同型である。
\end{enumerate}
\end{lemma}

\begin{proof}
(1) は、上で見た $\mathcal{I} = f^*\mathcal{J}$ と補題
\ref{defos-lemma-inf-obs-ext-rel-ringed-topoi} から従う。
(2) は補題 \ref{defos-lemma-inf-ext-rel-ringed-topoi} から従う。
(3) は補題 \ref{defos-lemma-inf-map-rel-ringed-topoi} から従う。
\end{proof}

\begin{situation}
\label{defos-situation-morphism-flat-thickenings-ringed-topoi}
$f : (\Sh(\mathcal{C}), \mathcal{O}) \to
(\Sh(\mathcal{B}), \mathcal{O}_\mathcal{B})$ を環付きトポスの射とする。
可換図式
$$
\xymatrix{
(\Sh(\mathcal{C}'_1), \mathcal{O}'_1) \ar[r]_h \ar[d]_{f'_1} &
(\Sh(\mathcal{C}'_2), \mathcal{O}'_2) \ar[d]_{f'_2} \\
(\Sh(\mathcal{B}'_1), \mathcal{O}_{\mathcal{B}'_1}) \ar[r] &
(\Sh(\mathcal{B}'_2), \mathcal{O}_{\mathcal{B}'_2})
}
$$
を考える。ここで $h$ は $(\Sh(\mathcal{C}), \mathcal{O})$ の
一次厚化の射であり、下側の横向きの矢印は
$(\Sh(\mathcal{B}), \mathcal{O}_\mathcal{B})$ の一次厚化の射であり、
各 $f'_i$ の制限は $f$ であり、二つの組 $(f, f_i')$ はいずれも
厚化の厳密な射であり、二つの $f'_i$ はいずれも平坦である。
最後に、$\mathcal{F}$ を $\mathcal{O}_\mathcal{B}$ 上平坦な
$\mathcal{O}$-加群とする。
\end{situation}

\begin{lemma}
\label{defos-lemma-functorial-ringed-topoi}
状況\ref{defos-situation-morphism-flat-thickenings-ringed-topoi}において、障害類
$o(\mathcal{F}, f^*\mathcal{J}_2 \otimes_\mathcal{O} \mathcal{F}, 1)$
は障害類
$o(\mathcal{F}, f^*\mathcal{J}_1 \otimes_\mathcal{O} \mathcal{F}, 1)$
に、標準的な写像
$$
\Ext^2_\mathcal{O}(
\mathcal{F}, f^*\mathcal{J}_2 \otimes_\mathcal{O} \mathcal{F})
\to \Ext^2_\mathcal{O}(
\mathcal{F}, f^*\mathcal{J}_1 \otimes_\mathcal{O} \mathcal{F})
$$
の下で写る。
\end{lemma}

\begin{proof}
注意 \ref{defos-remark-obstruction-extension-functorial-ringed-topoi} から従う。
\end{proof}

\begin{situation}
\label{defos-situation-ses-flat-thickenings-ringed-topoi}
$f : (\Sh(\mathcal{C}), \mathcal{O}) \to
(\Sh(\mathcal{B}), \mathcal{O}_\mathcal{B})$ を環付きトポスの射とする。
可換図式
$$
\xymatrix{
(\Sh(\mathcal{C}'_1), \mathcal{O}'_1) \ar[r]_h \ar[d]_{f'_1} &
(\Sh(\mathcal{C}'_2), \mathcal{O}'_2) \ar[r] \ar[d]_{f'_2} &
(\Sh(\mathcal{C}'_3), \mathcal{O}'_3) \ar[d]_{f'_3} \\
(\Sh(\mathcal{B}'_1), \mathcal{O}_{\mathcal{B}'_1}) \ar[r] &
(\Sh(\mathcal{B}'_2), \mathcal{O}_{\mathcal{B}'_2}) \ar[r] &
(\Sh(\mathcal{B}'_3), \mathcal{O}_{\mathcal{B}'_3})
}
$$
を考える。ここで、(a) 上段は $(\Sh(\mathcal{C}), \mathcal{O})$ の
一次厚化の短完全列であり、(b) 下段は
$(\Sh(\mathcal{B}), \mathcal{O}_\mathcal{B})$ の一次厚化の短完全列であり、
(c) 各 $f'_i$ の制限は $f$ であり、(d) 各組 $(f, f_i')$ は
厚化の厳密な射であり、(e) 各 $f'_i$ は平坦である。最後に、
$\mathcal{F}'_2$ を $\mathcal{O}_{\mathcal{B}'_2}$ 上平坦な
$\mathcal{O}'_2$-加群とし、$\mathcal{F} = \mathcal{F}'_2 \otimes \mathcal{O}$
と置く。注意 \ref{defos-remark-short-exact-sequence-thickenings-ringed-topoi}
の標準的な分裂を
$\pi : (\Sh(\mathcal{C}'_1), \mathcal{O}'_1) \to
(\Sh(\mathcal{C}), \mathcal{O})$
と書く。
\end{situation}

\begin{lemma}
\label{defos-lemma-verify-iv-ringed-topoi}
状況\ref{defos-situation-ses-flat-thickenings-ringed-topoi}において、
$\pi^*\mathcal{F}$ と $h^*\mathcal{F}'_2$ は
$\mathcal{O}_{\mathcal{B}'_1}$ 上平坦な $\mathcal{O}'_1$-加群であり、
$(\Sh(\mathcal{C}), \mathcal{O})$ への制限は $\mathcal{F}$ である。
両者の差（補題 \ref{defos-lemma-flat-ringed-topoi}）は
$\Ext^1_\mathcal{O}(\mathcal{F},
f^*\mathcal{J}_1 \otimes_\mathcal{O} \mathcal{F})$
の元 $\theta$ であり、その
$\Ext^2_\mathcal{O}(\mathcal{F},
f^*\mathcal{J}_3 \otimes_\mathcal{O} \mathcal{F})$
における境界は、$\mathcal{F}$ を $\mathcal{O}_{\mathcal{B}'_3}$ 上平坦な
$\mathcal{O}'_3$-加群へ持ち上げることに対する障害
（補題 \ref{defos-lemma-flat-ringed-topoi}）に等しい。
\end{lemma}

\begin{proof}
$\pi^*\mathcal{F}$ と $h^*\mathcal{F}'_2$ はいずれも
$(\Sh(\mathcal{C}), \mathcal{O})$ 上で $\mathcal{F}$ に制限され、
$\pi^*\mathcal{F} \to \mathcal{F}$ と
$h^*\mathcal{F}'_2 \to \mathcal{F}$ の核はいずれも
$f^*\mathcal{J}_1 \otimes_\mathcal{O} \mathcal{F}$ で与えられる。
したがって補題 \ref{defos-lemma-deform-module-ringed-topoi} により平坦である。
加群の列
$$
0 \to f^*\mathcal{J}_3 \otimes_\mathcal{O} \mathcal{F} \to
f^*\mathcal{J}_2 \otimes_\mathcal{O} \mathcal{F} \to
f^*\mathcal{J}_1 \otimes_\mathcal{O} \mathcal{F} \to 0
$$
は、状況\ref{defos-situation-ses-flat-thickenings-ringed-topoi}の仮定と
$\mathcal{F}$ が $\mathcal{O}_\mathcal{B}$ 上平坦であることにより短完全である。
したがって境界を取ることができる。障害類についての主張は、注意
\ref{defos-remark-complex-thickenings-and-ses-modules-ringed-topoi} の結果を
この特別な状況に直接言い換えたものである。
\end{proof}







\section{環付きトポスの変形と素朴な余接複体}
\label{defos-section-deformations-ringed-topoi}

\noindent
本節では、素朴な余接複体を用いて変形理論を少し展開する。まず、
環付きトポスの一次厚化
$t : (\Sh(\mathcal{B}), \mathcal{O}_\mathcal{B}) \to
(\Sh(\mathcal{B}'), \mathcal{O}_{\mathcal{B}'})$ から始める。
$\mathcal{J} = \Ker(t^\sharp)$ と記し、$\mathcal{B}$ と
$\mathcal{B}'$ の台となるトポスを同一視する。さらに、環付きトポスの射
$f : (\Sh(\mathcal{C}), \mathcal{O}) \to
(\Sh(\mathcal{B}), \mathcal{O}_\mathcal{B})$、$\mathcal{O}$-加群
$\mathcal{G}$、および $f^{-1}\mathcal{O}_\mathcal{B}$-加群の層の写像
$f^{-1}\mathcal{J} \to \mathcal{G}$ が与えられていると仮定する。
本節では、次の図式の疑問符に入るものを見いだせるかを問う。
\begin{equation}
\label{defos-equation-to-solve-ringed-topoi}
\vcenter{
\xymatrix{
0 \ar[r] & \mathcal{G} \ar[r] & {?} \ar[r] & \mathcal{O} \ar[r] & 0 \\
0 \ar[r] & f^{-1}\mathcal{J} \ar[u]^c \ar[r] &
f^{-1}\mathcal{O}_{\mathcal{B}'} \ar[u] \ar[r] &
f^{-1}\mathcal{O}_\mathcal{B} \ar[u] \ar[r] & 0
}
}
\end{equation}
また、解が存在する場合にそれがどの程度一意であるかも問う。
より正確には、一次厚化
$i : (\Sh(\mathcal{C}), \mathcal{O}) \to (\Sh(\mathcal{C}'), \mathcal{O}')$
と、(\ref{defos-equation-morphism-thickenings-ringed-topoi}) のような厚化の射
$(f, f')$ であって、$\Ker(i^\sharp)$ が $\mathcal{G}$ と同一視され、
$(f')^\sharp$ が与えられた写像 $c$ を誘導するものを求める。
$(\Sh(\mathcal{C}'), \mathcal{O}')$ を
(\ref{defos-equation-to-solve-ringed-topoi}) の {\it 解} と呼ぶ。


\begin{lemma}
\label{defos-lemma-huge-diagram-ringed-topoi}
環付きトポスの射の可換図式
\begin{equation}
\label{defos-equation-huge-1-ringed-topoi}
\vcenter{
\xymatrix{
& (\Sh(\mathcal{C}_2), \mathcal{O}_2) \ar[r]_{i_2} \ar[d]_{f_2} \ar[ddl]_g &
(\Sh(\mathcal{C}'_2), \mathcal{O}'_2) \ar[d]^{f'_2} \\
&
(\Sh(\mathcal{B}_2), \mathcal{O}_{\mathcal{B}_2}) \ar[r]^{t_2} \ar[ddl]|\hole &
(\Sh(\mathcal{B}'_2), \mathcal{O}_{\mathcal{B}'_2}) \ar[ddl] \\
(\Sh(\mathcal{C}_1), \mathcal{O}_1) \ar[r]_{i_1} \ar[d]_{f_1} &
(\Sh(\mathcal{C}'_1), \mathcal{O}'_1) \ar[d]^{f'_1} \\
(\Sh(\mathcal{B}_1), \mathcal{O}_{\mathcal{B}_1}) \ar[r]^{t_1} &
(\Sh(\mathcal{B}'_1), \mathcal{O}_{\mathcal{B}'_1})
}
}
\end{equation}
が与えられ、その横向きの矢印は一次厚化であると仮定する。
$\mathcal{G}_j = \Ker(i_j^\sharp)$ と置き、可換図式
\begin{equation}
\label{defos-equation-huge-2-ringed-topoi}
\vcenter{
\xymatrix{
& 0 \ar[r] & \mathcal{G}_2 \ar[r] &
\mathcal{O}'_2 \ar[r] &
\mathcal{O}_2 \ar[r] & 0 \\
& 0 \ar[r]|\hole &
f_2^{-1}\mathcal{J}_2 \ar[u]_{c_2} \ar[r] &
f_2^{-1}\mathcal{O}_{\mathcal{B}'_2} \ar[u] \ar[r]|\hole &
f_2^{-1}\mathcal{O}_{\mathcal{B}_2} \ar[u] \ar[r] & 0 \\
0 \ar[r] &
\mathcal{G}_1 \ar[ruu] \ar[r] &
\mathcal{O}'_1 \ar[r] &
\mathcal{O}_1 \ar[ruu] \ar[r] & 0 \\
0 \ar[r] &
f_1^{-1}\mathcal{J}_1 \ar[ruu]|\hole \ar[u]^{c_1} \ar[r] &
f_1^{-1}\mathcal{O}_{\mathcal{B}'_1} \ar[ruu]|\hole \ar[u] \ar[r] &
f_1^{-1}\mathcal{O}_{\mathcal{B}_1} \ar[ruu]|\hole \ar[u] \ar[r] & 0
}
}
\end{equation}
を誘導する $g^{-1}\mathcal{O}_1$-加群の写像
$\nu : g^{-1}\mathcal{G}_1 \to \mathcal{G}_2$ が与えられていると仮定する。
その手前と奥は (\ref{defos-equation-to-solve-ringed-topoi}) の解である。
（北北西向きの矢印は、始域に $g^{-1}$ を適用した後の
$\mathcal{C}_2$ 上の写像である。）
\begin{enumerate}
\item
$\Ext^1_{\mathcal{O}_2}(
Lg^*\NL_{\mathcal{O}_1/\mathcal{O}_{\mathcal{B}_1}}, \mathcal{G}_2)$
には正準元が存在し、その消滅は、$\nu$ と両立しながら
(\ref{defos-equation-huge-1-ringed-topoi}) に収まる環付きトポスの射
$(\Sh(\mathcal{C}'_2), \mathcal{O}'_2) \to
(\Sh(\mathcal{C}'_1), \mathcal{O}'_1)$ が存在するための必要十分条件である。
\item $\nu$ と両立しながら
(\ref{defos-equation-huge-1-ringed-topoi}) に収まる環付きトポスの射
$(\Sh(\mathcal{C}'_2), \mathcal{O}'_2) \to
(\Sh(\mathcal{C}'_1), \mathcal{O}'_1)$
が存在するならば、そのような射全体の集合は
$$
\Hom_{\mathcal{O}_1}(
\Omega_{\mathcal{O}_1/\mathcal{O}_{\mathcal{B}_1}}, g_*\mathcal{G}_2) =
\Hom_{\mathcal{O}_2}(
g^*\Omega_{\mathcal{O}_1/\mathcal{O}_{\mathcal{B}_1}}, \mathcal{G}_2) =
\Ext^0_{\mathcal{O}_2}(
Lg^*\NL_{\mathcal{O}_1/\mathcal{O}_{\mathcal{B}_1}}, \mathcal{G}_2).
$$
を作用群とする主等質空間である。
\end{enumerate}
\end{lemma}

\begin{proof}
この補題の証明は補題 \ref{defos-lemma-huge-diagram-ringed-spaces} の証明と
同一である。読者には以下ではなくそちらの証明を読むことを勧める。
以下に現れるすべての厚化について、台となるトポスを同一視する
（主張ですでにこの規約を用いた）。補題の最後の主張の等号は
定義から直ちに従う。したがって証明の残りでは群
$\Ext^k_{\mathcal{O}_2}(
Lg^*\NL_{\mathcal{O}_1/\mathcal{O}_{\mathcal{B}_1}}, \mathcal{G}_2)$,
$k = 0, 1$ を用いる。まず、補題に現れるすべての環付きトポスの
台となるトポスが同じ場合に帰着できることを示す。

\medskip\noindent
そのため、$g^{-1}\NL_{\mathcal{O}_1/\mathcal{O}_{\mathcal{B}_1}}$ は
環の層の準同型
$g^{-1}f_1^{-1}\mathcal{O}_{\mathcal{B}_1} \to g^{-1}\mathcal{O}_1$
の素朴な余接複体に等しいことに注意する。『サイト上の加群』、補題
\ref{sites-modules-lemma-pullback-differentials} を参照せよ。さらに、
$\NL_{\mathcal{O}_1/\mathcal{O}_{\mathcal{B}_1}}$ の次数 $0$ の項は
平坦 $\mathcal{O}_1$-加群なので、標準的な写像
$$
Lg^*\NL_{\mathcal{O}_1/\mathcal{O}_{\mathcal{B}_1}}
\longrightarrow
g^{-1}\NL_{\mathcal{O}_1/\mathcal{O}_{\mathcal{B}_1}}
\otimes_{g^{-1}\mathcal{O}_1} \mathcal{O}_2
$$
は次数 $0$ および $-1$ のコホモロジー層上で同型を誘導する。
したがって補題の Ext 群を
$$
\Ext^k_{g^{-1}\mathcal{O}_1}(
g^{-1}\NL_{\mathcal{O}_1/\mathcal{O}_{\mathcal{B}_1}}, \mathcal{G}_2) =
\Ext^k_{g^{-1}\mathcal{O}_1}(
\NL_{g^{-1}\mathcal{O}_1/g^{-1}f_1^{-1}\mathcal{O}_{\mathcal{B}_1}},
\mathcal{G}_2)
$$
で置き換えられる。$\nu$ と両立しながら
(\ref{defos-equation-huge-1-ringed-topoi}) に収まる環付きトポスの射
$(\Sh(\mathcal{C}'_2), \mathcal{O}'_2) \to
(\Sh(\mathcal{C}'_1), \mathcal{O}'_1)$ 全体の集合は、$f^\sharp$ および
$\nu$ と両立する
$g^{-1}f_1^{-1}\mathcal{O}_{\mathcal{B}'_1}$-代数の準同型
$g^{-1}\mathcal{O}'_1 \to \mathcal{O}'_2$ 全体の集合と全単射する。
このようにして、サイト $\mathcal{C}$ 上の層の図式
(\ref{defos-equation-huge-2-ringed-topoi}) があり（台となるトポス上で
$f_1 = f_2 = \text{id}$）、それに収まる環の層の準同型
$\mathcal{O}'_1 \to \mathcal{O}'_2$ を求めるものと仮定できる。

\medskip\noindent
補題の証明の残りでは、台となる位相空間はすべて同じであると仮定する。
すなわち、サイト $\mathcal{C}$ 上の層の図式
(\ref{defos-equation-huge-2-ringed-topoi}) があり（台となるトポス上で
$f_1 = f_2 = \text{id}$）、それに収まる環の層の準同型
$\mathcal{O}'_1 \to \mathcal{O}'_2$ を求める。Ext 群としては
$\Ext^k_{\mathcal{O}_1}(
\NL_{\mathcal{O}_1/\mathcal{O}_{\mathcal{B}_1}}, \mathcal{G}_2)$、$k = 0, 1$ を用いる。

\medskip\noindent
第1段階。障害類の構成。集合の層
$$
\mathcal{E} = \mathcal{O}'_1 \times_{\mathcal{O}_2} \mathcal{O}'_2
$$
を考える。これには全射 $\alpha : \mathcal{E} \to \mathcal{O}_1$ が
付随するので、$\NL_{\mathcal{O}_1/\mathcal{O}_{\mathcal{B}_1}}$
の代わりに $\NL(\alpha)$ を用いることができる。『サイト上の加群』、
補題 \ref{sites-modules-lemma-NL-up-to-qis} を参照せよ。次のように置く。
$$
\mathcal{I}' =
\Ker(\mathcal{O}_{\mathcal{B}'_1}[\mathcal{E}] \to \mathcal{O}_1)
\quad\text{および}\quad
\mathcal{I} =
\Ker(\mathcal{O}_{\mathcal{B}_1}[\mathcal{E}] \to \mathcal{O}_1)
$$
全射 $\mathcal{I}' \to \mathcal{I}$ があり、その核は
$\mathcal{J}_1\mathcal{O}_{\mathcal{B}'_1}[\mathcal{E}]$ である。
二つの $\mathcal{O}_{\mathcal{B}'_2}$-代数準同型
$$
a : \mathcal{O}_{\mathcal{B}'_1}[\mathcal{E}] \to \mathcal{O}'_1
\quad\text{および}\quad
b : \mathcal{O}_{\mathcal{B}'_1}[\mathcal{E}] \to \mathcal{O}'_2
$$
を得る。これらは写像
$a|_{\mathcal{I}'} : \mathcal{I}' \to \mathcal{G}_1$ および
$b|_{\mathcal{I}'} : \mathcal{I}' \to \mathcal{G}_2$ を誘導する。
$a$ と $b$ はともに $(\mathcal{I}')^2$ を零化する。さらに、
(\ref{defos-equation-huge-2-ringed-topoi}) の左側の正方形は可換なので、
$a$ と $b$ は $\mathcal{G}_2$ への写像として
$\mathcal{J}_1\mathcal{O}_{\mathcal{B}'_1}[\mathcal{E}]$ 上で一致する。
したがって差
$b|_{\mathcal{I}'} - \nu \circ a|_{\mathcal{I}'}$
は良定義な $\mathcal{O}_1$-線形写像
$$
\xi : \mathcal{I}/\mathcal{I}^2 \longrightarrow \mathcal{G}_2
$$
を誘導する。この写像は、$\mathcal{I}$ の局所切断 $f$ の類を
$a(f') - \nu(b(f'))$ に写す。ここで $f'$ は $f$ を
$\mathcal{I}'$ の局所切断に持ち上げたものである。その像（下記参照）を
$[\xi] \in \Ext^1_{\mathcal{O}_1}(\NL(\alpha), \mathcal{G}_2)$
と記す。

\medskip\noindent
第2段階。$[\xi]$ の消滅が必要であること。次のように記す。
$\Omega =
\Omega_{\mathcal{O}_{\mathcal{B}_1}[\mathcal{E}]/\mathcal{O}_{\mathcal{B}_1}}
\otimes_{\mathcal{O}_{\mathcal{B}_1}[\mathcal{E}]} \mathcal{O}_1$.
$\NL(\alpha) = (\mathcal{I}/\mathcal{I}^2 \to \Omega)$ は完全三角形
$$
\Omega[0] \to
\NL(\alpha) \to
\mathcal{I}/\mathcal{I}^2[1] \to
\Omega[1]
$$
に収まる。したがって $[\xi]$ が零であるための必要十分条件は、ある写像
$\Omega \to \mathcal{G}_2$ に対して $\xi$ が合成
$\mathcal{I}/\mathcal{I}^2 \to \Omega \to \mathcal{G}_2$
となることである。(\ref{defos-equation-huge-2-ringed-topoi}) に収まる
環の層の準同型
$\varphi : \mathcal{O}'_1 \to \mathcal{O}'_2$
が存在すると仮定する。この場合、写像
$\mathcal{O}'_1[\mathcal{E}] \to \mathcal{G}_2$,
$f' \mapsto b(f') - \varphi(a(f'))$ を考える。計算により、これは
$\mathcal{J}_1\mathcal{O}_{\mathcal{B}'_1}[\mathcal{E}]$ を零化し、導分
$\mathcal{O}_{\mathcal{B}_1}[\mathcal{E}] \to \mathcal{G}_2$ を誘導する。
こうして得られる線形写像 $\Omega \to \mathcal{G}_2$ により、
この場合 $[\xi] = 0$ であることが分かる。

\medskip\noindent
第3段階。$[\xi]$ の消滅が十分であること。
$\theta : \Omega \to \mathcal{G}_2$ を、$\xi$ が
$\theta \circ (\mathcal{I}/\mathcal{I}^2 \to \Omega)$ に等しくなるような
$\mathcal{O}_1$-線形写像とする。計算により、
$$
b + \theta \circ d :
\mathcal{O}_{\mathcal{B}'_1}[\mathcal{E}]
\longrightarrow
\mathcal{O}'_2
$$
は $\mathcal{I}'$ を零化し、したがって
(\ref{defos-equation-huge-2-ringed-topoi}) に収まる写像
$\mathcal{O}'_1 \to \mathcal{O}'_2$ を定める。

\medskip\noindent
上の特別な場合における (2) の証明。省略する。ヒント：補題
\ref{defos-lemma-huge-diagram} の (2) の証明とまったく同じである。
\end{proof}

\begin{lemma}
\label{defos-lemma-NL-represent-ext-class-ringed-topoi}
$\mathcal{C}$ をサイトとし、$\mathcal{A} \to \mathcal{B}$ を
$\mathcal{C}$ 上の環の層の準同型、$\mathcal{G}$ を
$\mathcal{B}$-加群とする。
$\xi \in \Ext^1_\mathcal{B}(\NL_{\mathcal{B}/\mathcal{A}}, \mathcal{G})$
とする。このとき集合の層の写像
$\alpha : \mathcal{E} \to \mathcal{B}$ であって、
$\xi \in \Ext^1_\mathcal{B}(\NL(\alpha), \mathcal{G})$ が写像
$\mathcal{I}/\mathcal{I}^2 \to \mathcal{G}$ の類となるものが存在する
（記法については証明を参照せよ）。
\end{lemma}

\begin{proof}
全射 $\mathcal{A}[\mathcal{E}] \to \mathcal{B}$ の核が
$\mathcal{I}$ となるような $\alpha : \mathcal{E} \to \mathcal{B}$
が与えられたとき、複体
$\NL(\alpha) = (\mathcal{I}/\mathcal{I}^2 \to 
\Omega_{\mathcal{A}[\mathcal{E}]/\mathcal{A}}
\otimes_{\mathcal{A}[\mathcal{E}]} \mathcal{B})$ は標準的に
$\NL_{\mathcal{B}/\mathcal{A}}$ と同型であることを思い出そう。
『サイト上の加群』、補題 \ref{sites-modules-lemma-NL-up-to-qis}
を参照せよ。さらに、
$\Omega = \Omega_{\mathcal{A}[\mathcal{E}]/\mathcal{A}}
\otimes_{\mathcal{A}[\mathcal{E}]} \mathcal{B}$ は前層
$U \mapsto \bigoplus_{e \in \mathcal{E}(U)} \mathcal{B}(U)$.
に付随する層である。言い換えれば、$\Omega$ は集合の層
$\mathcal{E}$ 上の自由 $\mathcal{B}$-加群であり、特に標準的な写像
$\mathcal{E} \to \Omega$ がある。

\medskip\noindent
以上を踏まえ、ある $\mathcal{E}$ を選ぶ（例えば、素朴な余接複体の
定義と同様に $\mathcal{E} = \mathcal{B}$ とする）。$\xi$ を写像
$\mathcal{I}/\mathcal{I}^2 \to \mathcal{G}$ の類として表すことへの障害は
$\Ext^1_\mathcal{B}(\Omega, \mathcal{G})$ の元である。これが
$\mathcal{B}$-加群の拡大
$0 \to \mathcal{G} \to \mathcal{H} \to \Omega \to 0$
で表されるとする。集合の層
$\mathcal{E}' = \mathcal{E} \times_\Omega \mathcal{H}$
を考える。これには誘導された写像
$\alpha' : \mathcal{E}' \to \mathcal{B}$ が付随する。
$\mathcal{I}' = \Ker(\mathcal{A}[\mathcal{E}'] \to \mathcal{B})$ および
$\Omega' = \Omega_{\mathcal{A}[\mathcal{E}']/\mathcal{A}}
\otimes_{\mathcal{A}[\mathcal{E}']} \mathcal{B}$.
と置く。擬同型 $\NL(\alpha') \to \NL(\alpha)$ による
$\xi$ の引き戻しは
$\Ext^1_\mathcal{B}(\Omega', \mathcal{G})$
において零に写る。実際、$\Omega'$ は集合の層 $\mathcal{E}'$ 上の
自由 $\mathcal{B}$-加群であり、構成により可換図式
$$
\xymatrix{
\mathcal{E}' \ar[r] \ar[d] & \mathcal{E} \ar[d] \\
\mathcal{H} \ar[r] & \Omega
}
$$
があるので、写像 $\Omega' \to \Omega$ による拡大
$\mathcal{H}$ の引き戻しは分裂する。これで証明が完了する。
\end{proof}

\begin{lemma}
\label{defos-lemma-choices-ringed-topoi}
\leavevmode\par\noindent
(\ref{defos-equation-to-solve-ringed-topoi}) の解が存在するならば、
解の同型類の集合は
$\Ext^1_\mathcal{O}(
\NL_{\mathcal{O}/\mathcal{O}_\mathcal{B}}, \mathcal{G})$
を作用群とする主等質空間である。
\end{lemma}

\begin{proof}
まず、(\ref{defos-equation-to-solve-ringed-topoi}) の二つの解
$\mathcal{O}'_1$ と $\mathcal{O}'_2$ が与えられると、補題
\ref{defos-lemma-huge-diagram-ringed-topoi} により、写像
$\mathcal{O}'_1 \to \mathcal{O}'_2$ の存在に対する障害元
$o(\mathcal{O}'_1, \mathcal{O}'_2) \in \Ext^1_\mathcal{O}(
\NL_{\mathcal{O}/\mathcal{O}_\mathcal{B}}, \mathcal{G})$
を得る。この元は明らかに同型の存在に対する障害であり、したがって
同型類を分離する。よって証明を完了するには、解 $\mathcal{O}'$ と元
$\xi \in \Ext^1_\mathcal{O}(
\NL_{\mathcal{O}/\mathcal{O}_\mathcal{B}}, \mathcal{G})$
が与えられたとき、$o(\mathcal{O}', \mathcal{O}'_\xi) = \xi$ を満たす
第二の解 $\mathcal{O}'_\xi$ を見いだせることを示せば十分である。

\medskip\noindent
類 $\xi$ に対し、補題
\ref{defos-lemma-NL-represent-ext-class-ringed-topoi} のような
$\alpha : \mathcal{E} \to \mathcal{O}$ を選ぶ。核を $\mathcal{I}$ とする全射
$f^{-1}\mathcal{O}_\mathcal{B}[\mathcal{E}] \to \mathcal{O}$
と、それに対応する素朴な余接複体
$\NL(\alpha) = (\mathcal{I}/\mathcal{I}^2 \to
\Omega_{f^{-1}\mathcal{O}_\mathcal{B}[\mathcal{E}]/
f^{-1}\mathcal{O}_\mathcal{B}}
\otimes_{f^{-1}\mathcal{O}_\mathcal{B}[\mathcal{E}]} \mathcal{O})$.
を考える。補題により、$\xi$ は射
$\delta : \mathcal{I}/\mathcal{I}^2 \to \mathcal{G}$ の類である。
$\mathcal{E}$ を $\mathcal{E} \times_\mathcal{O} \mathcal{O}'$ で
置き換えることにより、$\alpha$ が写像
$\alpha' : \mathcal{E} \to \mathcal{O}'$ を経由すると仮定してもよい。

\medskip\noindent
これらの選択は $f^{-1}\mathcal{O}_{\mathcal{B}'}$-代数準同型
$\varphi : \mathcal{O}_{\mathcal{B}'}[\mathcal{E}] \to \mathcal{O}'$.
を定める。$\mathcal{I}' = \Ker(\varphi)$ と置く。
$\varphi$ は写像
$\varphi|_{\mathcal{I}'} : \mathcal{I}' \to \mathcal{G}$
を誘導し、$\mathcal{O}'$ は次の図式における押し出しである。
$$
\xymatrix{
0 \ar[r] & \mathcal{G} \ar[r] & \mathcal{O}' \ar[r] &
\mathcal{O} \ar[r] & 0 \\
0 \ar[r] & \mathcal{I}' \ar[u]^{\varphi|_{\mathcal{I}'}} \ar[r] &
f^{-1}\mathcal{O}_{\mathcal{B}'}[\mathcal{E}] \ar[u] \ar[r] &
\mathcal{O} \ar[u]_{=} \ar[r] & 0
}
$$
$\psi : \mathcal{I}' \to \mathcal{G}$ を、写像
$\varphi|_{\mathcal{I}'}$ と合成
$$
\mathcal{I}' \to \mathcal{I}'/(\mathcal{I}')^2 \to
\mathcal{I}/\mathcal{I}^2 \xrightarrow{\delta} \mathcal{G}.
$$
との和とする。$\psi$ に沿う押し出しは、上のような図式に収まる
別の環の拡大 $\mathcal{O}'_\xi$ である。計算（省略）により、望みどおり
$o(\mathcal{O}', \mathcal{O}'_\xi) = \xi$ であることが分かる。
\end{proof}

\begin{lemma}
\label{defos-lemma-extensions-of-relative-ringed-topoi}
\leavevmode\par\noindent
$f : (\Sh(\mathcal{C}), \mathcal{O}) \to
(\Sh(\mathcal{B}), \mathcal{O}_\mathcal{B})$ を環付きトポスの射とし、
$\mathcal{G}$ を $\mathcal{O}$-加群とする。
$f^{-1}\mathcal{O}_\mathcal{B}$-代数の拡大
$$
0 \to \mathcal{G} \to \mathcal{O}' \to \mathcal{O} \to 0
$$
で $\mathcal{G}$ が平方零イデアルであるもの\footnote{言い換えれば、
$(\Sh(\mathcal{B}), \mathcal{O}_\mathcal{B})$ 上の一次厚化
$i : (\Sh(\mathcal{C}), \mathcal{O}) \to (\Sh(\mathcal{C}), \mathcal{O}')$
で、$\mathcal{O}$-加群の同型 $\mathcal{G} \to \Ker(i^\sharp)$ を
備えたものの同型類の集合である。}の同型類の集合は、標準的に
\par\noindent
$\Ext^1_\mathcal{O}(\NL_{\mathcal{O}/\mathcal{O}_\mathcal{B}}, \mathcal{G})$
と全単射する。
\end{lemma}

\begin{proof}
これを証明するため、(\ref{defos-equation-to-solve-ringed-topoi}) が図式
$$
\xymatrix{
0 \ar[r] &
\mathcal{G} \ar[r] &
{?} \ar[r] &
\mathcal{O} \ar[r] & 0 \\
0 \ar[r] &
0 \ar[u] \ar[r] &
f^{-1}\mathcal{O}_\mathcal{B} \ar[u] \ar[r]^{\text{id}} &
f^{-1}\mathcal{O}_\mathcal{B} \ar[u] \ar[r] & 0
}
$$
で与えられる場合に、上の結果を適用する。解
$\mathcal{G} \oplus \mathcal{O}$ が存在するという事実と、補題
\ref{defos-lemma-choices-ringed-topoi} から本補題が従う。
（全単射の直接的な構成については次の注意を参照せよ。）
\end{proof}

\begin{remark}
\label{defos-remark-extensions-of-relative-ringed-topoi}
$f : (\Sh(\mathcal{C}), \mathcal{O}) \to
(\mathcal{B}, \mathcal{O}_\mathcal{B})$ と $\mathcal{G}$ は補題
\ref{defos-lemma-extensions-of-relative-ringed-topoi} のとおりとする。
補題のような拡大
$0 \to \mathcal{G} \to \mathcal{O}' \to \mathcal{O} \to 0$
を考える。集合の層 $\mathcal{E}$ と可換図式
$$
\xymatrix{
\mathcal{E} \ar[d]_{\alpha'} \ar[rd]^\alpha \\
\mathcal{O}' \ar[r] & \mathcal{O}
}
$$
で、$f^{-1}\mathcal{O}_\mathcal{B}[\mathcal{E}] \to \mathcal{O}$ が
核 $\mathcal{J}$ をもつ全射となるものを選べる
（例えば $\mathcal{O}'$ へ全射する任意の集合の層を選べばよい）。このとき
$$
\NL_{\mathcal{O}/\mathcal{O}_\mathcal{B}} \cong \NL(\alpha) = 
\left(
\mathcal{J}/\mathcal{J}^2
\longrightarrow
\Omega_{f^{-1}\mathcal{O}_\mathcal{B}[\mathcal{E}]/
f^{-1}\mathcal{O}_\mathcal{B}}
\otimes_{f^{-1}\mathcal{O}_\mathcal{B}[\mathcal{E}]} \mathcal{O}\right)
$$
『サイト上の加群』第\ref{sites-modules-section-netherlander}節、
特に補題 \ref{sites-modules-lemma-NL-up-to-qis} を参照せよ。
もちろん $\alpha'$ は写像
$f^{-1}\mathcal{O}_\mathcal{B}[\mathcal{E}] \to \mathcal{O}'$
を定め、それはさらに写像
$$
\mathcal{J}/\mathcal{J}^2 \longrightarrow \mathcal{G}
$$
を定め、それはさらに
$\Ext^1_\mathcal{O}(\NL(\alpha), \mathcal{G}) =
\Ext^1_\mathcal{O}(\NL_{\mathcal{O}/\mathcal{O}_\mathcal{B}}, \mathcal{G})$
の元で、補題の全単射により $\mathcal{O}'$ に対応するものを定める。
\end{remark}

\begin{lemma}
\label{defos-lemma-extensions-of-relative-ringed-topoi-functorial}
環付きトポスの射
$f : (\Sh(\mathcal{C}), \mathcal{O}_\mathcal{C}) \to
(\Sh(\mathcal{B}), \mathcal{O}_\mathcal{B})$ および
$g : (\Sh(\mathcal{D}), \mathcal{O}_\mathcal{D}) \to
(\Sh(\mathcal{C}), \mathcal{O}_\mathcal{C})$ をとる。
$\mathcal{F}$ を $\mathcal{O}_\mathcal{C}$-加群、$\mathcal{G}$ を
$\mathcal{O}_\mathcal{D}$-加群とし、$c : g^*\mathcal{F} \to \mathcal{G}$
を $\mathcal{O}_\mathcal{D}$-線形写像とする。最後に次を考える。
\begin{enumerate}
\item[(a)]
\leavevmode\par\noindent
$0 \to \mathcal{F} \to \mathcal{O}_{\mathcal{C}'} \to
\mathcal{O}_\mathcal{C} \to 0$
$f^{-1}\mathcal{O}_\mathcal{B}$-代数の拡大で、
$\xi \in \Ext^1_{\mathcal{O}_\mathcal{C}}(
\NL_{\mathcal{O}_\mathcal{C}/\mathcal{O}_\mathcal{B}}, \mathcal{F})$
に対応するもの。
\item[(b)]
\leavevmode\par\noindent
$0 \to \mathcal{G} \to \mathcal{O}_{\mathcal{D}'} \to
\mathcal{O}_\mathcal{D} \to 0$
$g^{-1}f^{-1}\mathcal{O}_\mathcal{B}$-代数の拡大で、
\par\noindent
$\zeta \in \Ext^1_{\mathcal{O}_\mathcal{D}}(
\NL_{\mathcal{O}_\mathcal{D}/\mathcal{O}_\mathcal{B}}, \mathcal{G})$
に対応するもの。
\end{enumerate}
補題 \ref{defos-lemma-extensions-of-relative-ringed-topoi} を参照せよ。
このとき、$g$ および $c$ と両立する
$(\Sh(\mathcal{B}), \mathcal{O}_\mathcal{B})$ 上の環付きトポスの射
$$
g' :
(\Sh(\mathcal{D}), \mathcal{O}_{\mathcal{D}'})
\longrightarrow
(\Sh(\mathcal{C}), \mathcal{O}_{\mathcal{C}'})
$$
が存在するための必要十分条件は、$\xi$ と $\zeta$ が
$\Ext^1_{\mathcal{O}_\mathcal{D}}(
Lg^*\NL_{\mathcal{O}_\mathcal{C}/\mathcal{O}_\mathcal{B}}, \mathcal{G})$
の同じ元に写ることである。
\end{lemma}

\begin{proof}
この主張が意味をなすのは、写像
$$
\Ext^1_{\mathcal{O}_\mathcal{C}}(
\NL_{\mathcal{O}_\mathcal{C}/\mathcal{O}_\mathcal{B}}, \mathcal{F}) \to
\Ext^1_{\mathcal{O}_\mathcal{D}}(
Lg^*\NL_{\mathcal{O}_\mathcal{C}/\mathcal{O}_\mathcal{B}}, Lg^*\mathcal{F}) \to
\Ext^1_{\mathcal{O}_\mathcal{D}}
(Lg^*\NL_{\mathcal{O}_\mathcal{C}/\mathcal{O}_\mathcal{B}}, \mathcal{G})
$$
が写像 $Lg^*\mathcal{F} \to g^*\mathcal{F} \xrightarrow{c} \mathcal{G}$
を用いて得られ、また写像
$$
\Ext^1_{\mathcal{O}_Y}(
\NL_{\mathcal{O}_\mathcal{D}/\mathcal{O}_\mathcal{B}}, \mathcal{G}) \to
\Ext^1_{\mathcal{O}_Y}(
Lg^*\NL_{\mathcal{O}_\mathcal{C}/\mathcal{O}_\mathcal{B}}, \mathcal{G})
$$
が写像 $Lg^*\NL_{\mathcal{O}_\mathcal{C}/\mathcal{O}_\mathcal{B}} \to
\NL_{\mathcal{O}_\mathcal{D}/\mathcal{O}_\mathcal{B}}$ を用いて得られるからである。
補題 \ref{defos-lemma-huge-diagram-ringed-topoi} を図式
$$
\xymatrix{
& 0 \ar[r] &
\mathcal{G} \ar[r] &
\mathcal{O}_{\mathcal{D}'} \ar[r] &
\mathcal{O}_\mathcal{D} \ar[r] & 0 \\
& 0 \ar[r]|\hole & 0 \ar[u] \ar[r] &
g^{-1}f^{-1}\mathcal{O}_\mathcal{B} \ar[u] \ar[r]|\hole &
g^{-1}f^{-1}\mathcal{O}_\mathcal{B} \ar[u] \ar[r] & 0 \\
0 \ar[r] &
\mathcal{F} \ar[ruu] \ar[r] &
\mathcal{O}_{\mathcal{C}'} \ar[r] &
\mathcal{O}_\mathcal{C} \ar[ruu] \ar[r] & 0 \\
0 \ar[r] & 0 \ar[ruu]|\hole \ar[u] \ar[r] &
f^{-1}\mathcal{O}_\mathcal{B} \ar[ruu]|\hole \ar[u] \ar[r] &
f^{-1}\mathcal{O}_\mathcal{B} \ar[ruu]|\hole \ar[u] \ar[r] & 0
}
$$
に適用し、補題 \ref{defos-lemma-extensions-of-relative-ringed-topoi} と
\ref{defos-lemma-huge-diagram-ringed-topoi} の証明中の構成の両立性を用いれば、
本補題の主張が従う。その両立性の主張と証明は省略する
（直接の議論については次の注意を参照せよ）。
\end{proof}

\begin{remark}
\label{defos-remark-extensions-of-relative-ringed-topoi-functorial}
次のデータは補題
\ref{defos-lemma-extensions-of-relative-ringed-topoi-functorial} のとおりとする：
$f : (\Sh(\mathcal{C}), \mathcal{O}_\mathcal{C}) \to
(\Sh(\mathcal{B}), \mathcal{O}_\mathcal{B})$,
$g : (\Sh(\mathcal{D}), \mathcal{O}_\mathcal{D}) \to
(\Sh(\mathcal{C}), \mathcal{O}_\mathcal{C})$,
$\mathcal{F}$,
$\mathcal{G}$,
$c : g^*\mathcal{F} \to \mathcal{G}$,
$0 \to \mathcal{F} \to \mathcal{O}_{\mathcal{C}'} \to
\mathcal{O}_\mathcal{C} \to 0$,
$\xi \in \Ext^1_{\mathcal{O}_\mathcal{C}}(
\NL_{\mathcal{O}_\mathcal{C}/\mathcal{O}_\mathcal{B}}, \mathcal{F})$,
$0 \to \mathcal{G} \to \mathcal{O}_{\mathcal{D}'} \to
\mathcal{O}_\mathcal{D} \to 0$、および
$\zeta \in \Ext^1_{\mathcal{O}_\mathcal{D}}(
\NL_{\mathcal{O}_\mathcal{D}/\mathcal{O}_\mathcal{B}}, \mathcal{G})$。
写像 $c : g^{-1}\mathcal{F} \to \mathcal{G}$ に沿う押し出しにより、拡大
$$
\xymatrix{
0 \ar[r] &
\mathcal{G} \ar[r] &
\mathcal{O}'_1 \ar[r] &
g^{-1}\mathcal{O}_\mathcal{C} \ar[r] & 0 \\
0 \ar[r] &
g^{-1}\mathcal{F} \ar[u]^c \ar[r] &
g^{-1}\mathcal{O}_{\mathcal{C}'} \ar[u] \ar[r] &
g^{-1}\mathcal{O}_\mathcal{C} \ar@{=}[u] \ar[r] & 0
}
$$
を構成できる。写像
$g^\sharp : g^{-1}\mathcal{O}_\mathcal{C} \to \mathcal{O}_\mathcal{D}$
に沿う引き戻しにより、拡大
$$
\xymatrix{
0 \ar[r] &
\mathcal{G} \ar[r] &
\mathcal{O}_{\mathcal{D}'} \ar[r] &
\mathcal{O}_\mathcal{D} \ar[r] & 0 \\
0 \ar[r] &
\mathcal{G} \ar@{=}[u] \ar[r] &
\mathcal{O}'_2 \ar[u] \ar[r] &
g^{-1}\mathcal{O}_\mathcal{C} \ar[u] \ar[r] & 0
}
$$
を構成できる。
図式追跡により、$g$ および $c$ と両立する
$(\Sh(\mathcal{B}), \mathcal{O}_\mathcal{B})$ 上の射
$g' : (\Sh(\mathcal{D}), \mathcal{O}_{\mathcal{D}'}) \to
(\Sh(\mathcal{C}), \mathcal{O}_{\mathcal{C}'})$
が存在するための必要十分条件は、$\mathcal{O}'_1$ と
$\mathcal{O}'_2$ が $g^{-1}\mathcal{O}_\mathcal{C}$ の
$\mathcal{G}$ による $g^{-1}f^{-1}\mathcal{O}_\mathcal{B}$-代数の拡大として
同型であることである。補題
\ref{defos-lemma-extensions-of-relative-ringed-topoi} により、これらの拡大は
次の等式の左辺によって分類される。
$$
\Ext^1_{g^{-1}\mathcal{O}_\mathcal{C}}(
\NL_{g^{-1}\mathcal{O}_\mathcal{C}/g^{-1}f^{-1}\mathcal{O}_\mathcal{B}},
\mathcal{G}) =
\Ext^1_{\mathcal{O}_\mathcal{D}}(
Lg^*\NL_{\mathcal{O}_\mathcal{C}/\mathcal{O}_\mathcal{B}}, \mathcal{G})
$$
この等式はテンソル--Hom 随伴と等式
$$
\vcenter{\halign{\hfil\ensuremath{\displaystyle #}\hfil\cr
\NL_{g^{-1}\mathcal{O}_\mathcal{C}/g^{-1}f^{-1}\mathcal{O}_\mathcal{B}} =
g^{-1}\NL_{\mathcal{O}_\mathcal{C}/\mathcal{O}_\mathcal{B}}\cr
\text{および}\cr
Lg^*\NL_{\mathcal{O}_\mathcal{C}/\mathcal{O}_\mathcal{B}} =
g^{-1}\NL_{\mathcal{O}_\mathcal{C}/\mathcal{O}_\mathcal{B}}
\otimes_{g^{-1}\mathcal{O}_X}^\mathbf{L} \mathcal{O}_Y\cr}}
$$
から従う。第一の等式については『サイト上の加群』、補題
\ref{sites-modules-lemma-pullback-NL} を参照せよ。第二の等式は
導来引き戻しの定義から従う。したがって、補題
\ref{defos-lemma-extensions-of-relative-ringed-topoi-functorial} を示すには、
$\mathcal{O}'_1$ が $\xi$ の像に対応し、$\mathcal{O}'_2$ が
$\zeta$ の像に対応することを示せば十分である。
$\xi$ と $\mathcal{O}'_1$ の対応は、注意
\ref{defos-remark-extensions-of-relative-ringed-topoi} における類 $\xi$ の構成から
直ちに分かる。$\zeta$ と $\mathcal{O}'_2$ の対応については、まず可換図式
$$
\xymatrix{
\mathcal{E} \ar[d]_{\beta'} \ar[rd]^\beta \\
\mathcal{O}_{\mathcal{D}'} \ar[r] & \mathcal{O}_\mathcal{D}
}
$$
で、$g^{-1}f^{-1}\mathcal{O}_\mathcal{B}[\mathcal{E}] \to
\mathcal{O}_\mathcal{D}$ が核 $\mathcal{K}$ をもつ全射となるものを選ぶ。
次に可換図式
$$
\xymatrix{
\mathcal{E} \ar[d]_{\beta'} &
\mathcal{E}' \ar[l]^\varphi \ar[d]_{\alpha'} \ar[rd]^\alpha \\
\mathcal{O}_{\mathcal{D}'} &
\mathcal{O}'_2 \ar[l] \ar[r] &
g^{-1}\mathcal{O}_\mathcal{C}
}
$$
で、$g^{-1}f^{-1}\mathcal{O}_\mathcal{B}[\mathcal{E}'] \to
g^{-1}\mathcal{O}_\mathcal{C}$ が核 $\mathcal{J}$ をもつ全射となるものを選ぶ
（例えば集合の層として
$\mathcal{E}' = \mathcal{E} \amalg \mathcal{O}'_2$ とすればよい）。
写像 $\varphi$ は複体の写像 $\NL(\alpha) \to \NL(\beta)$
（記法は『加群』第\ref{modules-section-netherlander}節のとおり）を誘導し、
特に
$\bar\varphi : \mathcal{J}/\mathcal{J}^2 \to \mathcal{K}/\mathcal{K}^2$
を誘導する。ここで
$\NL(\alpha) \cong \NL_{\mathcal{O}_\mathcal{D}/\mathcal{O}_\mathcal{B}}$
および
$\NL(\beta) \cong
\NL_{g^{-1}\mathcal{O}_\mathcal{C}/g^{-1}f^{-1}\mathcal{O}_\mathcal{B}}$
であり、複体の写像 $\NL(\alpha) \to \NL(\beta)$ は、補題
\ref{defos-lemma-extensions-of-relative-ringed-topoi-functorial} の主張で用いた写像
$Lg^*\NL_{\mathcal{O}_\mathcal{C}/\mathcal{O}_\mathcal{B}} \to
\NL_{\mathcal{O}_\mathcal{D}/\mathcal{O}_\mathcal{B}}$
を表す（その証明の前半を参照せよ）。ここで $\zeta$ は、$\beta'$ が誘導する写像
$\mathcal{K}/\mathcal{K}^2 \to \mathcal{G}$ の類に対応する。注意
\ref{defos-remark-extensions-of-relative-ringed-topoi} を参照せよ。同様に、拡大
$\mathcal{O}'_2$ は $\alpha'$ が誘導する写像
$\mathcal{J}/\mathcal{J}^2 \to \mathcal{G}$ に対応する。
上の可換図式から、この写像は $\beta'$ が誘導する写像
$\mathcal{K}/\mathcal{K}^2 \to \mathcal{G}$ と写像
$\bar\varphi : \mathcal{J}/\mathcal{J}^2 \to \mathcal{K}/\mathcal{K}^2$
との合成である。これで求める両立性が証明された。
\end{remark}

\begin{lemma}
\label{defos-lemma-parametrize-solutions-ringed-topoi}
(\ref{defos-equation-to-solve-ringed-topoi}) のようなデータ
$t : (\Sh(\mathcal{B}), \mathcal{O}_\mathcal{B}) \to
(\Sh(\mathcal{B}'), \mathcal{O}_{\mathcal{B}'})$,
$\mathcal{J} = \Ker(t^\sharp)$,
$f : (\Sh(\mathcal{C}), \mathcal{O}) \to
(\Sh(\mathcal{B}), \mathcal{O}_\mathcal{B})$、$\mathcal{G}$、および
$c : \mathcal{J} \to \mathcal{G}$ をとる。
補題 \ref{defos-lemma-extensions-of-relative-ringed-topoi} により、
$\mathcal{O}_\mathcal{B}$ の $\mathcal{J}$ による拡大
$\mathcal{O}_{\mathcal{B}'}$ に対応する元を
$\xi \in \Ext^1_{\mathcal{O}_\mathcal{B}}(
\NL_{\mathcal{O}_\mathcal{B}/\mathcal{O}_{\mathcal{B}'}}, \mathcal{J})$
と記す。解の同型類の集合は、写像
$$
\Ext^1_\mathcal{O}(\NL_{\mathcal{O}/\mathcal{O}_{\mathcal{B}'}},
\mathcal{G})\to
\Ext^1_\mathcal{O}(
Lf^*\NL_{\mathcal{O}_\mathcal{B}/\mathcal{O}_{\mathcal{B}'}}, \mathcal{G})
$$
の $\xi$ の像上のファイバーと標準的に全単射する。
\end{lemma}

\begin{proof}
補題 \ref{defos-lemma-extensions-of-relative-ringed-topoi} を
$t \circ f : (\Sh(\mathcal{C}), \mathcal{O}) \to
(\Sh(\mathcal{B}'), \mathcal{O}_{\mathcal{B}'})$
と $\mathcal{O}$-加群 $\mathcal{G}$ に適用すると、
$\Ext^1_\mathcal{O}(\NL_{\mathcal{O}/\mathcal{O}_{\mathcal{B}'}},
\mathcal{G})$ の元 $\zeta$ は
$0 \to \mathcal{G} \to \mathcal{O}' \to \mathcal{O} \to 0$
という $f^{-1}\mathcal{O}_{\mathcal{B}'}$-代数の拡大をパラメータ付ける。
補題 \ref{defos-lemma-extensions-of-relative-ringed-topoi-functorial} を
$$
(\Sh(\mathcal{C}), \mathcal{O}) \xrightarrow{f}
(\Sh(\mathcal{B}), \mathcal{O}_\mathcal{B}) \xrightarrow{t}
(\Sh(\mathcal{B}'), \mathcal{O}_{\mathcal{B}'})
$$
と $c : \mathcal{J} \to \mathcal{G}$ に適用すると、$c$ および $f$ と
両立する $(\Sh(\mathcal{B}'), \mathcal{O}_{\mathcal{B}'})$ 上の射
$$
f' : 
(\Sh(\mathcal{C}), \mathcal{O}')
\longrightarrow
(\Sh(\mathcal{B}'), \mathcal{O}_{\mathcal{B}'})
$$
が存在するための必要十分条件は、$\zeta$ が $\xi$ に写ることである。
もちろん、これは $\mathcal{O}'$ が
(\ref{defos-equation-to-solve-ringed-topoi}) の解であるということにほかならない。
\end{proof}






















\section{代数空間の変形}
\label{defos-section-deformations-spaces}

\noindent
本節では、第\ref{defos-section-deformations-ringed-topoi}節の結果が
代数空間の変形に対して何を意味するかを詳述する。

\begin{lemma}
\label{defos-lemma-match-thickenings}
$S$ をスキームとし、$i : Z \to Z'$ を $S$ 上の代数空間の射とする。
次は同値である。
\begin{enumerate}
\item $i$ は『空間の射詳論』第
\ref{spaces-more-morphisms-section-thickenings}節で定義された
代数空間の厚化である。
\item 付随する環付きトポスの射
$i_{small} : (\Sh(Z_\etale), \mathcal{O}_Z) \to
(\Sh(Z'_\etale), \mathcal{O}_{Z'})$
（『空間の性質』、補題
\ref{spaces-properties-lemma-morphism-ringed-topoi}）は、第
\ref{defos-section-thickenings-ringed-topoi}節の意味で厚化である。
\end{enumerate}
\end{lemma}

\begin{proof}
これは自明なことではないと強調しておく。

\medskip\noindent
(1) を仮定する。『空間の射詳論』、補題
\ref{spaces-more-morphisms-lemma-thickening-equivalence}
により、射 $i$ は小エタールサイトの同値、したがって特にトポスの同値を
誘導する。もちろん厚化の定義により、$i^\sharp$ は全射であり、
その核は局所冪零である。

\medskip\noindent
(2) を仮定する。（この向きは重要性が低く、むしろ興味深い余談である。）
任意のエタール射 $Y' \to Z'$ に対して、
$Y = Z \times_{Z'} Y'$ のエタール・トポスは $Y'$ のものと同じである。
特に、準コンパクトであることはトポス論的な概念なので
（『サイト』、補題 \ref{sites-lemma-quasi-compact}）、
$Y'$ が準コンパクトであることと $Y$ が準コンパクトであることは同値である。
このことから、$Y'$ が準コンパクトかつ準分離であることと、
$Y$ が準コンパクトかつ準分離であることも同値である
（実際、$Y'$ が準分離であることは、$Y'$ 上エタールな任意の
準コンパクト代数空間 $Y'_1, Y'_2$ に対して
$Y'_1 \times_{Y'} Y'_2$ が準コンパクトであることによって特徴づけられる）。
$Y'$ をアフィンとする。このとき代数空間 $Y$ は準コンパクトかつ準分離である。
任意の準連接 $\mathcal{O}_Y$-加群 $\mathcal{F}$ に対して、
エタール・トポスが同じであることから
$H^q(Y, \mathcal{F}) = H^q(Y', (Y \to Y')_*\mathcal{F})$ である。
さらに、押し出しが準連接であり
（『空間の射』、補題 \ref{spaces-morphisms-lemma-pushforward}）、
$Y$ がアフィンであることから
$H^q(Y', (Y \to Y')_*\mathcal{F}) = 0$ である。
したがって『空間のコホモロジー』、命題
\ref{spaces-cohomology-proposition-vanishing-affine} により $Y'$ はアフィンである
（この補題の向きには、この参照を避ける証明もきっとあるはずである）。
ゆえに $i$ はアフィン射である。アフィンの場合には、第
\ref{defos-section-thickenings-ringed-topoi}節の条件から、
$i$ が代数空間の厚化であることが容易に従う。
\end{proof}

\begin{lemma}
\label{defos-lemma-deform-spaces}
$S$ をスキームとする。
$Y \subset Y'$ を $S$ 上の代数空間の一次厚化とし、
$f : X \to Y$ を $S$ 上の代数空間の平坦射とする。
$S$ 上の代数空間の平坦射 $f' : X' \to Y'$ と、$Y$ 上の同型
$a : X \to X' \times_{Y'} Y$ が存在するとき、次が成り立つ。
\begin{enumerate}
\item 組 $(f' : X' \to Y', a)$ の同型類の集合は
$\Ext^1_{\mathcal{O}_X}(\NL_{X/Y}, f^*\mathcal{C}_{Y/Y'})$
の下で主等質である。
\item $Y'$ 上の自己同型 $\varphi : X' \to X'$ で、
$X' \times_{Y'} Y$ 上では恒等射に帰着するもの全体の集合は
$\Ext^0_{\mathcal{O}_X}(\NL_{X/Y}, f^*\mathcal{C}_{Y/Y'})$ である。
\end{enumerate}
\end{lemma}

\begin{proof}
環付きトポスの変形に関する結果を、この補題に現れる代数空間の
小エタール・トポスに適用する。
$X$ を $X'$ の閉部分空間とみなせば、
$(f, f') : (X \subset X') \to (Y \subset Y')$ は一次厚化の射となる。
補題 \ref{defos-lemma-match-thickenings} により、これは環付きトポスの厚化の射に移される。
すると『空間の射詳論』、補題
\ref{spaces-more-morphisms-lemma-deform}
（または、より一般的な補題 \ref{defos-lemma-deform-module-ringed-topoi}）から、
$X'$ における $X$ のイデアル層は $f^*\mathcal{C}_{Y'/Y}$ に等しく、
しかもこのことは $X'$ の $Y'$ 上での平坦性と同値であることが分かる。
したがって、次の可換図式を得る。
$$
\xymatrix{
0 \ar[r] & f^*\mathcal{C}_{Y/Y'} \ar[r] &
\mathcal{O}_{X'} \ar[r] &
\mathcal{O}_X \ar[r] & 0 \\
0 \ar[r] &
f_{small}^{-1}\mathcal{C}_{Y/Y'} \ar[u] \ar[r] &
f_{small}^{-1}\mathcal{O}_{Y'} \ar[u] \ar[r] &
f_{small}^{-1}\mathcal{O}_Y \ar[u] \ar[r] & 0
}
$$
(\ref{defos-equation-to-solve-ringed-topoi}) と比較されたい。
(2) の自己同型 $\varphi$ は、上の図式を可換にする自己同型
$\varphi^\sharp : \mathcal{O}_{X'} \to \mathcal{O}_{X'}$ を与える。
逆に、上の層の図式を可換にする自己同型
$\alpha : \mathcal{O}_{X'} \to \mathcal{O}_{X'}$ は、
『空間の射詳論』、補題
\ref{spaces-more-morphisms-lemma-first-order-thickening-maps} により、
(2) のある自己同型 $\varphi$ に対する $\varphi^\sharp$ に等しい。
最後に、『空間の射詳論』、補題
\ref{spaces-more-morphisms-lemma-first-order-thickening} によれば、
$X_\etale$ 上の別の環の層 $\mathcal{A}$ で次の図式を可換にするものが見つかれば、
$$
\xymatrix{
0 \ar[r] & f^*\mathcal{C}_{Y/Y'} \ar[r] &
\mathcal{A} \ar[r] &
\mathcal{O}_X \ar[r] & 0 \\
0 \ar[r] &
f_{small}^{-1}\mathcal{C}_{Y/Y'} \ar[u] \ar[r] &
f_{small}^{-1}\mathcal{O}_{Y'} \ar[u] \ar[r] &
f_{small}^{-1}\mathcal{O}_Y \ar[u] \ar[r] & 0
}
$$
すると $\mathcal{O}_{X''} = \mathcal{A}$ を満たす一次厚化
$X \subset X''$ が存在する。さらに『空間の射詳論』、補題
\ref{spaces-more-morphisms-lemma-first-order-thickening-maps} をもう一度適用すると、
望む性質をすべて備えた射
$(f, f'') : (X \subset X'') \to (Y \subset Y')$ を得る。
したがって (1) は補題 \ref{defos-lemma-choices-ringed-topoi} から従い、
(2) は補題 \ref{defos-lemma-huge-diagram-ringed-topoi} の (2) から従う。
（『空間の射詳論』第
\ref{spaces-more-morphisms-section-netherlander}節で代数空間の射に対して定義された
$\NL_{X/Y}$ は、第\ref{defos-section-deformations-ringed-topoi}節で用いられる
$\NL_{X/Y}$ と一致することに注意せよ。）
\end{proof}

\noindent
$S$ をスキームとし、$f : X \to B$ を $S$ 上の代数空間の射とする。
$\mathcal{F} \to \mathcal{G}$ を $\mathcal{O}_X$-加群の準同型とする
（準連接であるとは限らない）。次の関手を考える。
$$
F :
\left\{
\begin{matrix}
f^{-1}\mathcal{O}_B\text{-代数の拡大}\\
0 \to \mathcal{F} \to \mathcal{O}' \to \mathcal{O}_X \to 0\\
\mathcal{F}\text{ は平方零イデアル}
\end{matrix}
\right\}
\longrightarrow
\left\{
\begin{matrix}
f^{-1}\mathcal{O}_B\text{-代数の拡大}\\
0 \to \mathcal{G} \to \mathcal{O}' \to \mathcal{O}_X \to 0\\
\mathcal{G}\text{ は平方零イデアル}
\end{matrix}
\right\}
$$
これは押し出しによって定まる。

\begin{lemma}
\label{defos-lemma-thickening-space-quasi-coherent}
上の状況で、$X$ は準コンパクトかつ準分離であり、
$DQ_X(\mathcal{F}) \to DQ_X(\mathcal{G})$
（『空間の導来圏』第
\ref{spaces-perfect-section-better-coherator}節）は同型であると仮定する。
このとき関手 $F$ は圏同値である。
\end{lemma}

\begin{proof}
$\NL_{X/B}$ は $D_\QCoh(\mathcal{O}_X)$ の対象であることを想起しよう。
『空間の射詳論』、補題
\ref{spaces-more-morphisms-lemma-netherlander-quasi-coherent} を参照せよ。
したがって仮定から、すべての $i$ に対して写像
$$
\Ext^i_X(\NL_{X/B}, \mathcal{F}) \longrightarrow
\Ext^i_X(\NL_{X/B}, \mathcal{G})
$$
が同型であることが従う。補題 \ref{defos-lemma-huge-diagram-ringed-topoi} により、
これは関手が充満忠実であることを意味する。
一方、両方の圏に解 $\mathcal{O}_X \oplus \mathcal{F}$ と
$\mathcal{O}_X \oplus \mathcal{G}$ があるので、補題
\ref{defos-lemma-choices-ringed-topoi} により、この関手は本質的全射でもある。
\end{proof}

\noindent
$S$ をスキームとする。$B \subset B'$ を $S$ 上の代数空間の一次厚化とし、
そのイデアル層を $\mathcal{J}$ とする。これを準連接
$\mathcal{O}_B$-加群とも、$B'$ 上の準連接イデアル層ともみなす。
『空間の射詳論』第
\ref{spaces-more-morphisms-section-thickenings}節を参照せよ。
$f : X \to B$ を $S$ 上の代数空間の射とし、
$\mathcal{F} \to \mathcal{G}$ を $\mathcal{O}_X$-加群の準同型とする
（準連接であるとは限らない）。$c : f^{-1}\mathcal{J} \to \mathcal{F}$ を
$f^{-1}\mathcal{O}_B$-加群の写像とし、その合成を
$c' : f^{-1}\mathcal{J} \to \mathcal{G}$ と記す。次の関手を考える。
$$
FT :
\{(\ref{defos-equation-to-solve-ringed-topoi})\text{ の }\mathcal{F}\text{ と }c
\text{ に対する解}\}
\longrightarrow
\{(\ref{defos-equation-to-solve-ringed-topoi})\text{ の }\mathcal{G}\text{ と }c'
\text{ に対する解}\}
$$
これは押し出しによって定まる。

\begin{lemma}
\label{defos-lemma-thickening-over-thickening-space-quasi-coherent}
上の状況で、$X$ は準コンパクトかつ準分離であり、
$DQ_X(\mathcal{F}) \to DQ_X(\mathcal{G})$
（『空間の導来圏』第
\ref{spaces-perfect-section-better-coherator}節）は同型であると仮定する。
このとき関手 $FT$ は圏同値である。
\end{lemma}

\begin{proof}
(\ref{defos-equation-to-solve-ringed-topoi}) の $\mathcal{F}$ に対する解は、特に
$f^{-1}\mathcal{O}_{B'}$-代数の拡大
$$
0 \to \mathcal{F} \to \mathcal{O}' \to \mathcal{O}_X \to 0
$$
を与え、ここで $\mathcal{F}$ は平方零イデアルである。
$\mathcal{G}$ についても同様である。また、このような拡大が与えられると、
写像 $c_{\mathcal{O}'} : f^{-1}\mathcal{J} \to \mathcal{F}$ を得る。
したがって、考えているのは、このような $f^{-1}\mathcal{O}_{B'}$-代数の拡大のうち
$c = c_{\mathcal{O}'}$ を満たすものからなる充満部分圏である。
$F$ を補題 \ref{defos-lemma-thickening-space-quasi-coherent} の同値
（今回は $X \to B'$ に適用する）として
$\mathcal{O}'' = F(\mathcal{O}')$ ならば、明らかに
$c_{\mathcal{O}''}$ は $c_{\mathcal{O}'}$ と写像
$\mathcal{F} \to \mathcal{G}$ の合成である。これで補題が証明された。
\end{proof}





\section{複体の変形}
\label{defos-section-deformations-complexes}

\noindent
本節は次節への準備である。
可能な限り、可換代数に関する各章の結果を利用する。

\begin{lemma}
\label{defos-lemma-canonical-class-algebra}
$R' \to R$ を環の全射とし、その核を平方零イデアル $I$ とする。
任意の $K \in D^-(R)$ に対して、$D(R)$ に標準射
$$
\omega(K) : K \longrightarrow K \otimes_R^\mathbf{L} I[2]
$$
が存在し、次の性質をもつ。
\begin{enumerate}
\item $\omega(K) = 0$ であることと、
$K' \otimes_{R'}^\mathbf{L} R = K$ を満たす $K' \in D(R')$ が存在することは同値である。
\item $D^-(R)$ における $K \to L$ が与えられると、図式
$$
\xymatrix{
K \ar[d] \ar[rr]_-{\omega(K)} & &
K \otimes^\mathbf{L}_R I[2] \ar[d] \\
L \ar[rr]^-{\omega(L)} & &
L \otimes^\mathbf{L}_R I[2]
}
$$
は可換である。
\item $\omega(K)$ の構成は環準同型 $R' \to S'$ と両立する
（正確な主張は証明を参照せよ）。
\end{enumerate}
\end{lemma}

\begin{proof}
$K$ を表す自由 $R$-加群の上に有界な複体 $K^\bullet$ を選ぶ。
このとき $K^n$ を持ち上げる自由 $R'$-加群 $(K')^n$ を選べる。
さらに、$K^\bullet$ の微分 $d^n_K : K^n \to K^{n + 1}$ を持ち上げる
$R'$-加群の写像 $(d')^n_K : (K')^n \to (K')^{n + 1}$ を選べる。
合成
$$
(d')^{n + 1}_K \circ (d')^n_K : (K')^n \to (K')^{n + 2}
$$
は零であるとは限らないが、$d^{n + 1} \circ d^n = 0$ であるため、実際には
$$
(K')^n \to K^n \xrightarrow{\omega^n_K}
K^{n + 2} \otimes_R I = I(K')^{n + 2} \to (K')^{n + 2}
$$
と分解する。計算により $\omega^n_K$ が複体の写像を定めることが分かる。
この複体の写像が $\omega(K)$ を定める。

\medskip\noindent
この構成が、上に有界な自由 $R$-加群の複体の写像
$\alpha^\bullet : K^\bullet \to L^\bullet$ および持ち上げの選択
$(K')^n, (L')^n, (d')^n_K, (d')^n_L$ と両立することを示そう。
そこで、成分 $\alpha^n : K^n \to L^n$ を持ち上げる
$(\alpha')^n : (K')^n \to (L')^n$ を選ぶ。先ほどと同様に、
$(d')^n_L \circ (\alpha')^n - (\alpha')^{n + 1} \circ (d')_K^n$ の分解
$$
(K')^n \to K^n \xrightarrow{h^n}
L^{n + 1} \otimes_R I = I(L')^{n + 1} \to (L')^{n + 2}
$$
を得る。このとき、計算により
$$
\omega^n_L \circ \alpha^n =
(d_L^{n + 1} \otimes \text{id}_I) \circ h^n + h^{n + 1} \circ d_K^n +
(\alpha^{n + 2} \otimes \text{id}_I) \circ \omega^n_K
$$
が示される。これにより、この特別な場合について補題の (2) の図式の
可換性が証明される。$K$ を表す上に有界な自由複体の二つの異なる選択に
これを適用すると、$\omega(K)$ が良定義であることが分かる。
もちろん、一般の場合にも (2) が成り立つ。

\medskip\noindent
$K$ が $D^-(R')$ における $K'$ に持ち上がるならば、
$K'$ を自由 $R'$-加群の上に有界な複体で表すことができ、
$\omega(K) = 0$ であることが直ちに分かる。
逆に、選択 $K^\bullet$, $(K')^n$, $(d')^n_K$ に戻る。
$\omega(K) = 0$ ならば、次を満たす
$g^n : K^n \to K^{n + 1} \otimes_R I$ を見つけられる。
$$
\omega^n = (d_K^{n + 1} \otimes \text{id}_I) \circ g^n +
g^{n + 1} \circ d_K^n
$$
これは、微分 $(d')^n_K - g^n : (K')^n \to (K')^{n + 1}$ によって、
$K^\bullet$ を持ち上げる自由 $R'$-加群の複体が得られることを意味する。
これで (1) が証明された。

\medskip\noindent
最後に、(3) の意味は次のとおりである。$R' \to S'$ を環準同型とする。
$S = S' \otimes_{R'} R$ と置き、$S' \to S$ の平方零な核を
$J = IS' \subset S'$ と記す。このとき $K \in D^-(R)$ が与えられると、
次の可換図式を得るというのが主張である。
$$
\xymatrix{
K \otimes_R^\mathbf{L} S \ar[d] \ar[rr]_-{\omega(K) \otimes \text{id}} & &
(K \otimes^\mathbf{L}_R I[2]) \otimes_R^\mathbf{L} S \ar[d] \\
K \otimes_R^\mathbf{L} S \ar[rr]^-{\omega(K \otimes_R^\mathbf{L} S)} & &
(K \otimes_R^\mathbf{L} S) \otimes^\mathbf{L}_S J[2]
}
$$
ここで右の垂直射は
$$
(K \otimes^\mathbf{L}_R I[2]) \otimes_R^\mathbf{L} S =
(K \otimes_R^\mathbf{L} S) \otimes_S^\mathbf{L}
(I \otimes_R^\mathbf{L} S)[2] \longrightarrow
(K \otimes_R^\mathbf{L} S) \otimes_S^\mathbf{L} J[2]
$$
から得られる。上と同様に $K^\bullet$, $(K')^n$, $(d')^n_K$ を選ぶ。
すると $\omega(K \otimes_R^\mathbf{L} S)$ の構成には、
$K^\bullet \otimes_R S$, $(K')^n \otimes_{R'} S'$, および
$(d')^n_K \otimes \text{id}_{S'}$ を用いることができる。
これらの選択の下では、複体の写像の段階で可換性が直ちに確かめられる。
\end{proof}







\section{環付きトポス上の複体の変形}
\label{defos-section-thickenings-complexes}

\noindent
本節の内容は \cite{lieblich-complexes} に基づく。

\medskip\noindent
本節の内容は、第\ref{defos-section-thickenings-ringed-topoi}節で定義した
環付きトポスの一次厚化の設定で成り立つ。ただし記法を簡単にするため、
基礎となるサイト $\mathcal{C}$ と $\mathcal{D}$ は同じであると仮定する。
さらに、環の層の全射準同型 $\mathcal{O}' \to \mathcal{O}$ を、
文献でより一般的と思われる記法に従って
$\mathcal{O} \to \mathcal{O}_0$ と記す。

\begin{lemma}
\label{defos-lemma-lift-complex}
$\mathcal{C}$ をサイトとし、$\mathcal{O} \to \mathcal{O}_0$ を
環の層の全射とする。次のデータが与えられていると仮定する。
\begin{enumerate}
\item 平坦 $\mathcal{O}$-加群 $\mathcal{G}^n$。
\item $\mathcal{O}$-加群の写像 $\mathcal{G}^n \to \mathcal{G}^{n + 1}$。
\item $\mathcal{O}_0$-加群の複体 $\mathcal{K}_0^\bullet$。
\item $\mathcal{O}$-加群の写像 $\mathcal{G}^n \to \mathcal{K}_0^n$。
\end{enumerate}
これらは次を満たすものとする。
\begin{enumerate}
\item[(a)] $n \gg 0$ に対して $H^n(\mathcal{K}_0^\bullet) = 0$。
\item[(b)] $n \gg 0$ に対して $\mathcal{G}^n = 0$。
\item[(c)]
$\mathcal{G}^n_0 = \mathcal{G}^n \otimes_\mathcal{O} \mathcal{O}_0$
と置くと、誘導される写像は複体 $\mathcal{G}_0^\bullet$ と
複体の写像 $\mathcal{G}_0^\bullet \to \mathcal{K}_0^\bullet$ を定める。
\end{enumerate}
このとき次が存在する。
\begin{enumerate}
\item[(\romannumeral1)]
平坦 $\mathcal{O}$-加群 $\mathcal{F}^n$。
\item[(\romannumeral2)]
$\mathcal{O}$-加群の写像 $\mathcal{F}^n \to \mathcal{F}^{n + 1}$。
\item[(\romannumeral3)]
$\mathcal{O}$-加群の写像 $\mathcal{F}^n \to \mathcal{K}_0^n$。
\item[(\romannumeral4)]
$\mathcal{O}$-加群の写像 $\mathcal{G}^n \to \mathcal{F}^n$。
\end{enumerate}
$n \gg 0$ に対して $\mathcal{F}^n = 0$ であり、すべての $n$ に対して図式
$$
\xymatrix{
\mathcal{G}^n \ar[r] \ar[d] & \mathcal{G}^{n + 1} \ar[d] \\
\mathcal{F}^n \ar[r] & \mathcal{F}^{n + 1}
}
$$
が可換であり、合成
$\mathcal{G}^n \to \mathcal{F}^n \to \mathcal{K}_0^n$ は与えられた写像
$\mathcal{G}^n \to \mathcal{K}_0^n$ である。さらに
$\mathcal{F}^n_0 = \mathcal{F}^n \otimes_\mathcal{O} \mathcal{O}_0$ と置くと、
複体 $\mathcal{F}_0^\bullet$ と、擬同型である複体の写像
$\mathcal{F}_0^\bullet \to \mathcal{K}_0^\bullet$ が得られる。
\end{lemma}

\begin{proof}
$e$ に関する降下帰納法により、$n \geq e$ に対する $\mathcal{F}^n$、
$\mathcal{G}^n \to \mathcal{F}^n$、および
$\mathcal{F}^n \to \mathcal{F}^{n + 1}$ で、次の可換図式に入るものを
見つけられることを示す。
$$
\xymatrix{
\ldots \ar[r] &
\mathcal{G}^{e - 1} \ar[r] \ar@/_2pc/[dd] &
\mathcal{G}^e \ar[d] \ar[r] \ar@/_2pc/[dd] &
\mathcal{G}^{e + 1} \ar[d] \ar[r] \ar@/_2pc/[dd]|\hole &
\ldots \\
& &
\mathcal{F}^e \ar[d] \ar[r] &
\mathcal{F}^{e + 1} \ar[d] \ar[r] & \ldots \\
\ldots  \ar[r] &
\mathcal{K}_0^{e - 1} \ar[r] &
\mathcal{K}_0^e \ar[r] &
\mathcal{K}_0^{e + 1} \ar[r] & \ldots
}
$$
ここで $\mathcal{F}_0^\bullet$ は複体であり、誘導される写像
$\mathcal{F}_0^\bullet \to \mathcal{K}_0^\bullet$ は、$n > e$ では
$H^n$ 上の同型を、$n = e$ では全射を誘導するものとする。
$e \gg 0$ の場合、仮定 (a), (b) により $n \geq e$ に対して
$\mathcal{F}^n = 0$ と取れるので、これは成り立つ。

\medskip\noindent
帰納段階。$\mathcal{F}^{e - 1}$ と写像
$\mathcal{G}^{e - 1} \to \mathcal{F}^{e - 1}$、
$\mathcal{F}^{e - 1} \to \mathcal{F}^e$、
$\mathcal{F}^{e - 1} \to \mathcal{K}_0^{e - 1}$ を構成しなければならない。
$\mathcal{F}^{e - 1} = A \oplus B \oplus C$ を三つの部分の直和として選ぶ。

\medskip\noindent
第一の部分には $A = \mathcal{G}^{e - 1}$ と取り、写像
$\mathcal{G}^{e - 1} \to \mathcal{F}^{e - 1}$ を第一直和因子への包含とする。
写像 $A \to \mathcal{K}^{e - 1}_0$ と $A \to \mathcal{F}^e$ は明らかなものとする。

\medskip\noindent
$B$ を選ぶため、帰納法の仮定による全射
$$
\gamma :
\Ker(\mathcal{F}^e_0 \to \mathcal{F}^{e + 1}_0)
\longrightarrow
\Ker(\mathcal{K}^e_0 \to \mathcal{K}^{e + 1}_0)/
\Im(\mathcal{K}^{e - 1}_0 \to \mathcal{K}^e_0)
$$
を考える。集合 $I$、各 $i \in I$ に対する $\mathcal{C}$ の対象 $U_i$、および切断
$s_i \in \mathcal{F}^e(U_i)$, $t_i \in \mathcal{K}^{e - 1}_0(U_i)$ を、
次を満たすように選べる。
\begin{enumerate}
\item $s_i$ は $\Ker(\gamma) \subset
\Ker(\mathcal{F}^e_0 \to \mathcal{F}^{e + 1}_0)$ の切断に写る。
\item $s_i$ と $t_i$ は $\mathcal{K}^e_0$ の同じ切断に写る。
\item 切断 $s_i$ は $\mathcal{O}_0$-加群として $\Ker(\gamma)$ を生成する。
\end{enumerate}
このことの完全な正当化は省略する。ここでは
$\mathcal{F}^e \to \mathcal{F}^e_0$ が集合の層の全射であることを用いる。
そこで
$$
B = \bigoplus\nolimits_{i \in I} j_{U_i!}\mathcal{O}_{U_i}
$$
と置き、$s_i$ と $t_i$ によって直和因子
$j_{U_i!}\mathcal{O}_{U_i}$ の行き先を定めることで、写像
$B \to \mathcal{F}^e$ と $B \to \mathcal{K}_0^{e - 1}$ を定義する。

\medskip\noindent
$\mathcal{F}^{e - 1} = A \oplus B$ とし、写像を上のように取ると、
$e - 1$ に対する上記の図式が得られ、
$\mathcal{F}_0^\bullet \to \mathcal{K}_0^\bullet$ は $n \geq e$ で
$H^n$ 上の同型を誘導する。この写像を $H^{e - 1}$ 上で全射にするため、
直和因子 $C$ を次のように選ぶ。集合 $J$、各 $j \in J$ に対する
$\mathcal{C}$ の対象 $U_j$、および $U_j$ 上の
$\Ker(\mathcal{K}^{e - 1}_0 \to \mathcal{K}^e_0)$ の切断 $t_j$ を、
これらの切断がこの核を $\mathcal{O}_0$ 上生成するように選ぶ。そこで
$$
C = \bigoplus\nolimits_{j \in J} j_{U_j!}\mathcal{O}_{U_j}
$$
と置く。$C \to \mathcal{F}^e$ は零写像とし、$s_j$ を用いて直和因子
$j_{U_j!}\mathcal{O}_{U_j}$ の行き先を定めることで、写像
$C \to \mathcal{K}_0^{e - 1}$ を定める。
$\mathcal{F}^{e - 1} = A \oplus B \oplus C$ とし、写像を上記のように取れば、
帰納段階が完了する。
\end{proof}

\begin{lemma}
\label{defos-lemma-canonical-class}
$\mathcal{C}$ をサイトとする。$\mathcal{O} \to \mathcal{O}_0$ を環の層の全射とし、
その核を平方零イデアル層 $\mathcal{I}$ とする。
$D^-(\mathcal{O}_0)$ の任意の対象 $K_0$ に対して、$D(\mathcal{O}_0)$ に標準射
$$
\omega(K_0) :
K_0 \longrightarrow
K_0 \otimes_{\mathcal{O}_0}^\mathbf{L} \mathcal{I}[2]
$$
が存在し、$D^-(\mathcal{O}_0)$ の任意の写像 $K_0 \to L_0$ に対して図式
$$
\xymatrix{
K_0 \ar[d] \ar[rr]_-{\omega(K_0)} & &
(K_0 \otimes^\mathbf{L}_{\mathcal{O}_0} \mathcal{I})[2] \ar[d] \\
L_0 \ar[rr]^-{\omega(L_0)} & &
(L_0 \otimes^\mathbf{L}_{\mathcal{O}_0} \mathcal{I})[2]
}
$$
は可換である。
\end{lemma}

\begin{proof}
$K_0$ を任意の $\mathcal{O}_0$-加群の複体
$\mathcal{K}_0^\bullet$ で表す。すべての $n$ に対して
$\mathcal{G}^n = 0$ として補題 \ref{defos-lemma-lift-complex} を適用する。
補題によって得られる写像を
$d : \mathcal{F}^n \to \mathcal{F}^{n + 1}$ と記す。このとき
$d \circ d : \mathcal{F}^n \to \mathcal{F}^{n + 2}$ は
$\mathcal{I}$ を法として零である。$\mathcal{F}^n$ は平坦なので、
$\mathcal{I}\mathcal{F}^n =
\mathcal{F}^n \otimes_{\mathcal{O}} \mathcal{I} =
\mathcal{F}^n_0 \otimes_{\mathcal{O}_0} \mathcal{I}$.
である。したがって複体の標準射
$$
d \circ d : \mathcal{F}_0^\bullet
\longrightarrow
(\mathcal{F}_0^\bullet \otimes_{\mathcal{O}_0} \mathcal{I})[2]
$$
を得る。$\mathcal{F}_0^\bullet$ は平坦 $\mathcal{O}_0$-加群の
上に有界な複体なので K-平坦であり、導来テンソル積の計算に使える。
さらに、複体の写像 $\mathcal{F}_0^\bullet \to \mathcal{K}_0^\bullet$ は
構成により擬同型である。したがって、構成した写像の始域と終域はそれぞれ
$K_0$ と $K_0 \otimes_{\mathcal{O}_0}^\mathbf{L} \mathcal{I}[2]$ を表し、
写像 $\omega(K_0)$ が得られる。

\medskip\noindent
この手順が複体の写像と両立することを示そう。
$\mathcal{L}_0^\bullet$ を $D^-(\mathcal{O}_0)$ の別の対象を表す複体とし、
$$
\mathcal{K}_0^\bullet \longrightarrow \mathcal{L}_0^\bullet
$$
が複体の写像であるとする。複体 $\mathcal{L}_0^\bullet$、平坦加群
$\mathcal{F}^n$、写像 $\mathcal{F}^n \to \mathcal{F}^{n + 1}$、および合成
$\mathcal{F}^n \to \mathcal{K}_0^n \to \mathcal{L}_0^n$
に補題 \ref{defos-lemma-lift-complex} を適用する
（文字の使い方が逆転していることをお詫びする）。補題の性質をすべて満たす
平坦加群 $\mathcal{G}^n$、写像 $\mathcal{F}^n \to \mathcal{G}^n$、
$\mathcal{G}^n \to \mathcal{G}^{n + 1}$、および
$\mathcal{G}^n \to \mathcal{L}_0^n$ を得る。このとき明らかに
$$
\xymatrix{
\mathcal{F}_0^\bullet \ar[d] \ar[r] &
(\mathcal{F}_0^\bullet \otimes_{\mathcal{O}_0} \mathcal{I})[2] \ar[d] \\
\mathcal{G}_0^\bullet \ar[r] &
(\mathcal{G}_0^\bullet \otimes_{\mathcal{O}_0} \mathcal{I})[2]
}
$$
は複体の可換図式である。

\medskip\noindent
次に $\omega(K_0)$ が良定義であることを示す。$K_0$ を表す二つの
$\mathcal{O}_0$-加群の複体 $\mathcal{K}_0^\bullet$ と
$(\mathcal{K}'_0)^\bullet$、および上記のような二つの系
$(\mathcal{F}^n, d : \mathcal{F}^n \to \mathcal{F}^{n + 1},
\mathcal{F}^n \to \mathcal{K}_0^n)$
および
$((\mathcal{F}')^n, d : (\mathcal{F}')^n \to (\mathcal{F}')^{n + 1},
(\mathcal{F}')^n \to \mathcal{K}_0^n)$
が与えられたとする。この二つの複体が導来圏で $K_0$ を表すことを実現する
複体 $(\mathcal{K}''_0)^\bullet$ と擬同型
$\mathcal{K}_0^\bullet  \to (\mathcal{K}''_0)^\bullet$
および
$(\mathcal{K}'_0)^\bullet  \to (\mathcal{K}''_0)^\bullet$
を選べる。次に、前段落の結果を
$$
(\mathcal{K}_0)^\bullet \oplus (\mathcal{K}'_0)^\bullet
\longrightarrow
(\mathcal{K}''_0)^\bullet
$$
に適用する。これにより可換図式
$$
\xymatrix{
\mathcal{F}_0^\bullet \oplus (\mathcal{F}'_0)^\bullet
\ar[d] \ar[r] &
(\mathcal{F}_0^\bullet \otimes_{\mathcal{O}_0} \mathcal{I})[2] \oplus
((\mathcal{F}'_0)^\bullet \otimes_{\mathcal{O}_0} \mathcal{I})[2] \ar[d] \\
\mathcal{G}_0^\bullet \ar[r] &
(\mathcal{G}_0^\bullet \otimes_{\mathcal{O}_0} \mathcal{I})[2]
}
$$
を得る。垂直射は各直和因子上で擬同型を与えるので、
$D(\mathcal{O}_0)$ における所望の可換性が従う。

\medskip\noindent
良定義性が確立されたので、写像との両立性に関する主張は第二段落の結果から従う。
\end{proof}

\begin{lemma}
\label{defos-lemma-induced-map}
$(\mathcal{C}, \mathcal{O})$ を環付きサイトとする。
$\alpha : K \to L$ を $D^-(\mathcal{O})$ の写像とし、
$\mathcal{F}$ を $\mathcal{O}$-加群の層とする。$n \in \mathbf{Z}$ とする。
\begin{enumerate}
\item $i \geq n$ に対して $H^i(\alpha)$ が同型ならば、
$i \geq n$ に対して
$H^i(\alpha \otimes_\mathcal{O}^\mathbf{L} \text{id}_\mathcal{F})$ も同型である。
\item $i > n$ に対して $H^i(\alpha)$ が同型で、$i = n$ に対して全射ならば、
$i > n$ に対して
$H^i(\alpha \otimes_\mathcal{O}^\mathbf{L} \text{id}_\mathcal{F})$ は同型であり、
$i = n$ に対して全射である。
\end{enumerate}
\end{lemma}

\begin{proof}
区別三角形
$$
K \to L \to C \to K[1]
$$
を選ぶ。(2) の場合、$i \geq n$ に対して $H^i(C) = 0$ である。
したがって『導来圏』、補題 \ref{derived-lemma-negative-vanishing} の双対により、
$i \geq n$ に対して
$H^i(C \otimes_\mathcal{O}^\mathbf{L} \mathcal{F}) = 0$ である。
これから、$i > n$ に対して
$H^i(\alpha \otimes_\mathcal{O}^\mathbf{L} \text{id}_\mathcal{F})$ は同型であり、
$i = n$ に対して全射であることが従う。
(1) の場合には、さらに $H^{n - 1}(L) \to H^{n - 1}(C)$ が全射である。
図式
$$
\xymatrix{
H^{n - 1}(L) \otimes_\mathcal{O} \mathcal{F} \ar[r] \ar[d] &
H^{n - 1}(C) \otimes_\mathcal{O} \mathcal{F} \ar@{=}[d] \\
H^{n - 1}(L \otimes_\mathcal{O}^\mathbf{L} \mathcal{F}) \ar[r] &
H^{n - 1}(C \otimes_\mathcal{O}^\mathbf{L} \mathcal{F})
}
$$
を考えると、下の水平射が全射であることが分かる。
これと先の結果を合わせると、
$H^n(\alpha \otimes_\mathcal{O}^\mathbf{L} \text{id}_\mathcal{F})$ が同型である。
\end{proof}

\begin{lemma}
\label{defos-lemma-canonical-class-obstruction}
$\mathcal{C}$ をサイトとする。$\mathcal{O} \to \mathcal{O}_0$ を環の層の全射とし、
その核を平方零イデアル層 $\mathcal{I}$ とする。
$D^-(\mathcal{O}_0)$ の任意の対象 $K_0$ に対して、次は同値である。
\begin{enumerate}
\item 補題 \ref{defos-lemma-canonical-class} で構成された類
$\omega(K_0) \in
\Ext^2_{\mathcal{O}_0}(K_0, K_0 \otimes_{\mathcal{O}_0} \mathcal{I})$
は零である。
\item $D(\mathcal{O}_0)$ において
$K \otimes_\mathcal{O}^\mathbf{L} \mathcal{O}_0 = K_0$
を満たす $K \in D^-(\mathcal{O})$ が存在する。
\end{enumerate}
\end{lemma}

\begin{proof}
$K$ を (2) のものとする。$K$ を平坦 $\mathcal{O}$-加群の上に有界な複体
$\mathcal{F}^\bullet$ で表せる。このとき
$\mathcal{F}_0^\bullet =
\mathcal{F}^\bullet \otimes_{\mathcal{O}} \mathcal{O}_0$ は
$D(\mathcal{O}_0)$ において $K_0$ を表す。
$\mathcal{F}^\bullet$ は複体なので
$d_{\mathcal{F}^\bullet} \circ d_{\mathcal{F}^\bullet} = 0$ であり、
$\omega(K_0)$ の構成そのものから、これが零であることが分かる。

\medskip\noindent
(1) を仮定する。$\mathcal{F}^n$ と
$d : \mathcal{F}^n \to \mathcal{F}^{n + 1}$ を
$\omega(K_0)$ の構成におけるものとする。$\omega(K_0)$ が零であることから、写像
$$
\omega = d \circ d : \mathcal{F}_0^\bullet
\longrightarrow
(\mathcal{F}_0^\bullet \otimes_{\mathcal{O}_0} \mathcal{I})[2]
$$
は $D(\mathcal{O}_0)$ において零である。導来圏を
$\mathcal{O}_0$-加群の複体のホモトピー圏の局所化として定義することにより、
擬同型 $\alpha : \mathcal{G}_0^\bullet \to \mathcal{F}_0^\bullet$ と
$\mathcal{O}_0$-加群の写像
$h^n : \mathcal{G}_0^n \to
\mathcal{F}_0^{n + 1} \otimes_\mathcal{O} \mathcal{I}$
で次を満たすものが存在する。
$$
\omega \circ \alpha =
d_{\mathcal{F}_0^\bullet \otimes \mathcal{I}} \circ h +
h \circ d_{\mathcal{G}_0^\bullet}
$$
次のように置く。
$$
\mathcal{H}^n = \mathcal{F}^n \times_{\mathcal{F}^n_0} \mathcal{G}_0^n
$$
また、次のように定義する。
$$
d' : \mathcal{H}^n \longrightarrow \mathcal{H}^{n + 1},\quad
(f^n, g_0^n) \longmapsto (d(f^n) - h^n(g_0^n), d(g_0^n))
$$
ここでは、
$\mathcal{F}_0^{n + 1} \otimes_{\mathcal{O}_0} \mathcal{I} =
\mathcal{F}^{n + 1} \otimes_\mathcal{O} \mathcal{I} =
\mathcal{I}\mathcal{F}^{n + 1} \subset \mathcal{F}^{n + 1}$.
であることを用いて明らかな記法を採用した。$h^n$ の選択と $\omega$ の定義から、
$d' \circ d' = 0$ が確かめられる。したがって $\mathcal{H}^\bullet$ は
$D(\mathcal{O})$ の対象を定める。一方、$\mathcal{O}$-加群の複体の短完全列
$$
0 \to \mathcal{F}_0^\bullet \otimes_{\mathcal{O}_0} \mathcal{I} \to
\mathcal{H}^\bullet \to \mathcal{G}_0^\bullet \to 0
$$
がある。なお、
$\mathcal{H}^\bullet \otimes_\mathcal{O}^\mathbf{L} \mathcal{O}_0$ が
$K_0$ と同型であることを示す必要がある。
$\mathcal{E}^\bullet$ を平坦 $\mathcal{O}$-加群の上に有界な複体として、
擬同型 $\mathcal{E}^\bullet \to \mathcal{H}^\bullet$ を選ぶ。
すると可換図式
$$
\xymatrix{
0 \ar[r] &
\mathcal{E}^\bullet \otimes_\mathcal{O} \mathcal{I} \ar[d]^\beta \ar[r] &
\mathcal{E}^\bullet \ar[d]^\gamma \ar[r] &
\mathcal{E}_0^\bullet \ar[d]^\delta \ar[r] &
0 \\
0 \ar[r] &
\mathcal{F}_0^\bullet \otimes_{\mathcal{O}_0} \mathcal{I} \ar[r] &
\mathcal{H}^\bullet \ar[r] &
\mathcal{G}_0^\bullet \ar[r] &
0
}
$$
を得る。$\delta$ は擬同型であると主張する。$i \gg 0$ に対して
$H^i(\delta)$ は同型なので、$i \geq n$ に対して $H^i(\delta)$ が
同型であるような $n$ について降下帰納法を使える。ここで、
$\mathcal{E}^\bullet \otimes_\mathcal{O} \mathcal{I}$
は
$\mathcal{E}_0^\bullet \otimes_{\mathcal{O}_0}^\mathbf{L} \mathcal{I}$,
を表し、
$\mathcal{F}_0^\bullet \otimes_{\mathcal{O}_0} \mathcal{I}$
は
$\mathcal{G}_0^\bullet \otimes_{\mathcal{O}_0}^\mathbf{L} \mathcal{I}$,
を表す。また $D(\mathcal{O}_0)$ の写像として
$\beta = \delta \otimes_{\mathcal{O}_0}^\mathbf{L} \text{id}_\mathcal{I}$
である。実際、
$\beta =
(\alpha \otimes \text{id}_\mathcal{I})
\circ
(\delta \otimes \text{id}_\mathcal{I})$.
だからである。次数 $\geq n$ で $H^i(\delta)$ が同型であると仮定する。
上で述べたことと補題 \ref{defos-lemma-induced-map} により、$\beta$ についても同じである。
そこで図式
$$
\xymatrix@C=1.2em{
H^{n - 1}(\mathcal{E}^\bullet \otimes_\mathcal{O} \mathcal{I})
\ar[r] \ar[d]^{H^{n - 1}(\beta)} &
H^{n - 1}(\mathcal{E}^\bullet) \ar[r] \ar[d] &
H^{n - 1}(\mathcal{E}_0^\bullet) \ar[r] \ar[d]^{H^{n - 1}(\delta)} &
H^n(\mathcal{E}^\bullet \otimes_\mathcal{O} \mathcal{I})
\ar[r] \ar[d]^{H^n(\beta)} &
H^n(\mathcal{E}^\bullet) \ar[d] \\
H^{n - 1}(\mathcal{F}_0^\bullet \otimes_\mathcal{O} \mathcal{I}) \ar[r] &
H^{n - 1}(\mathcal{H}^\bullet) \ar[r] &
H^{n - 1}(\mathcal{G}_0^\bullet) \ar[r] &
H^n(\mathcal{F}_0^\bullet \otimes_\mathcal{O} \mathcal{I}) \ar[r] &
H^n(\mathcal{H}^\bullet)
}
$$
を考える。『ホモロジー』、補題 \ref{homology-lemma-four-lemma} により、
$H^{n - 1}(\delta)$ は全射である。すると補題 \ref{defos-lemma-induced-map} により、
$H^{n - 1}(\beta)$ も全射である。もう一度『ホモロジー』、補題
\ref{homology-lemma-four-lemma} を使うと、$H^{n - 1}(\delta)$ は同型である。
帰納法により主張が成り立つので、$\delta$ は擬同型である。
これが示すべきことであった。
\end{proof}

\begin{lemma}
\label{defos-lemma-lift-map-complexes}
$\mathcal{C}$ をサイトとし、$\mathcal{O} \to \mathcal{O}_0$ を
環の層の全射とする。次のデータが与えられていると仮定する。
\begin{enumerate}
\item $\mathcal{O}$-加群の複体 $\mathcal{F}^\bullet$。
\item $\mathcal{O}_0$-加群の複体 $\mathcal{K}_0^\bullet$。
\item 擬同型 $\mathcal{K}_0^\bullet \to
\mathcal{F}^\bullet \otimes_\mathcal{O} \mathcal{O}_0$。
\end{enumerate}
このとき擬同型 $\mathcal{G}^\bullet \to \mathcal{F}^\bullet$ で、複体の写像
$\mathcal{G}^\bullet  \otimes_\mathcal{O} \mathcal{O}_0 \to
\mathcal{F}^\bullet \otimes_\mathcal{O} \mathcal{O}_0$ が
$\mathcal{O}_0$-加群の複体のホモトピー圏において
$\mathcal{K}_0^\bullet$ を経由するようなものが存在する。
\end{lemma}

\begin{proof}
$\mathcal{F}_0^\bullet =
\mathcal{F}^\bullet \otimes_\mathcal{O} \mathcal{O}_0$ と置く。
『導来圏』、補題 \ref{derived-lemma-make-surjective} により、与えられた写像には分解
$$
\mathcal{K}_0^\bullet \to \mathcal{L}_0^\bullet \to \mathcal{F}_0^\bullet
$$
が存在し、第一の射はホモトピーを除いて逆射をもち、第二の射は各次数で分裂全射である。
したがって $\mathcal{K}_0^\bullet \to \mathcal{F}_0^\bullet$ が
各次数で全射であると仮定してよい。この場合、
$$
\mathcal{G}^n = \mathcal{F}^n \times_{\mathcal{F}^n_0} \mathcal{K}_0^n
$$
と取れば、すべて明らかである。
\end{proof}

\begin{lemma}
\label{defos-lemma-inf-obs-map-defo-complex}
$\mathcal{C}$ をサイトとする。$\mathcal{O} \to \mathcal{O}_0$ を環の層の全射とし、
その核を平方零イデアル層 $\mathcal{I}$ とする。
$K, L \in D^-(\mathcal{O})$ とし、$D^-(\mathcal{O}_0)$ において
$K_0 = K \otimes_\mathcal{O}^\mathbf{L} \mathcal{O}_0$ および
$L_0 = L \otimes_\mathcal{O}^\mathbf{L} \mathcal{O}_0$ と置く。
$D(\mathcal{O}_0)$ における $\alpha_0 : K_0 \to L_0$ が与えられると、標準元
$$
o(\alpha_0) \in \Ext^1_{\mathcal{O}_0}(K_0,
L_0 \otimes_{\mathcal{O}_0}^\mathbf{L} \mathcal{I})
$$
が存在する。この元が消えることは、
$\alpha_0 = \alpha \otimes_\mathcal{O}^\mathbf{L} \text{id}$ を満たす
$D(\mathcal{O})$ の写像 $\alpha : K \to L$ が存在するための必要十分条件である。
\end{lemma}

\begin{proof}
$\alpha_0$ を持ち上げる $\alpha : K \to L$ を求めることは、合成
$K \xrightarrow{\alpha} L \to L_0$ が合成
$K \to K_0 \xrightarrow{\alpha_0} L_0$ に等しくなるような
$\alpha : K \to L$ を求めることと同じである。短完全列
$0 \to \mathcal{I} \to \mathcal{O} \to \mathcal{O}_0 \to 0$ は、
$D(\mathcal{O})$ に標準的な区別三角形
$$
L \otimes_\mathcal{O}^\mathbf{L} \mathcal{I} \to
L \to
L_0 \to
(L \otimes_\mathcal{O}^\mathbf{L} \mathcal{I})[1]
$$
を与える。『導来圏』、補題 \ref{derived-lemma-representable-homological} により、合成
$$
K \to K_0 \xrightarrow{\alpha_0} L_0 \to
(L \otimes_\mathcal{O}^\mathbf{L} \mathcal{I})[1]
$$
が零であることと、$\alpha_0$ を持ち上げる $\alpha : K \to L$ を
見つけられることは同値である。この合成は、随伴により
$$
\Hom_{D(\mathcal{O})}(K, (L \otimes_\mathcal{O}^\mathbf{L} \mathcal{I})[1]) =
\Hom_{D(\mathcal{O}_0)}(K_0,
(L \otimes_\mathcal{O}^\mathbf{L} \mathcal{I})[1]) =
\Ext^1_{\mathcal{O}_0}(K_0,
L_0 \otimes_{\mathcal{O}_0}^\mathbf{L} \mathcal{I})
$$
の元である。
\end{proof}

\begin{lemma}
\label{defos-lemma-first-order-defos-complex}
$\mathcal{C}$ をサイトとする。$\mathcal{O} \to \mathcal{O}_0$ を環の層の全射とし、
その核を平方零イデアル層 $\mathcal{I}$ とする。
$K_0 \in D^-(\mathcal{O})$ とする。$K_0$ の持ち上げとは、
$D^-(\mathcal{O})$ の対象 $K$ と $D(\mathcal{O}_0)$ の同型
$\alpha_0 : K \otimes_\mathcal{O}^\mathbf{L} \mathcal{O}_0 \to K_0$
からなる組 $(K, \alpha_0)$ のことである。
\begin{enumerate}
\item 持ち上げ $(K, \alpha)$ が与えられると、この組の自己同型群は標準的に写像
$$
\Ext^{-1}_{\mathcal{O}_0}(K_0, K_0)
\longrightarrow
\Hom_{\mathcal{O}_0}(K_0, K_0 \otimes_{\mathcal{O}_0}^\mathbf{L} \mathcal{I})
$$
の余核である。
\item 持ち上げが存在するならば、持ち上げの同型類の集合は
\par\noindent
$\Ext^1_{\mathcal{O}_0}(K_0,
K_0 \otimes_{\mathcal{O}_0}^\mathbf{L} \mathcal{I})$ の下で主等質である。
\end{enumerate}
\end{lemma}

\begin{proof}
$(K, \alpha)$ の自己同型とは、
$\varphi \otimes_\mathcal{O} \text{id}_{\mathcal{O}_0} = \text{id}$ を満たす
$D(\mathcal{O})$ の写像 $\varphi : K \to K$ のことである。
これは
$$
K \xrightarrow{\varphi - \text{id}} K \to
K \otimes_\mathcal{O}^\mathbf{L} \mathcal{O}_0
$$
が零であることと同じである。したがって、自己同型群は写像
$$
\Hom_\mathcal{O}(K, K_0[-1])
\longrightarrow
\Hom_\mathcal{O}(K, K_0 \otimes_{\mathcal{O}_0}^\mathbf{L} \mathcal{I})
$$
の余核であると結論する。実際、区別三角形
$$
K \otimes_\mathcal{O}^\mathbf{L} \mathcal{I} \to
K \to
K \otimes_\mathcal{O}^\mathbf{L} \mathcal{O}_0 \to
(K \otimes_\mathcal{O}^\mathbf{L} \mathcal{I})[1]
$$
は $D(\mathcal{O})$ におけるものであり、さらに『導来圏』、補題
\ref{derived-lemma-representable-homological} を用いた。
これを補題に現れる群へ移すには、制限関手
$D(\mathcal{O}_0) \to D(\mathcal{O})$ と
$- \otimes_\mathcal{O} \mathcal{O}_0 : D(\mathcal{O}) \to D(\mathcal{O}_0)$
の随伴を用いる。これで (1) が証明された。

\medskip\noindent
(2) の証明。
$D(\mathcal{O})$ において
$K_0 = K \otimes_\mathcal{O}^\mathbf{L} \mathcal{O}_0$ であると仮定する。
補題 \ref{defos-lemma-inf-obs-map-defo-complex} により、持ち上げ
$(K', \alpha_0)$ を $\alpha_0$ の持ち上げに対する障害 $o(\alpha_0)$ へ送る写像は、
組の同型類の集合から
$\Ext^1_{\mathcal{O}_0}(K_0,
K_0 \otimes_{\mathcal{O}_0}^\mathbf{L} \mathcal{I})$ への標準的な単射を定める。
証明を終えるため、これが全射であることを示す。
補題の $\Ext^1$ に属する
$\xi : K_0 \to (K_0 \otimes_{\mathcal{O}_0}^\mathbf{L} \mathcal{I})[1]$ を取る。
$K$ を表す平坦 $\mathcal{O}$-加群の上に有界な複体
$\mathcal{F}^\bullet$ を選ぶ。写像 $\xi$ は $t \circ s^{-1}$ と表せる。
ここで $s : \mathcal{K}_0^\bullet \to \mathcal{F}_0^\bullet$ は擬同型であり、
$t : \mathcal{K}_0^\bullet \to
\mathcal{F}_0^\bullet \otimes_{\mathcal{O}_0} \mathcal{I}[1]$ は複体の写像である。
補題 \ref{defos-lemma-lift-map-complexes} により、$\mathcal{O}$-加群の複体の擬同型
$\mathcal{G}^\bullet \to \mathcal{F}^\bullet$ で、
$\mathcal{G}_0^\bullet \to \mathcal{F}_0^\bullet$ がホモトピーを除いて
$s$ を経由するようなものが存在すると仮定してよい。
$\mathcal{G}^\bullet$ は平坦 $\mathcal{O}$-加群の上に有界な複体で置き換えてよいし、
実際そのように置き換える（そのような複体から $\mathcal{G}^\bullet$ への
擬同型を選んで置き換える）。すると $\xi$ は、複体の写像
$t : \mathcal{G}_0^\bullet \to
\mathcal{F}_0^\bullet \otimes_{\mathcal{O}_0} \mathcal{I}[1]$ と擬同型
$\mathcal{G}_0^\bullet \to \mathcal{F}_0^\bullet$ によって表される。
次のように置く。
$$
\mathcal{H}^n = \mathcal{F}^n \times_{\mathcal{F}_0^n} \mathcal{G}_0^n
$$
微分は
$$
\mathcal{H}^n \to \mathcal{H}^{n + 1},\quad
(f^n, g_0^n) \mapsto (d(f^n) + t(g_0^n), d(g_0^n))
$$
で定める。これは
$\mathcal{F}_0^{n + 1} \otimes_{\mathcal{O}_0} \mathcal{I} =
\mathcal{F}^{n + 1} \otimes_\mathcal{O} \mathcal{I} =
\mathcal{I}\mathcal{F}^{n + 1} \subset \mathcal{F}^{n + 1}$.
であるため意味をもつ。$\mathcal{H}^\bullet$ が $\mathcal{O}$-加群の複体であることを
示す計算は省略する。構成により、$\mathcal{O}$-加群の複体の短完全列
$$
0 \to \mathcal{F}_0^\bullet \otimes_{\mathcal{O}_0} \mathcal{I} \to
\mathcal{H}^\bullet \to \mathcal{G}_0^\bullet \to 0
$$
がある。補題 \ref{defos-lemma-canonical-class-obstruction} の証明とまったく同様に、
この列が $D(\mathcal{O}_0)$ の同型
$\alpha_0 :
\mathcal{H}^\bullet \otimes_\mathcal{O}^\mathbf{L} \mathcal{O}_0 \to
\mathcal{G}_0^\bullet$ を誘導することが示される。言い換えれば、組
$(\mathcal{H}^\bullet, \alpha_0)$ を構成したことになる。
$o(\alpha_0) = \xi$ の確認は省略する。ヒントとして、この場合には
$\alpha_0$ の成分を持ち上げる写像
$\mathcal{H}^n \to \mathcal{F}^n$（微分とは両立しない）があるので、
$o(\alpha_0)$ を明示的に計算できる。これで証明が完了した。
\end{proof}
